\documentclass[11pt]{article}
\usepackage[a4paper,margin=27mm,top=28mm,bottom=30mm]{geometry}
\usepackage[T1]{fontenc}
\usepackage[utf8]{inputenc}
\usepackage{mathpazo}
\usepackage[scaled=0.92]{helvet}
\usepackage{microtype}
\usepackage{amsmath,amssymb,amsthm,mathtools}
\usepackage{booktabs,array,tabularx}
\usepackage{enumitem}
\usepackage{tikz}
\usetikzlibrary{calc,arrows.meta}
\usepackage{aliascnt}
\usepackage[colorlinks=true,linkcolor=black!70!blue,citecolor=black!60!blue,urlcolor=black!60!blue]{hyperref}
\usepackage[capitalise,noabbrev]{cleveref}
\usepackage{caption}
\captionsetup{font=small,labelfont=bf,width=0.92\textwidth}
\setlist{itemsep=2pt,topsep=3pt}
\linespread{1.04}
\theoremstyle{plain}
\newtheorem{theorem}{Theorem}[section]
\newaliascnt{proposition}{theorem}
\newtheorem{proposition}[proposition]{Proposition}
\aliascntresetthe{proposition}
\newaliascnt{lemma}{theorem}
\newtheorem{lemma}[lemma]{Lemma}
\aliascntresetthe{lemma}
\newaliascnt{corollary}{theorem}
\newtheorem{corollary}[corollary]{Corollary}
\aliascntresetthe{corollary}
\theoremstyle{definition}
\newaliascnt{definition}{theorem}
\newtheorem{definition}[definition]{Definition}
\aliascntresetthe{definition}
\theoremstyle{remark}
\newaliascnt{remark}{theorem}
\newtheorem{remark}[remark]{Remark}
\aliascntresetthe{remark}
\newtheorem*{mainA}{Theorem A}
\newtheorem*{gibbsB}{Corollary B}
\newtheorem*{mainC}{Theorem C}
\newtheorem*{gibbsD}{Corollary D}
\newtheorem*{mainE}{Theorem E}
\newtheorem*{gibbsF}{Corollary F}
\crefname{theorem}{Theorem}{Theorems}
\crefname{proposition}{Proposition}{Propositions}
\crefname{lemma}{Lemma}{Lemmas}
\crefname{corollary}{Corollary}{Corollaries}
\crefname{definition}{Definition}{Definitions}
\crefname{remark}{Remark}{Remarks}
\newcommand{\Z}{\mathbb Z}
\newcommand{\N}{\mathbb N}
\newcommand{\Prob}{\mathbb P}
\newcommand{\eps}{\varepsilon}
\newcommand{\ind}{\mathbf 1}
\newcommand{\botv}{\bot}
\newcommand{\Mo}{\{-1,0,1\}^2}
\newcommand{\HH}{\mathcal H}
\newcommand{\MM}{\mathcal M}
\newcommand{\RR}{\mathcal R}
\newcommand{\osc}{\operatorname{osc}}
\newcommand{\init}{\operatorname{init}}
\newcommand{\End}{\operatorname{End}}
\newcommand{\Orig}{\operatorname{Orig}}
\newcommand{\Pin}{\operatorname{Pin}}
\newcommand{\proved}{}
\newcommand{\recalled}{}
\newcommand{\fail}{\mathsf{FAIL}}
\newcommand{\halt}{\mathsf{HALT}}
\hypersetup{pdftitle={One noisy cellular automaton with an uncomputable noise diagram},pdfauthor={Seth Douglas and Nidhal Mghirbi}}
\title{\vspace{-1.2cm}\LARGE One noisy cellular automaton\\ with an uncomputable noise diagram}
\author{Seth Douglas \and Nidhal Mghirbi}
\date{\small Version 1.4.0, November 2026}
\begin{document}
\maketitle
\begin{abstract}
\noindent
We construct one computably specified finite-alphabet, finite-range cellular automaton on $\Z^3$ with computable disjoint nonempty rational rate windows. Throughout the $n$th window, uniform whole-symbol replacement is ergodic exactly when a fixed universal machine halts on $n$. Halting gives exponential full-symbol coalescence over all initial configurations, with one fixed positive computable decay coefficient and an effective halting-dependent prefactor uniform on each window; nonhalting gives two stationary laws invariant under every spatial translation. A corner-deterministic chart, causal L-frame computations, history strips, private interval sensors and a packed wrong-gate boundary give the construction. Its space-time interaction has fixed finite range and $\{0,1\}$ energies on $\Z^4$: on computable rational inverse-temperature windows, halting is equivalent to uniqueness among all Gibbs measures and gives exponential weak spatial mixing in arbitrary volumes. The earlier results on $\Z^4$, $\Z^5$, $\Z^6$ and $\Z^7$ are retained. The Lean verification covers the stated original six-dimensional results; the three- and four-dimensional automaton results are not formalized.

\end{abstract}
\begin{center}\small AI systems assisted with this work; see ``How this work was done.''\end{center}
\setcounter{tocdepth}{1}
\makeatletter
{\footnotesize\linespread{1}\selectfont
\renewcommand{\l@section}[2]{\@dottedtocline{1}{0em}{2em}{\bfseries #1}{\bfseries #2}}
\tableofcontents}
\makeatother
\clearpage

\section{Introduction}\label{sec:intro}

Can the long-run behaviour of one cellular automaton, subjected to the same noise mechanism at different rates, encode an undecidable problem? Independent replacement gives every symbol positive probability at every update. It nevertheless permits persistent information. We show that changing only the replacement rate can select which of two behaviours occurs: forgetting every initial state, or retaining a bit in two stationary laws.

Here is the main result in plain words. Fix a universal deterministic machine $U$ whose halting set on positive unary inputs is undecidable.
\begin{mainA}
\proved\ There is one computably specified cellular automaton on $\Z^6$, with a finite alphabet and a finite neighbourhood, and computable disjoint rational open intervals $I_n$, such that at every rate in $I_n$ uniform whole-symbol replacement makes the automaton ergodic exactly when $U(n)$ halts. Nonhalting gives at least two translation-invariant stationary laws. Halting gives a unique stationary law attracting every initial law.
\end{mainA}

In particular, the noise diagram, the set of rates at which $F_\eps$ is ergodic, is undecidable at rational rates for this rule (\cref{cor:undecidable}).

The dimension can be lowered while retaining an exponential common-map bound.
\begin{mainC}
There is one computably specified finite-alphabet, finite-range automaton on $\Z^4$, with fixed rational $C_4>0$ and integer $n_4$, and computable pairwise disjoint nonempty rational open intervals
\[
 I_n^{(4)}=\left(\frac{2C_4}{2\cdot4^{n+n_4}-1},
 \frac{21C_4}{10(2\cdot4^{n+n_4}-1)}\right),\qquad n\ge1.
\]
At every rate in $I_n^{(4)}$ uniform whole-symbol replacement is ergodic if and only if $U(n)$ halts. Nonhalting gives two stationary laws invariant under every spatial translation. Halting gives finite-window coupling from the past over all initial configurations, attraction from every initial law, and, for every finite $F\subset\Z^4$ and integer $t\ge0$,
\[
 \begin{gathered}
 \Prob\{\text{the shared depth-$t$ full-symbol map is nonconstant on }F\}\\
 \le |F|C_n\exp\!\left[-\frac{t}{4\cdot701(4^{n+n_4}-1+52)}\right].
 \end{gathered}
\]
The integer prefactor $C_n$ is effectively findable on halting inputs and is uniform over the whole window; the displayed rate is total computable from $n$. The construction uses three transverse $16\times16$ boxes, threshold $512$ of $768$, depth $500$ and drift $(58,58)$, with all fixed constants specified in \cref{sec:lower}.
\end{mainC}

The local rule receives no numerical noise parameter. Rulers of different lengths see different numbers of random replacements. Their responses select an integer, and finite computation boards read it. All boards use the same machine and the same alphabet. The computation is in the spatial constraints of a chart; a halting computation is excluded by forbidding its halting state. \Cref{fig:mechanism} shows the two parts of this selection.

Theorem~A (in the form of \cref{thm:main}), \cref{cor:undecidable} and the one-rule conclusion above are proved in Lean; \cref{app:lean} states the scope and the checks.

The numbered additions strengthen this dichotomy: halting gives exponential all-input coalescence; nonhalting gives continuum many extremal stationary laws (\cref{thm:exponential,thm:continuum}). The two translation-invariant bit laws in Theorem~A remain a separate conclusion.

Through the space--time correspondence and \cite[Corollary~13.7]{DMst}, Corollary~B in \cref{sec:parameters} gives one computably specified translation-invariant finite-range $\{0,1\}$-energy interaction on $\Z^7$. On computable rational inverse-temperature windows $\mathcal J_n$, halting gives exactly one Gibbs measure, with no invariance assumed, and exponential weak spatial mixing in arbitrary volumes; nonhalting gives at least two translation-invariant Gibbs measures and continuum many extremal Gibbs measures. The latter extremes need not be translation invariant. A fixed field direction has a derivative at zero exactly on the halting windows, although the zero-field pressure is analytic (\cref{thm:field}).

\begin{gibbsD}
There is one computably specified translation-invariant finite-range interaction on $\Z^5$, with a finite alphabet and individual energies in $\{0,1\}$, and computable pairwise disjoint nonempty rational open inverse-temperature intervals $\mathcal J_n^{(5)}$. For every $n\ge1$ and every $\beta\in\mathcal J_n^{(5)}$, uniqueness among all Gibbs measures is equivalent to halting of $U(n)$. Halting also gives exponential weak spatial mixing in every finite volume, uniformly over all exterior configurations and including volumes with holes or disconnected components. Nonhalting gives two Gibbs measures invariant under every translation of $\Z^5$. The mixing constants may depend on $\beta$; the interaction is fixed independently of $n$ and $\beta$.
\end{gibbsD}

\begin{mainE}
There is one computably specified finite-alphabet finite-range automaton on $\Z^3$ with computable pairwise disjoint nonempty rational open rate windows $I_n^{(3)}$, on which uniform whole-symbol replacement is ergodic if and only if $U(n)$ halts. Nonhalting gives two stationary laws invariant under every spatial translation. Halting gives finite-window coupling from the past over all starts, attraction from every initial law, and
\[
 \Prob\{\text{the shared depth-$t$ full-symbol map is nonconstant on }F\}
 \le |F|C_ne^{-ct}
\]
for every finite $F\subset\Z^3$ and integer $t\ge0$, uniformly throughout the window. Here $c>0$ is the fixed computable rational in \eqref{eq:three-rate}; an integer $C_n$ is effectively found on halting inputs. All alphabets, neighbourhoods and parameters are fixed before selecting $n$.
\end{mainE}

\begin{gibbsF}
One computably specified translation-invariant finite-range $\{0,1\}$-energy interaction on $\Z^4$ has computable pairwise disjoint nonempty rational inverse-temperature windows $\mathcal J_n^{(4)}$ on which uniqueness among all Gibbs measures is equivalent to halting of $U(n)$. Halting gives exponential weak spatial mixing in arbitrary finite volumes, uniformly over all exterior configurations. Nonhalting gives two Gibbs measures invariant under every translation of $\Z^4$.
\end{gibbsF}

Theorem~E and Corollary~F are proved in \cref{sec:three}; neither is formalized in Lean.
Theorem~C and Corollary~D are proved in \cref{sec:lower}; neither is formalized in Lean. The six-dimensional Theorem~A, seven-dimensional Corollary~B and the stated Lean verification of the earlier results are retained.

The hard-model paper asks the following questions \cite[Open problems~9.1--9.2]{DMhard}:
\begin{quote}
\lq\lq Is there one translation-invariant finite-range interaction on $\Z^d$ with finite rational energies, and computable rational inverse temperatures $\beta_n$, such that uniqueness of the translation-invariant Gibbs measure at $\beta_n$ is undecidable?\rq\rq

\lq\lq Is there one interaction, with or without hard constraints, for which uniqueness among all Gibbs measures, as a function of $\beta$, is undecidable?\rq\rq
\end{quote}
Version~1.0.0's Corollary~B already answered Open problem~9.1, by taking $\beta_n$ to be the midpoint of $\mathcal J_n$. The upgraded corollary answers Open problem~9.2 in dimension seven with finite energies. Corollary~D answers the same question in dimension five with finite energies, and Corollary~F answers it in dimension four.

\paragraph{Changes from version 1.3.0.} This version adds Theorem~E on $\Z^3$ and Corollary~F on $\Z^4$, with complete proofs in \cref{sec:three}. A deterministic cover includes the canonical routing; a causal L-frame tableau replaces the old computation; history strips and relocated lanes preserve the sensors; complete boundary witnesses and one candidate field per history yield a fixed positive decay coefficient. The seed routing is checked directly without a typical-phase assumption. Two independent exact-arithmetic programs certify the fresh scale. All existing statements and numbers, Section~10 and the Lean appendix are retained. The new results are not formalized.
\paragraph{Changes from version 1.2.0.} Version~1.3.0 added Theorem~C on $\Z^4$ and Corollary~D on $\Z^5$, with the full construction, finite-history all-input proof and all-DLR transfer in \cref{sec:lower}. Three transverse $16\times16$ boxes at threshold $512$ of $768$ and marked-parent density replace the five-box geometry and every-step coherence for this second rule. The planar chart, computation boards and sensors are retained as finite algorithms with newly selected fixed parameters. No existing statement or statement number changes; the constants of the earlier rule and the Lean appendix are unchanged. The two lower-dimensional results are not formalized.

\paragraph{Changes from version 1.1.0.} Version~1.2.0 added \cref{app:lean}, reporting the Lean proofs of Theorem~A (in the form of Theorem~9.3), Corollary~9.4 and the one-rule conclusion after Theorem~A, and updated the verification account. It also reconciled the historical count of corrections. No mathematical statement, statement number or construction constant changes.

\paragraph{Changes from version 1.0.0.} Corrections: (i) in \cref{lem:pole} the time-span clause now requires unary links to preserve time, and the two-terminal clause requires a positive support edge; (ii) in \cref{sec:chart} the gate classification is restricted to finite bottom rows, and (iii) the phase decoder uses finite approximants instead of liftability; (iv) the birth-contour clock in \cref{lem:fullcoupling} is stated once; (v) the proof of \cref{prop:local} counts each helper parent once, with $\kappa$ unchanged. Corollary~B is strengthened to uniqueness among all Gibbs measures, with exponential weak spatial mixing, on halting windows and continuum many extremal Gibbs measures on nonhalting windows; its version~1.0.0 sentence on uniformly recurrent and periodic Gibbs measures is subsumed and removed. Separate numbered results add exponential coalescence, extremal continua, field response, a computable robustness radius and effective corollaries. Theorem~A, Theorem~9.3, Corollary~9.4 and the construction constants keep their statements and values.

\begin{figure}[ht]
\centering
\begin{tikzpicture}[font=\small,>=Stealth]
\draw[->] (0,3.4)--(12.4,3.4) node[right] {$\eps$};
\foreach \x/\lab in {1.1/{I_{n+2}},3.8/{I_{n+1}},7.1/{I_n}} {
 \draw[line width=4pt,blue!45] (\x,3.4)--(\x+1.5,3.4);
 \node at (\x+0.75,3.9) {$\lab$};
}
\node[anchor=east] at (0,3.4) {$0$};
\node[align=center] at (6.3,2.7) {Each window: halting $\Rightarrow$ forgetting;\quad nonhalting $\Rightarrow$ two stationary laws};
\node[draw,rounded corners,align=center,minimum height=1cm] (r) at (1.4,1) {finite rulers\\sample noise};
\node[draw,rounded corners,align=center,minimum height=1cm] (c) at (5.1,1) {rank cutoff\\selects input};
\node[draw,rounded corners,align=center,minimum height=1cm] (b) at (9.6,1) {finite boards\\simulate $U$};
\draw[->,thick] (r)--(c); \draw[->,thick] (c)--(b);
\end{tikzpicture}
\caption{The intervals and their labels are schematic. The interval endpoints are computable; their ergodicity labels encode halting. The local rule uses finite sensor fields and the chart's finite geometry.}
\label{fig:mechanism}
\end{figure}

\paragraph{Relation to earlier work.}\recalled\ Toom's stability theorem and G\'acs's reliable automata established memory under positive noise \cite{Toom80,Gacs01}. Undecidability of ergodicity across descriptions goes back to Kurdyumov \cite{Kur78}; see \cite[Chapter~14]{Toometal90}. Toom proved it for binary PCA with transition probabilities in $\{0,\tfrac12,1\}$ and for continuous-time particle systems \cite{Toom00a,Toom00b}. Bu\v{s}i\'c, Mairesse and Marcovici showed that a deterministic one-dimensional CA is ergodic exactly when it is nilpotent \cite[Theorem~3.3]{BMM13}; Kari proved nilpotency undecidable \cite{Kari92}. Marcovici, Sablik and Taati developed ergodicity criteria for noisy classes \cite{MST19}. Marsan, Sablik and T\"orm\"a constructed a fixed one-dimensional automaton with at least two ergodicity transitions and asked which sets of rates can occur \cite[Section~6]{MST26}: \lq\lq which sets can be obtained as\rq\rq\ $\{\eps\in[0,1]:F_\eps\text{ is ergodic}\}$. For continuous Cantor maps under uniform noise, Marsan and Sablik realize every $G_\delta$ subset of $[0,1]$ containing $1$ simultaneously as the uniqueness set and the ergodicity set \cite[Theorem~3.19]{MS25}. They state that the uniqueness set of a computable map is $\Pi_2$-computable in their sense, that every $\Pi_2$-computable $G_\delta$ subset of $[0,1]$ containing $1$ is the uniqueness set of a computable map, and that for cellular automata it is still open to characterize the possible uniqueness sets \cite[Section~5]{MS25}. Their earlier work realizes sets of zero-noise limiting measures, using computations whose time budget is set by the noise rate \cite[Sections~5.1--5.2.1]{MS24}. Marcovici and Taati prove exponential forgetting for positive-rate PCA admitting a stationary Bernoulli law \cite{MT26}. The distinction between a family of rules and one fixed rule is recorded in \cref{tab:landscape}.

\begin{table}[t]
\centering\small
\begin{tabularx}{\textwidth}{@{}>{\raggedright\arraybackslash}p{.26\textwidth}>{\raggedright\arraybackslash}X>{\raggedright\arraybackslash}X@{}}
\toprule
Result & Object held fixed & Parameter or input\\
\midrule
Bu\v{s}i\'c--Mairesse--Marcovici \cite{BMM13} & Description varies & Undecidable ergodicity, already for deterministic CA\\[3pt]
Marsan--Sablik--T\"orm\"a \cite{MST26} & One CA on $\Z$ & At least two noise transitions\\[3pt]
Marsan--Sablik \cite{MS25} & One continuous Cantor map; shift commutation not required & Every $G_\delta$ set containing $1$, as uniqueness and ergodicity set; every $\Pi_2$-computable one as the uniqueness set of a computable map (their Section~5)\\[3pt]
Forgetting \cite{DMforget} & Computable family on $\Z^8$ & Halting selects ergodicity at fixed noise\\[3pt]
Windows \cite{DMwindows} & One CA on $\Z^4$ & Infinitely many alternating computable windows\\[3pt]
Hard-constraint \cite{DMhard} & One classical hard-constraint model & Halting-labelled Gibbs uniqueness windows\\[3pt]
Theorem A and \cref{thm:exponential} & One CA on $\Z^6$ & Halting-labelled windows; exponential all-input coalescence when halting\\
\bottomrule
\end{tabularx}
\caption{The closest results and the parameter in which their conclusions vary.}
\label{tab:landscape}
\end{table}

\recalled\ Hochman and Meyerovitch characterized multidimensional SFT entropies \cite{HM10}. Chazottes and Hochman, and Gayral, Sablik and Taati studied zero-temperature Gibbs limits and their computability \cite{CH10,GST25}. Cubitt, Perez-Garcia and Wolf proved spectral-gap undecidability \cite{CPGW15,CPGW22}; Bausch, Cubitt and Watson gave a continuous quantum parameter family with an uncomputable phase diagram \cite{BCW21}, and Purcell, Li and Cubitt placed a transition at an uncomputable number \cite{PLC24}.

\paragraph{What is ours.}\proved\ Theorem~A realizes halting-labelled windows for a cellular automaton: a shift-commuting map with one finite neighbourhood, within the class whose possible diagrams are left open in \cite[Section~5]{MS25}. Our construction combines a fixed computational chart, a finite-range interval controller, a joint estimate for sensor feedback, and a common-noise coupling. The hierarchical scaffold is from our hard-constraint paper \cite{DMhard}. Our forgetting paper supplies the checker/quorum approach \cite{DMforget}; our windows paper asks the fixed-rule question answered here \cite{DMwindows}. Our nilpotent paper gives, for each $n$, a nilpotent rule with $n$ separated memory windows and the mechanism of independent gates read from diagonal copies \cite[Section~3, Exact gate law]{DMnil}; here the intrinsic gate sources follow diagonal characteristics as well. Corollary~B uses the correspondence between translation-invariant stationary laws and translation-invariant space--time Gibbs measures due to Goldstein, Kuik, Lebowitz and Maes \cite{GKLM89}, in the form proved in our space--time paper \cite[Lemma~2.4 and Corollary~10.4]{DMst}, and obtains uniqueness among all Gibbs measures and weak spatial mixing from \cref{thm:exponential} through \cite[Corollary~13.7]{DMst}. The proof below establishes the stochastic statements needed for this construction and proves finite-board forcing in every phase. The new finite-history argument uses a finite motif certificate and coherence within one trajectory. Its direct precursor is the varying-rule all-Gibbs result of \cite[Corollary~8.8 and Remark~8.14]{DMforget}; common-map certificates and their effective search are \cite[Propositions~4.2--4.3]{DMst}. Here the same rule is selected by its noise rate. Chart conditioning gives the extremal continuum without stacking, unlike \cite[Corollary~8.17]{DMforget}. The field-response argument follows \cite[Corollary~10.7]{DMthermal}, and the zero-pressure normalization is recalled in \cite[Remark~10.13]{DMst}.

\paragraph{Dimension and scale.} One fixed rule now suffices on $\Z^3$, and its Gibbs interaction lives on $\Z^4$. Theorem~E uses one chart-space coordinate and two transverse coordinates, viewing the planar chart diagonally through physical time. Theorem~C keeps two chart-space coordinates and two transverse coordinates; Theorem~A uses two and four. Fixed-rule questions on $\Z^2$ and $\Z$ remain open. Each rule has its own terminating parameter recipe: \cref{sec:parameters}, \cref{sec:lower} and \cref{sec:three} specify the six-, four- and three-dimensional scales respectively.

\paragraph{Reading guide.} The dimension-three construction and dimension-four Gibbs transfer are in \cref{sec:three}. The four-dimensional construction and five-dimensional Gibbs transfer are in \cref{sec:lower}. \Cref{sec:parameters} contains the Gibbs corollary on $\Z^7$, field response and the effective consequences. The finite-history proof is in \cref{app:finite}, and the local-kernel radius is in \cref{app:robust}. \Cref{sec:rule} gives the noisy dynamics and its complete local rule, with constants selected in \cref{sec:parameters}. \Cref{sec:chart,sec:machine} construct the chart and prove the all-phase board lemma. \Cref{sec:poles,sec:sensors} give the contour and sensor estimates. \Cref{sec:memory} proves memory; \cref{sec:forget} proves forgetting from every initial state. The imported scaffold facts are isolated in \cref{prop:scaffold}. Readers interested in the stochastic mechanism can take that proposition as the cited input and read \cref{sec:rule,sec:poles,sec:sensors,sec:memory,sec:forget}.

\section{Dynamics and the fixed local rule}\label{sec:rule}

A cellular automaton on a finite alphabet $\Sigma$ and lattice $\Z^d$ is a map $F$ given by one finite neighbourhood $N_F$ and one function $f:\Sigma^{N_F}\to\Sigma$. Its uniform replacement perturbation has independent variables $E_t(w)\sim\mathrm{Bernoulli}(\eps)$ and $Z_t(w)$ uniform on $\Sigma$, at all sites and integer times. The update is
\[
 X_t(w)=\begin{cases}Z_t(w),&E_t(w)=1,\\F(X_{t-1})(w),&E_t(w)=0.\end{cases}
\]
The replacement is of the entire product symbol. All coupled trajectories use these same variables. A stationary law is fixed by this Markov kernel. Ergodicity means a unique stationary law and weak convergence to it from every initial law. Translation invariance means invariance under every spatial translation.

\subsection{The planar checker and provider threshold}

For a finite Wang set $R$, let $\mathcal L_R$ consist of pairs $(c,b)$, where $c$ is an edge-matching $R$-pattern on $\Mo$ and $b\in\{0,1\}$. These patterns are enumerated by finite local checks. Two letters at displacement $\zeta\in\Mo\setminus\{0\}$ are compatible when their $R$-blocks agree on the overlap after translation by $\zeta$ and their bits agree. Put $\mathcal A_R=\{\botv\}\cup\mathcal L_R$. A letter other than $\botv$ is live. A full legal chart means an assignment of live letters compatible in all eight Moore directions. Equivalently it is the $3\times3$ block presentation of an $R$-tiling with one global bit. The controller reads the central projected tile of a live letter.

Write a site as $(u,P)\in\Z^2\times\Z^4$. Fix
\begin{equation}\label{eq:geometry}
 m=16,\quad L=160,\quad
 K=16^4,\quad q=5K,\quad h=4K.
\end{equation}
Let $C_0=0$, $C_i=L\mathbf e_i$ for $1\le i\le4$, $S_i=C_i+\{0,\ldots,15\}^4$ and $S=\bigcup_{i=0}^4S_i$. For a data configuration $D_0$, check each provider $(u,P+s)$ by retaining its central live letter exactly when its eight Moore neighbours are live and compatible with that centre at their respective displacements. A failed check returns $\botv$. Define $F_R(D_0)(u,P)$ as a live letter occurring in at least $h$ checked providers, and $\botv$ if none does. Since $2h>q$, that letter is unique. A legal chart cloned over $P$ is fixed by $F_R$.

\subsection{Interval fields}

Set $T=500$ and $v=(32,32,32,32)$. The remaining fixed constants $C$, $B=4^k$, $s$, and $n_0$ are selected in \cref{sec:parameters}. In particular $s=\lceil3CB/4\rceil$. The helper alphabet $\HH$ has 100 letters: LOW represents $\{0\}$, UNKNOWN represents $\{0,1\}$, and 98 HIGH aliases represent $\{1\}$. Fix one representative for each of these three intervals.

The physical alphabet is
\begin{equation}\label{eq:alphabet}
 \Sigma=\mathcal A_R\times\mathcal A_R\times\HH^B.
\end{equation}
Its coordinates are data $D$, a second coordinate $H$, and fields $I_0,\ldots,I_{B-1}$. The second coordinate records the ordinary proposal and is retained in the physical symbol. All drivers read old data. At proposal time $t$ put
\begin{equation}\label{eq:drivers}
 A_t=F_R(D_{t-1}),\qquad
 Q_t(u,P)=F_R^T(D_{t-1})(u,P-(T-1)v).
\end{equation}
The chart alphabet provides predicates $\Orig$, $\End$, a row flag $H_{\rm row}$, and a lane phase $p\in\Z/B\Z$, defined in \cref{sec:chart}. Let $\mathbf e_x,\mathbf e_y$ denote the planar unit vectors; a provider vector in a site expression changes only $P$.

For a physical site $w$, field $j$ always uses the base $z=w-j\mathbf e_y$. Its provisional interval is
\begin{equation}\label{eq:helper}
 J_{j,t}(w)=
 \begin{cases}
 \{0,1\},&A_t(z)=\botv,\\
 \{0\}\cup I_{j,t-1}(w-\mathbf e_x+v),&A_t(z)\ne\botv,\ Q_t(z)=\botv,\\
 \{0\},&A_t(z),Q_t(z)\ne\botv,\ \Orig(Q_t(z)),\\
 I_{j,t-1}(w-\mathbf e_x+v),&\text{otherwise}.
 \end{cases}
\end{equation}
Intervals on the right are interpreted through their represented subsets. At a base $w$, let $h_w$ count the provisional intervals $J_{j,t}(w+j\mathbf e_y)$ containing $1$, and $l_w$ those equal to $\{1\}$. All of them use the same $A_t(w),Q_t(w)$. Define
\begin{equation}\label{eq:pin}
\begin{split}
 \Pin_t(w)=1\quad\Longleftrightarrow\quad
 &Q_t(w)=\botv\quad\text{or}\\
 &\bigl(Q_t(w)\ne\botv,\ p(Q_t(w))=0,\ \End(Q_t(w)),\ \text{and}\\
 &\quad\bigl[(H_{\rm row}(Q_t(w))=1\ \&\ h_w\ge s)\ \text{or}\\
 &\qquad\ (H_{\rm row}(Q_t(w))=0\ \&\ l_w<s)\bigr]\bigr).
\end{split}
\end{equation}
The deterministic proposed symbol is
\begin{equation}\label{eq:proposal}
 \left(\ \begin{cases}\botv,&\Pin_t(w)=1,\\A_t(w),&\Pin_t(w)=0,\end{cases}
 \quad A_t(w),\quad (\operatorname{rep}J_{j,t}(w))_{j<B}\right).
\end{equation}
One uniform whole-symbol replacement follows. These cases also specify the rule on malformed configurations and on every exceptional wall: the predicates are local tile functions. Every offset and every iteration depth in \eqref{eq:drivers}--\eqref{eq:proposal} is fixed.

\begin{proposition}\label{prop:local}
\proved\ Once $R,B,s$ are fixed, \eqref{eq:proposal} defines one computably specified finite-range cellular automaton. The number of parent sites in both dependency directions together is at most
\[
 \kappa=2B(n_A+n_Q+2),\qquad n_A=9q,\qquad
 n_Q=(2T+1)^2[T(L+m-1)+1]^4.
\]
Each output uses $B$ predecessor helper sites. This bound is independent of $|R|$ and $n_0$.
\end{proposition}
\begin{proof}
An ordinary query reads at most $9q$ old data sites. In $T$ iterations the planar displacement lies in $[-T,T]^2$ and each provider coordinate in $[0,T(L+m-1)]$; recentering translates this box. At output $w$ the $A,Q$ bases are the $B$ sites $w-j\mathbf e_y$, including $w$. The gate's provisional fields use the already counted base $w$ and the $B$ predecessor helper sites $w+j\mathbf e_y-\mathbf e_x+v$, $j<B$. The helper outputs read only $w-\mathbf e_x+v$, the case $j=0$. Thus there are at most $B(n_A+n_Q)+B\le B(n_A+n_Q+2)$ parents. Translation invariance gives the same reverse degree. Alphabet enlargement changes finite lookup tables, with the checker radius remaining one.
\end{proof}

\section{Charts, frames and routing}\label{sec:chart}

\subsection{The cited scaffold}

We use the scaffold of \cite[Section 3 and Appendix A]{DMhard}, with the base alphabet, edge endpoints and substitution tables specified there. Its base substitution is Ollinger's \cite{Ollinger08}, with the ruler decoration from \cite{DMforget}. Its 276-letter decorated Wang set is denoted $\mathbb A$. For a physical coordinate $z=\delta_i+2u+\sigma$, $\sigma\in\{0,1\}$, parent rank means $\nu_2(u-\varpi_i)$, with rank infinity when $u=\varpi_i$. Its thin, thick and black components, the ruler bits $\chi_1,\chi_2$, and its quarter alphabet have the meanings of those tables. The following is the precise imported input.

\begin{proposition}[Scaffold facts]\label{prop:scaffold}
\recalled\ The scaffold has an injective expansion-two substitution, preserves legality in both directions, and every tiling has a unique desubstitution tower. Its quarter tilings have a unique quarter offset $\delta\in\{0,1\}^2$ and phase $\tau=\delta+2\varpi\in\Z_2^2$. Every phase occurs; its residue modulo $2^l$ is a local function of the tiling. For Haar-almost every phase the tiling is determined by the phase, every coordinate has finite parent rank and belongs to a finite frame.

On axis $i$ put $a_d^{(i)}=\tau_i\bmod4^{d+1}$ in $[0,4^{d+1})$. The depth-$d$ frames are
\begin{equation}\label{eq:frames}
 I_d^{(i)}(z)=
 [a_d^{(i)}+(4z+1)4^d+1,\ a_d^{(i)}+(4z+3)4^d].
\end{equation}
Same-depth frames are disjoint; frames of different depths are disjoint or strictly nested. Every deeper frame contains a frame of each smaller positive depth. The free coordinates $\mathcal F_d^{(i)}(z)$ are the interior coordinates outside all smaller-depth frames. They number $N_d=2^{d+2}-2$, and their extremes are the frame endpoints plus one and minus one.

A column of finite parent rank $r$ carries light vertical wall sites only if $r=2e-1$ and it is an endpoint of a depth-$e$ column frame; it then carries them exactly at the rows in depth-$e$ row frames, as a left or right side according as it is a left or right endpoint. The horizontal statement follows by exchanging axes. A depth-$e$ endpoint has parent rank $2e-1$ and ruler bit $1$. More generally $\chi_i$ equals parent-rank parity at finite rank.
\end{proposition}

The last paragraph uses the Ruler bits and Red paths lemmas \cite[Lemmas~A.6 and A.7]{DMhard}. Their finite-rank statements hold for all tilings. In particular the column statement assumes only the column's finite rank. The global light-wall statement in that paper assumes finite ranks on both axes; below we apply the finite-rank statements within one closed frame.

At a finite-rank row/column intersection the black type is $\mathsf A$ for equal ranks, $\mathsf B$ or $\mathsf C$ if the column rank is smaller, and $\mathsf D$ or $\mathsf E$ if it is larger. Indeed, the substitution table \cite[Table~4]{DMhard} makes child $(0,0)$ keep the black type, child $(1,1)$ have black $\mathsf A$, child $(1,0)$ have black in $\{\mathsf B,\mathsf C\}$ and child $(0,1)$ in $\{\mathsf D,\mathsf E\}$. The letter descends through children $(0,0)$ from its ancestor at the first nonzero digit pair of the parent coordinate $u-\varpi$; that pair is $(1,1)$, $(1,0)$ or $(0,1)$ according as the column rank equals, is smaller than, or is larger than the row rank.

\subsection{Rank flags and lane phase}

Add a column bit $V$ and a row bit $H_{\rm row}$ before quartering. Copy $V$ on north/south edge colours and $H_{\rm row}$ on east/west colours. Impose
\begin{equation}\label{eq:rankfilter}
 \mathsf A:\ V=H_{\rm row};\qquad
 \mathsf B,\mathsf C:\ V\ge H_{\rm row};\qquad
 \mathsf D,\mathsf E:\ V\le H_{\rm row}.
\end{equation}
Internal quarter edges compare the whole decorated parent; external edges compare its edge colours. This is a finite Wang construction.

\begin{lemma}\label{lem:cutoff}
\proved\ In any legal decorated chart the flags at finite ranks have one common nonincreasing value sequence $(b_r)_{r\ge0}$. Every finite cutoff $b_r=\ind\{r<c\}$ occurs over every typical scaffold chart.
\end{lemma}
\begin{proof}
Rows and columns carry constant bits by edge matching. Equal-rank intersections have black type $\mathsf A$, so all rank-$r$ rows and columns share one value $b_r$. Intersecting rank $r<s$ in the two possible orders and applying \eqref{eq:rankfilter} gives $b_r\ge b_s$. Conversely assign the displayed cutoff to each row and column. The three black-type rank comparisons just proved verify every local filter. All ranks are finite in a typical chart.
\end{proof}

The lane phase is
\begin{equation}\label{eq:phase}
 p(x,y)=y-\tau_2-1\pmod B.
\end{equation}
It is implemented by a higher-block Wang extension. Here is an effective construction of its lookup table. For residues modulo $B$ encode the quarter offset and $2k-1$ base substitution digits. Enumerate finite hierarchy patches of increasing radius, with the source's explicit substitution and matching constraints. For a residue $\pi$, central radius $r$ and hierarchy radius $R_{\rm h}$, let $\mathcal A(\pi,r,R_{\rm h})$ be the central words projected from legal radius-$R_{\rm h}$ hierarchy patches with the prescribed digits. These finite sets decrease with $R_{\rm h}$, and their intersections are the sector languages by compactness. Search for $(r,R_{\rm h})$ for which the sets belonging to different residues are pairwise disjoint. Each test is finite exhaustive enumeration. The search terminates: phase sectors are disjoint compact sets by the tower and local phase facts of \cref{prop:scaffold}; if every larger hierarchy patch allowed an ambiguous fixed central word, compactness would give ambiguity in two full sectors. There are finitely many central words and sector pairs. Decode a word by the residue of the unique finite approximant containing it, and give a word lying in no approximant a fixed default; words of legal charts lie in their sector's approximant. Require $p$ to agree with that table, to copy horizontally and increase by one vertically. The usual overlapping-block edge colours recode these finite constraints as a Wang set, retaining the central tile projections.

Define $\Orig$ as the intersection of a bottom light horizontal wall with a left light vertical wall, and $\End$ with a right light vertical wall. These are literal functions of a quarter tile's red atoms, quadrant and ruler bits from \cite[Definition~3.11]{DMhard}. On a finite bottom row of depth $d$, origins and ends are respectively
\[
 (L_d(z),\operatorname{Bot}_d(z')),\quad (R_d(z),\operatorname{Bot}_d(z')),
 \qquad \ell_d=R_d-L_d=2\cdot4^d-1.
\]
Here $\operatorname{Bot}_d$ denotes the lower endpoint of a row frame. Finite-rank Red paths gives this assertion even if the other phase coordinate is exceptional. At this row
\[
 p(\operatorname{Bot}_d(z'))\equiv(4z'+1)4^d\pmod{4^{\min(d+1,k)}},
\]
so a finite bottom row is a phase-zero gate row exactly when $d\ge k$. (Exceptional phases may contain further gate sites of infinite rank. The local rule acts there too, and no argument below uses them.)

\subsection{Local routing rules}

Retain the orientation/event fields of \cite[Section 4, (O1)--(E3)]{DMhard}. We give their constraints here. At each side there is an orientation pair: entries in $\{\mathrm L,\mathrm R\}$ on horizontal edges, in $\{\mathrm B,\mathrm T\}$ on vertical edges. There is also an event tag in $\{\mathrm{cell},\mathrm{wall}\}^2\cup\{*\}$. Matching edges identify the entire tags. Let $\operatorname{VW}\in\{-,\mathrm L,\mathrm R\}$ and $\operatorname{HW}\in\{-,\mathrm B,\mathrm T\}$ be the local wall predicates.

If $\operatorname{VW}=-$, the west and east orientation pairs agree. At a wall of side $a$, the plus entry on its west edge and minus entry on its east edge both equal $a$. Apply the transposed rule vertically. An event tag is a pair precisely when its orientation is $(\mathrm L,\mathrm R)$ horizontally or $(\mathrm B,\mathrm T)$ vertically, and is $*$ otherwise.

A site is a horizontal carrier when $\operatorname{VW}=-$ and the west orientation is $(\mathrm L,\mathrm R)$; define vertical carriers by transposition. An active site is both. A horizontal-only site copies both horizontal event entries; an active site puts cell in its west plus and east minus entries. At an inward-facing left wall put wall in the east minus entry, and at an inward-facing right wall put wall in the west plus entry. The vertical rule puts cell in the south plus and north minus entries of active sites, copies at vertical-only sites, and puts wall at inward-facing bottom/top walls. These are (E2)/(E3).

On carriers introduce horizontal logical colours in $\{\circ\}\cup S_P$ and vertical ones in $\{\circ\}\cup S_P^2$, where $S_P$ is the machine-content alphabet of \cref{sec:machine}. Nonhorizontal carriers have empty horizontal colours; horizontal-only sites copy those colours. Impose the transposed vertical conditions. Records are empty off active sites; at active sites they are genuine triples $(t_S,t_C,t_N)$ with $t_C\in S_P$, $t_S,t_N\in S_P\cup\{\botv\}$. Define the four flags at an active site by
\[
 f_W=[e^W_-={\rm wall}],\quad f_E=[e^E_+={\rm wall}],\quad
 f_S=[e^S_-={\rm wall}],\quad f_N=[e^N_+={\rm wall}].
\]
Require $\ell^X=\circ$ exactly when $f_X=1$, for every side $X$. The remaining constraints are
\begin{align}\label{eq:tableau}
 &t_S=\botv\ (f_S=1),\quad \ell^S=(t_S,t_C)\ (f_S=0),\nonumber\\
 &t_N=\botv\ (f_N=1),\quad \ell^N=(t_C,t_N)\ (f_N=0),\nonumber\\
 &t_C=\init(H_{\rm row},\chi_2,f_S)\ (f_W=1),\quad t_C=\ell^W\ (f_W=0),\\
 &\ell^E=f_P(t_S,t_C,t_N)\ (f_E=0),\qquad
 t_C\text{ does not carry }\halt.\nonumber
\end{align}
The read-only input tracks in $t_C$ agree with the local $H_{\rm row},\chi_2$. This retains precisely (K1)--(K3), the boundary rule (C0), and the displayed replacement for (C5). The finite alphabet $R$ is the set of tuples passing these tests, followed by the phase extension above. Each tuple test and each colour equality is effective.

\subsection{Forcing a finite board in every phase}

\begin{lemma}[Finite-board forcing]\label{lem:allphase}
\proved\ In any scaffold phase, every extension satisfying the preceding rules contains, for each $d\ge1$ and each pair of depth-$d$ frames $[L_*,R_*]$, $[\operatorname{Bot}_*,\operatorname{Top}_*]$, the board
\[
 \{(x_i,y_j):0\le i,j<N_d\}
 =\mathcal F_d^{(1)}\times\mathcal F_d^{(2)}.
\]
Here $x_0<\cdots<x_{N_d-1}$ and $y_0<\cdots<y_{N_d-1}$ list the free coordinates. Its consecutive horizontal and vertical channels are single carriers. Its boundary flags are exactly
\[
 f_W=[i=0],\quad f_E=[i=N_d-1],\quad
 f_S=[j=0],\quad f_N=[j=N_d-1].
\]
Its central contents $s_i(j)=t_C(x_i,y_j)$ obey
\begin{align*}
 s_0(j)&=\init(H_{\rm row}(y_j),\chi_2(y_j),[j=0]),\\
 s_{i+1}(j)&=f_P(s_i(j-1),s_i(j),s_i(j+1)),\quad i<N_d-1,\\
 s_i(-1)&=s_i(N_d)=\botv.
\end{align*}
Thus it contains the unique $N_d$-cell, $(N_d-1)$-step finite-tape computation.
\end{lemma}
\begin{proof}
Write $z-\delta_i=2u+\sigma$ for a physical coordinate and its quarter digit $\sigma\in\{0,1\}$. Its parent rank is $\nu_2(u-\varpi_i)$. If that rank is at least $2d+1$, including infinity, then
\[
 z-\tau_i=2(u-\varpi_i)+\sigma\equiv\sigma\pmod{4^{d+1}}.
\]
Every coordinate in the closed frame has residue in $[4^d+1,3\cdot4^d]$, which contains neither $0$ nor $1$. All coordinates of the closed rectangle therefore have finite parent rank at most $2d$. Equivalently the rectangle sits strictly inside one quartered level-$(2d+1)$ supertile. We may apply the finite-rank wall classification at every site of its walls and channels.

Fix a free row $y_j$. The only vertical light walls on $[L_*,R_*]\times\{y_j\}$ are its two endpoints. Shallower walls cannot occur because $y_j$ lies in no shallower row frame. Same-depth column frames are disjoint. A deeper endpoint cannot lie in the shallower frame by strict nesting. Orientation matching across the open strip is constant; the left wall fixes its minus entry to L and the right wall its plus entry to R. The strip is therefore horizontally carried with orientation $(\mathrm L,\mathrm R)$.

We must determine its vertical carriers as well. For $z\in(L_*,R_*)$ let
\[
 D_z=\{e:1\le e\le d,\ z\text{ lies in a depth-}e\text{ column frame}\},
 \qquad e_z=\min D_z.
\]
This set contains $d$. Consider all depth-$e$ row frames with $e\in D_z$ meeting $[\operatorname{Bot}_*,\operatorname{Top}_*]$. They form a finite strictly laminar family containing $[\operatorname{Bot}_*,\operatorname{Top}_*]$. By the wall classification their endpoints are exactly the horizontal walls on this strip. Its minimal members are exactly its depth-$e_z$ members: every deeper member contains a depth-$e_z$ frame, and none of smaller eligible depth exists. Consecutive endpoints labelled B then T bound exactly a minimal member. Indeed, if the interval starting at the earlier endpoint ended later than the second, the interval ending at the second would have to start earlier than the first, contradicting disjointness or strict nesting. The converse follows because a minimal interval contains no endpoint in its interior. The vertical orientation rule forces each gap's pair from its two endpoint labels. Hence vertical carriers are exactly the interiors of those minimal frames.

If $z$ is free at depth $d$, then $e_z=d$ and every $(z,y_j)$ is a vertical carrier. If $z$ is not free, then $e_z<d$, while $y_j$ lies in no depth-$e_z$ frame, so $(z,y_j)$ is horizontal only. This proves that the active sites on each free row are exactly the $x_i$, with horizontal-only interiors between consecutive ones. Transposing proves the vertical-only assertion. The extreme free coordinates are $L_*+1,R_*-1,\operatorname{Bot}_*+1,\operatorname{Top}_*-1$ by \cref{prop:scaffold}; their adjacent exterior sites are the four light walls.

All horizontal event tags in this strip are pairs. The left wall sets the inward minus entry to wall; the right wall sets the inward plus entry to wall. Each active site supplies the facing cell entries, and the single-carrier channels copy both entries. The successive gap pairs, including gaps consisting of a single shared edge, are
\[
 ({\rm wall},{\rm cell}),\qquad ({\rm cell},{\rm cell}),\qquad
 ({\rm cell},{\rm wall}).
\]
Thus the west/east flags are exactly the two extreme-index indicators. The transposed event argument proves the south/north flags.

Horizontal-only sites transmit horizontal colours and have empty vertical colours and records. The left/right walls are not horizontal carriers, so their horizontal colours are empty. Consequently consecutive logical cells share one horizontal colour and the two exterior faces are empty. Apply the same argument to vertical channels and bottom/top walls. The boundary rule and the exact flags now force every internal logical face to be nonempty and every outer one empty. Every active site has a triple record. Exterior choices cannot disconnect or inject records into this network.

The shared vertical colour at $(i,j)$ and $(i,j+1)$ has descriptions
\[
 (s_i(j),t_N(i,j))=(t_S(i,j+1),s_i(j+1)).
\]
Thus $t_N(i,j)=s_i(j+1)$ and $t_S(i,j+1)=s_i(j)$, with bottom symbols on the outer boundaries. At $i=0$ the west flag initializes the contents. At $i<N_d-1$ the east colour is $f_P(t_S,t_C,t_N)$ and is passed to the next cell's central content. Substitution gives the stated recursion. It determines each time column, then all remaining record entries and colours. No successor is demanded on the last east boundary. The machine conversion in \cref{lem:machine} identifies this recursion with the one-head computation and sends attempted exits to $\fail$.
\end{proof}

\section{The computation carried by a chart}\label{sec:machine}

Take $U$ to be a fixed deterministic one-way-tape universal machine on positive unary inputs, with finite tape alphabet $\Gamma_U$ containing blank and $1$. Its finite transition table is part of the specification. Absorbing states $\fail,\halt$ are distinct from its working states. The integer $n_0$ is fixed by \cref{sec:parameters}.

\subsection{A direct finite machine}

The machine $P$ has read-only bits $h_{\rm in},\pi,z$, writable bits $p_*,e_*$, a writable track $n_*\in\{0,1,2\}$ and a writable track $u_*\in\Gamma_U$. Initially $h_{\rm in}=H_{\rm row}$, $\pi=\chi_2$, $z=[j=0]$, the work tracks are zero or blank, and the sole head is at cell zero in SCAN. A transition preserves every track not explicitly written. All unspecified cases and every attempted tape exit go to absorbing $\fail$. Right means increasing logical tape index $j$ (north). The following clauses are a finite transition prescription. A return to the origin with continuation $Z$ means state $\mathrm{RET}_Z$, moving left while $z=0$ and staying with state $Z$ at $z=1$.

\begin{enumerate}
\item In SCAN, on $h_{\rm in}=1$ write $p_*=1$ and move right in SCAN. On $h_{\rm in}=0,\pi=0$ fail. On $h_{\rm in}=0,\pi=1$ write $p_*=1$, move right in PLUS. PLUS writes $p_*=1$, moves right in SETEND. SETEND writes $e_*=1$, moves left, and returns with continuation $R_{00}$.
\item States $R_{bc}$ have $b\in\{0,1\}$, $c\in\{0,1,2\}$. At $e_*=0$, a zero $p_*$ is skipped. A one $p_*$ is retained if $b=0$ and erased if $b=1$; move right, flip $b$ and replace $c$ by $\min(c+1,2)$. At $e_*=1$ apply these clauses in order: $c=0$ fails; $c=1$ returns with continuation $\mathrm{SKIP}_0$; $c=2,b=1$ fails; $c=2,b=0$ returns with continuation APPENDN.
\item APPENDN scans right over $n_*=1$. At the initial zero it writes one and returns with continuation $R_{00}$; a two fails. Thus every even halving appends one token.
\item For $0\le c<n_0+1$, $\mathrm{SKIP}_c$ requires $n_*=1$, changes it to two, and moves right to $\mathrm{SKIP}_{c+1}$. State $\mathrm{SKIP}_{n_0+1}$ requires another one and returns with continuation FINDN.
\item FINDN scans over twos. On a one it writes two and returns with continuation APPENDU. On a zero it returns with continuation $\mathrm{RUN}_{q_{U,0}}$. APPENDU scans right over $u_*=1$, writes one at the initial blank and returns with continuation FINDN; every other tape symbol fails.
\item $\mathrm{RUN}_q$ executes the literal transition table of $U$ on $u_*$, preserving the other tracks. A halting successor, or an initially halting state, enters $\halt$. Both $\halt$ and $\fail$ stay forever.
\end{enumerate}

Let $S_P$ be the resulting finite alphabet of tape symbols with one state or a no-head mark. Define $f_P(a,b,c)$ by this ordered rule: if $b$ carries a head, perform its transition, retaining a failing head if the move points to a $\botv$ neighbour (a tape end); otherwise retain its new head for a stay move or drop it for a move. If $b$ has no head but $a$ moves a head north, retain $b$'s symbol and install that new head state. Otherwise do the same for a head from $c$ moving south. In the remaining case preserve $b$ without a head. This is total even on multiple-head triples, with the stated priority.

\begin{lemma}\label{lem:machine}
\proved\ On a finite tape initialized above, $f_P$ gives exactly the sequential $P$ computation, with one head at every time and absorbing failure on an attempted exit. If the initial zero of $h_{\rm in}$ has index $j_*$, $P$ can begin simulating $U$ only if its parity bit is one and $j_*+2=2^e$ with $e>n_0+1$. It then runs exactly $U(e-n_0-1)$.
\end{lemma}
\begin{proof}
A stay step changes just the head cell. For an interior move the old head cell loses the head and exactly its destination receives the new state, retaining the destination's own tape symbol. Every other cell preserves its symbol. An exit leaves the head at the boundary in $\fail$. Induction gives the one-head assertion and excludes the malformed priorities.

The scan makes exactly $j_*+2$ marks, with $e_*$ on the next cell. Each $R$ pass keeps the odd-numbered surviving marks. Its $c$ records their original count capped at two, and $b$ its parity. The ordered termination clauses accept a singleton before rejecting an odd number exceeding one. Consequently an original power $2^e$ has exactly $e$ successful even passes; any other positive count eventually fails. APPENDN therefore writes $e$ tokens. SKIP consumes $n_0+1$ of them and requires a further token. FINDN and APPENDU copy precisely the remaining $e-n_0-1$ tokens to the initially blank $u_*$ track. Only RUN can enter $\halt$, and it executes $U$ on that unary input. Any shortage of tape during any routine fails.
\end{proof}

\subsection{Which input is read}

\begin{lemma}\label{lem:input}
\proved\ Suppose $a\ge2$ and the finite-rank flags satisfy $b_{2a-3}=1$, $b_{2a-1}=0$. Every depth-$d$ board with $d>a$ has its initial zero at logical index $2^{a+1}-2$, with odd parent rank $2a-1$. A board with $d\le a$ either has no zero before its boundary or has a zero with even rank before any odd-rank zero. Thus it cannot simulate a different input.
\end{lemma}
\begin{proof}
By \cref{lem:cutoff} the initial zero rank is either $2a-2$ or $2a-1$. In a depth-$d$ frame of phase $\tau_i$, write a free coordinate as $a_d^{(i)}+(4z+1)4^d+t$, starting with $t=2$. In its left half, before $t=4^d$, the parent rank is $\nu_2(\lfloor t/2\rfloor)$. An even rank $2e$, $1\le e<d$, gives residues $2\cdot4^e$ or $2\cdot4^e+1$ inside a smaller excluded frame. Hence noncentral free ranks are zero or odd. The two central coordinates have rank $2d$.

For $d>a$, the initial free coordinate of rank at least $2a-2$ is $t=4^a$, of rank $2a-1$. To count its index let $A_a$ be the residues in $[0,4^a)$ excluding every smaller-frame window. There are four at $a=1$. The four next-digit blocks contribute $|A_a|,1,0,|A_a|-1$, so $|A_a|=2^{a+1}$. Allowed coordinates in $[2,4^a]$ remove zero and one and add $4^a$, giving $2^{a+1}-1$ coordinates. The initial zero index is one less.

If $d<a-1$, all board ranks are below the cutoff. If $d=a-1$ an initial zero can occur only at the even central rank $2a-2$; if $d=a$ it again occurs at the even centre. The parity guard rejects these cases. This uses only relative finite residues and holds in every phase.
\end{proof}

\begin{proposition}[The computational chart]\label{prop:chart}
\proved\ For $a=n+n_0$, if $U(n)$ does not halt there is a spatially translation-invariant law of legal charts with either global bit and with row flag one on finite bottom depths $d<a$ and zero on depths $d\ge a$. If $U(n)$ halts, every legal chart in every phase has a depth-$(a-1)$ bottom end with flag zero or a depth-$a$ bottom end with flag one.
\end{proposition}
\begin{proof}
For nonhalting choose a typical phase and cutoff $2a-2$ or $2a-1$. In a typical phase, \cite[Propositions~5.1 and 5.7 and Lemma~5.3]{DMhard} give a legal assignment of the orientation and event tags, with active sites exactly the grids $\mathcal F_d^{(1)}\times\mathcal F_d^{(2)}$, each active site in one grid, and the boundary flags of \cref{lem:allphase}. These assignments use only (O1)--(E3), which we retain. Fill every board with its initialized finite run. Large boards decode $n$ by \cref{lem:input,lem:machine}; small boards fail or reach a boundary. A finite exit fails, so none contains $\halt$. Fill channel colours from their board and put empty records elsewhere. These assignments satisfy all tuple filters, and the phase extension is unique. Both global bits are compatible. Translating and taking a weak limit of spatial averages gives an invariant law. Its phase is Haar because each finite residue law is invariant under the transitive integer translations of that residue torus. Hence it gives full measure to typical phases and retains the cutoff.

For halting, avoiding both asserted wrong flags would impose the two rank conditions of \cref{lem:input}, since every finite depth has bottom corners in every phase. Every $d>a$ board would decode $n$. The decoding and the halting computation have finite time and space budgets. Choose $d$ with $N_d$ larger than the elapsed time and largest visited address. \Cref{lem:allphase} forces its computation, including $\halt$ at a stored time column, contradicting \eqref{eq:tableau}. Both critical depths are gated since $a-1\ge k$.
\end{proof}

\section{Contour estimates for the data checker}\label{sec:poles}

\subsection{Finite supports and charge arms}

The charge-arm form of Toom's contour method is described by G\'acs \cite{Gacs21}, building on the proof of Berman and Simon \cite{BermanSimon88}. We prove the finite estimate used here. A two-cluster support consists of one selected parent in each of two distinct provider clusters, with a Moore planar displacement allowed. For each charge it provides a good parent as follows: charge zero chooses a nonzero cluster; charge $i\ge1$ chooses a cluster other than $i$. For backward physical steps define
\[
 \psi_0(P,t)=5\sum_iP_i+640t,\qquad
 \psi_i(P,t)=-5P_i-160t\quad(1\le i\le4).
\]
Their sum is zero. On $P\mapsto P+s$, $t\mapsto t-1$ the good gains are at least $160$ for charge zero and
\begin{equation}\label{eq:poleconstants}
 \varrho_+=160-5\cdot15=85
\end{equation}
for the others. Every gain is at least $-\varrho_-$, where $\varrho_-=4\cdot160+5\cdot15=715$. Put
\[
 A_0=5(1+\varrho_-/\varrho_+)=4000/85<48.
\]
Planar displacements do not affect these potentials.

\begin{lemma}[Grouped terminal estimate]\label{lem:pole}
\proved\ In a finite root-reachable acyclic support graph, let every nonterminal have two selected supports and a good support for each of the five charges. Partition its terminal leaves into groups. Suppose the sum of the five potential oscillations within any group is at most $\Omega$. One can select $g$ distinct groups and a connected copied-arm witness containing the root with
\begin{equation}\label{eq:polebound}
 e\le A_0(g-1)+\Omega g/\varrho_+
\end{equation}
positive support edges, counted with multiplicity. Unary zero-potential links are permitted; they create no additional charge sources and do not enter $e$. With singleton groups, a root witness that contains a positive support edge has at least two terminal flags. If support edges decrease time by one, and unary links and group identifications preserve time, the span is at most $e$.
\end{lemma}
\begin{proof}
Give each vertex its set of descendant terminal-group names. Expose the graph in topological order, replacing a frontier nonterminal by its supports. Join frontier vertices when their descendant sets overlap. Replacing a binary vertex splits its component into at most two: each remaining vertex remains linked to one support by the old shortest overlap chain. It joins no outside component because support descendant sets are subsets of the removed set. A unary replacement cannot split a component. At the end each component is one terminal group.

A copied-arm witness is the graph obtained by taking separate copies of the arm interiors, identifying their sources and identifying terminal leaves in the same group. Start five good arms from the root. Maintain exactly one arm endpoint of each charge in each frontier overlap component. When a component splits, add a source at the split vertex. For each charge choose its initial support in the new component lacking that charge endpoint, then follow good supports. The split leaves each charge's old endpoint in one of the two components, and the new source supplies its missing endpoint in the other; a unary exposure moves an endpoint within its component. Only this initial edge of a nonroot source can be bad. Copy arm interiors, merge at sources and merge leaves by group name, and take the connected component containing the root. For each charge its sources and terminal groups are in bijection. For each charge, following its arms pairs every source with exactly one terminal group, and every selected group has exactly one such endpoint; thus the root component has $g$ sources and $g$ groups, with one incidence of each charge at each.

At most $5(g-1)$ edges are bad. Summing charge increments, the source sums vanish. Comparing each group's charge endpoints to one reference leaf bounds the total terminal sum by $g\Omega$. Thus
\[
 \varrho_+ e-(\varrho_++\varrho_-)5(g-1)\le g\Omega,
\]
which proves \eqref{eq:polebound}. If the root witness contains a positive support edge, then $e>0$, so singleton groups require $g\ge2$. Under the stated clock hypotheses a connected path from the root to any vertex changes time only on positive support edges. A simple such path uses at most $e$ of these edges; unary links and group identifications change time by zero. This proves the span assertion. The argument for the edge count used acyclicity and potential oscillations, so it also applies to groups at unequal times when their potentials agree.
\end{proof}

\subsection{The raw marked set and coherence}

Let $\MM^{t_0}$ start empty at physical time $t_0$. Mark each replacement flag and each site having a marked Moore parent in two distinct provider clusters. Let $\MM$ be its increasing limit as $t_0\to-\infty$. Membership has a finite derivation. This is a set determined by flags alone.

\begin{lemma}\label{lem:rawpaths}
\proved\ For the threshold in \cref{sec:parameters}, almost surely no infinite backward ordinary dependency path is contained in $\MM$. The assertion holds simultaneously at every root.
\end{lemma}
\begin{proof}
A marked point has a contour by \cref{lem:pole}, with $f$ distinct flags, at most $50(f-1)$ arm edges and span at most that number. A terminal root has $f=1,e=0$. Along a marked path attach a witness at its current root and restart one step before its earliest terminal. Selected time slabs are disjoint, so their flags are distinct. At most $50(f-1)+1\le50f$ path steps are covered by a witness; together its support and path edges cost at most $101f$ edges. The graph has fixed degree at most the $\Delta$ of \cref{sec:parameters}. A doubled-edge traversal and a terminal subset encode at most $K_{\rm mem}^f$ choices by \cref{lem:counts}. Their probability is $\eps^f$. A path of length $n$ forces $f\ge n/50$, and the geometric tail tends to zero. Countability handles all roots.
\end{proof}

\begin{lemma}[Coherence of unmarked live letters]\label{lem:coherence}
\proved\ In every complete trajectory of \eqref{eq:proposal}, all unmarked live data or second-coordinate readings at a fixed planar site and time have the same letter over all provider positions. Readings at neighbouring planar sites are compatible in all eight Moore directions. If any unmarked live reading exists, one full legal, time-constant planar chart agrees with every such reading at every time.
\end{lemma}
\begin{proof}
Off a replacement, data is an erasure of $F_R(D)$ and the second coordinate equals $F_R(D)$. An unmarked output has marked Moore parents in at most one cluster, hence at most $K$ affected provider offsets. A live output has at most $K=q-h$ failed checked providers.

For two roots at one planar position, reduce their provider $\ell^1$ separation by one coordinate step. Within each cluster pair offsets with the corresponding unit shift; there are $5K(1-1/16)$ injective pairs. For each root discard at most $K$ failed supports and $K$ marked-parent offsets; injectivity means each discard removes at most one candidate pair. At least
\begin{equation}\label{eq:spancount}
 [5-4-5/16]K=45056>0
\end{equation}
remain. Each remaining pair gives unmarked live data parents with the respective retained labels and a provider separation smaller by one. At zero separation data/second-coordinate agreement follows from the common proposal off flags; for mixed roots use one additional update before applying the same argument. Induction on the separation proves equal labels. For a Moore-neighbour pair at the same provider position, the two checked-provider sets overlap in at least $2h-q=3K$ offsets. At a common offset one check certifies that its retained centre letter is compatible with the second retained centre letter. This is the zero-provider-span compatibility base, including all four diagonals. For different provider positions, the same injective pairing and discards give unmarked DATA parents at the two planar sites with smaller provider span. Induction therefore proves compatibility at every span. Using $\MM^{t_0}$, the induction applies once elapsed time exceeds the span plus one. In a complete trajectory send $t_0$ to $-\infty$; unmarked parents remain unmarked at every finite stage.

Let $V_t$ contain the planar sites with an unmarked live reading. Among $h$ checked supports at least $h-K$ have no marked Moore parent. One supplies unmarked data at all nine planar sites at time $t-1$, retaining the central label. Therefore
\[
 V_t+\Mo\subseteq V_{t-1},
\]
and labels persist backwards on overlaps. A nonempty $V_t$ expands to cover every planar site in the union over earlier times. The ancient union $\bigcup_{s\le t}V_s$ is therefore the whole plane. The nested domains and persistent labels define a full chart; each compatibility relation is checked at a common sufficiently early time. Future domains use those same labels. All statements are deterministic on the common flag set, so they apply also to noncausally selected complete trajectories.
\end{proof}

\begin{lemma}[Deep wrong readings]\label{lem:deep}
\proved\ Relative to the chart in \cref{lem:coherence}, a wrong live depth-$T$ query in \eqref{eq:drivers} has a finite contour with at least eleven distinct raw flags and at most $50(f-1)$ support edges. All its flags are no later than the queried input time. At lag $\ell\le W$ from a current output time $t$, its earliest flag time $t_0$ satisfies $t-t_0\le W+50(f-1)$.
\end{lemma}
\begin{proof}
Every wrong live actual reading is marked by \cref{lem:coherence}. A wrong live forecast retains its letter at at least $h$ center providers. At each of its $T$ virtual levels these wrong centers occupy two distinct clusters; append their finite marked derivations. Give a virtual $F_R^j(D_s)(u,P)$ the provider characteristic coordinate $P+vs+jv$, and actual $D_s(u,P)$ the coordinate $P+vs$. Every virtual support has displacement $s'-v$, $s'\in S$, as does the final step to actual data; its potentials are the five spatial potentials $5\sum \chi_{P,i},-5\chi_{P,i}$. The gains are \eqref{eq:poleconstants} and $-715$. Recentring in \eqref{eq:drivers} places the query root at the current characteristic coordinate.

Apply \cref{lem:pole}. Every root-to-flag arm crosses all $T$ virtual levels, so $e\ge T=500$, and $e\le50(f-1)$ gives $f\ge11$. Virtual levels precede no actual input flag in physical time; their bookkeeping levels account for part of $e$. For the span use formal time $s+j$ at virtual level $j$ over $D_s$, and physical time on actual data and flags. Each support edge drops this clock by one and there are no unary links. Recentring puts the query root on the output characteristic, so $t-t_0\le\ell+1-T+e\le\ell+e$. This gives the stated bound.
\end{proof}

\section{Noise sensors and their joint source law}\label{sec:sensors}

Fix a legal reference chart and drive ideal interval fields by its exact instructions. A true origin resets to zero; other sites copy. The replacement mark is uniform on $\HH$ because \eqref{eq:alphabet} is a product, even though the fields at a physical site share their flag. The upper scalar mark is zero for LOW and one for UNKNOWN or HIGH; the lower scalar mark is one for HIGH and zero otherwise.

The true ride of a finite ruler is the segment of its base characteristic from the origin output to the output preceding its end. On a finite ruler of depth $d$, the provisional end reads exactly $\ell_d$ update atoms, from the origin output through the output preceding the end. The end's own replacement is later than its test. Thus the ideal upper success probability is
\begin{equation}\label{eq:upperlaw}
 \bar p_d=.99[1-(1-\eps)^{\ell_d}].
\end{equation}

For a recent window of $W$ updates a lane bit is one exactly when its most recent replacement mark is a HIGH alias; it is zero for LOW, UNKNOWN, or no replacement. Its probability is
\begin{equation}\label{eq:lowerlaw}
 p_W=.98[1-(1-\eps)^W].
\end{equation}
A short bad event means at least $s$ upper successes. A long bad event means fewer than $s$ successes in this recent-window experiment.

\begin{lemma}[Ownership of noise atoms]\label{lem:ownership}
\proved\ For one fixed chart phase, the intrinsic source supports of distinct finite phase-zero gates are disjoint, over all space and time. Their $B$ lanes are independent. The same assertion holds if a long support is restricted to its recent $W$ atoms.
\end{lemma}
\begin{proof}
A phase-zero row $y$ owns the physical rows $y,\ldots,y+B-1$, using field $j$ at row $y+j$. These blocks partition the rows. Thus a physical update atom is relevant to one field and one base row. Within a lane the invariant coordinates of a backward copy are
\begin{equation}\label{eq:characteristic}
 \chi_{\rm ch}=(x-t,y,P+vt).
\end{equation}
For fixed $\chi_{\rm ch}$ the end position determines the time and provider location. Its finite source segment is $[L_d,R_d-1]$ in the $x$ coordinate. Successive true ruler intervals on that row are disjoint. Hence each atom belongs to at most one gate's source. Restricting supports preserves disjointness.
\end{proof}

For a selected input $n\ge1$ put $a=n+n_0$ and
\begin{equation}\label{eq:windows}
 I_n=\left(\frac{2C}{\ell_a},\frac{21C}{10\ell_a}\right),
 \qquad W=\lfloor C/\eps\rfloor\quad(\eps\in I_n).
\end{equation}
The parameters guarantee $a\ge k+3\ge4$, $\eps\le C/100$, $0<C<.01$, and $B\ge20/C$. In particular
\begin{equation}\label{eq:rulermargins}
 \ell_{a-1}<W,\qquad \ell_a\ge2W.
\end{equation}
$W$ is used only in the analysis; the automaton does not store it.

\begin{lemma}[Intrinsic tails]\label{lem:tails}
\proved\ Define the rational numbers
\[
 u_B=2^{-\lfloor3BC/100\rfloor},\qquad
 z_B=2^{-\lfloor BC/3\rfloor},\qquad
 w_B=2^{-\lfloor BC^2/50\rfloor},\qquad
 \widehat\eta_B=\frac{3u_B}{1-4z_B}+2w_B.
\]
For the selected $B$, $4z_B<1$. In any $W$-time bin on one base characteristic, the probability that some finite gate has an intrinsic bad event is at most $\widehat\eta_B$. At a single critical gate of depth $a-1$ or $a$ the corresponding probability $\rho_B$ is also at most $\widehat\eta_B$.
\end{lemma}
\begin{proof}
For $d=a-1-r_{\rm off}$ the conservative support bound is
\begin{align*}
 \eps(\ell_d+1)
 &<\frac{21}{40}C4^{-r_{\rm off}}\frac{2\cdot4^a}{2\cdot4^a-1}\\
 &\le\frac{21}{40}\frac{128}{127}C4^{-r_{\rm off}}<.53C4^{-r_{\rm off}}.
\end{align*}
The displayed rounding estimate uses $a\ge3$. The exact $\ell_d$-atom bound is smaller. For independent Bernoulli variables of mean at most $p_{r_{\rm off}}=.53C4^{-r_{\rm off}}$, the exponential moment inequality $1-p+pe^\vartheta\le\exp(p(e^\vartheta-1))$ and the Chernoff method \cite{Chernoff52} give
\[
 \Prob\{\text{at least }B\theta\text{ successes}\}
 \le\exp[-B(\theta\log(\theta/p_{r_{\rm off}})-\theta+p_{r_{\rm off}})],
 \qquad\theta=.75C.
\]
This derivation also applies to the higher integer threshold $s$. At $r_{\rm off}=0$ its exponent exceeds $.03BC$: $\log(75/53)\ge2(75-53)/(75+53)=11/32$ and $.75(11/32)-.22>.03$. The exponent increases by at least $BC/3$ per rank step because $.75\log4-.3975>1/3$. At most $3\cdot4^{r_{\rm off}}$ short ends occur in a bin, since their spacing is $4^{d+1}$ and $W<4^a$. The sum is at most $3u_B/(1-4z_B)$, using $e^{-x}\le2^{-\lfloor x\rfloor}$.

For a long ruler, $\eps W\ge C-\eps\ge.99C$, so \eqref{eq:lowerlaw} gives
\[
 p_W\ge.98(1-e^{-.99C})>.9C.
\]
The last inequality follows from $1-e^{-x}\ge x-x^2/2$ and $C\le.01$. Also $s/B\le.75C+1/B\le.8C$. Hoeffding's inequality \cite[Theorem~2]{Hoeffding63} gives the lower-tail bound $e^{-2B(.1C)^2}=e^{-BC^2/50}$. At most two long ends occur in a bin. Their contribution is at most $2w_B$. This proves both assertions.
\end{proof}

Partition each characteristic \eqref{eq:characteristic} into bins $i_{\rm bin}W\le t<(i_{\rm bin}+1)W$. Write $R(\chi_{\rm ch},i_{\rm bin})$ for the event that some raw flag occurs on that characteristic in the bin. On a phase-zero base characteristic write $S(\chi_{\rm ch},i_{\rm bin})$ for some intrinsic finite-gate failure in that bin, and put $S=\varnothing$ on other rows. This gives one source label $(\chi_{\rm ch},i_{\rm bin})$ for each bin.

\begin{lemma}[Joint source bound]\label{lem:joint}
\proved\ For any $g$ distinct labels,
\begin{equation}\label{eq:joint}
 \Prob\{R\cup S\text{ at all }g\text{ labels}\}
 \le\gamma^g,\qquad \gamma=\sqrt C+\widehat\eta_B^{1/6}.
\end{equation}
\end{lemma}
\begin{proof}
Raw bins use disjoint atoms and have probability at most $\eps W\le C$. A statistical bin uses its $B$ private lane characteristics and at most its own and preceding bins. Distinct base characteristic labels have disjoint lane characteristics by \cref{lem:ownership}. On each characteristic bins in one residue class modulo three have disjoint source supports. A largest class therefore bounds an intersection of $g_S$ statistical events by $\widehat\eta_B^{g_S/3}$. For $g_R$ distinct raw events, Cauchy--Schwarz gives the mixed bound $C^{g_R/2}\widehat\eta_B^{g_S/6}$, without assuming that raw and statistical events are independent. Sum over the two source-type choices at each label to obtain \eqref{eq:joint}.
\end{proof}

\section{Memory when the computation does not halt}\label{sec:memory}

Fix a correct cutoff chart from \cref{prop:chart} and one of its two global bit lifts. Start data and the second coordinate at that chart, cloned over providers. Initialize each helper by its ideal chart-driven copying process with independent negative-time noise. Since every point of a lane has an earlier reset or an earlier replacement almost surely, this stationary initialization exists. Spatially average the joint chart/helper law; its phase is Haar and typical almost surely. The positive-time noise is independent of this initial law.

Call a data value or driver bad when it differs from the reference live letter, including when it is bottom. If an ordinary forecast is bad, at least $q-h+1=K+1$ checked providers fail to give the correct letter. Each such check has a bad Moore input. They occupy at least two clusters. Thus a bad forecast admits a two-cluster bad support. The same assertion holds at each virtual iteration of a bad deep driver.

\begin{lemma}[Error decomposition]\label{lem:decomposition}
\proved\ Every finite-time data or second-coordinate error has a finite acyclic derivation whose terminal sources are raw or statistical bins, whose branching steps are ordinary or virtual two-cluster supports, and whose remaining links follow one base characteristic for a lag at most $W$. Each such characteristic link is followed by a positive support step before another source-free link is needed.
\end{lemma}
\begin{proof}
A replaced output is a raw source. At an unflagged output, a bad ordinary $A_t$ supplies its two-cluster support. If $A_t$ is correct, the second coordinate is correct. A data error must then be an erasure. A bad current $Q_t$, including bottom, supplies a depth-$T$ bad support. A correct current $Q_t$ can erase only at a genuine gate of the chart.

At a short gate with correct flag one, its true origin lies less than $W$ updates earlier. If every ordinary and deep driver since that origin is correct, the reset followed by copying and raw marks gives exactly the ideal upper experiment. An erasure therefore requires its intrinsic short failure. Otherwise use a bad $A$ or $Q$ on this base characteristic as the next error. At a long gate with correct flag zero, the true origin is outside the recent $W$ window. If all drivers there are correct, each last HIGH mark forces a singleton one, whatever the older helper state. Absence of the intrinsic long failure gives at least $s$ such singleton fields, preventing an erasure. Again a false erasure supplies a recent driver error.

The $B$ physical lanes share the same driver base at every lag by \eqref{eq:helper}; there is one possible characteristic of driver errors. A bottom ordinary instruction counts as an ordinary error, and a bottom deep instruction as a deep error. Each leads to a two-cluster support. Recursing decreases physical time, or decreases the current finite virtual level before decreasing physical time. Before time zero data drivers may be taken as the perfect chart for this initialization convention. Thus the derivation is finite and acyclic, and a characteristic link always lands on a driver with a support step.
\end{proof}

\begin{lemma}[Counting the memory witnesses]\label{lem:memorycount}
\proved\ The error event at a prescribed site and time is contained in a union of at most $K_{\rm mem}^g$ events for each $g\ge1$, each requiring $g$ distinct source labels. The constants $K_{\rm mem}$ in \cref{sec:parameters} depend on geometry and $T$, and are independent of $B,W,n,n_0$ and $R$.
\end{lemma}
\begin{proof}
Use the finite derivation before any projection. For an actual $D_s(u,P)$ assign characteristic provider coordinate $P+vs$, planar coordinate $(u_x-s,u_y)$ and bin $\lfloor s/W\rfloor$. For virtual $F_R^j(D_s)(u,P)$ use $P+vs+jv$, planar coordinate $(u_x-s-1,u_y)$, that bin and its level $j$. A recentered $Q$ root has exactly its output's characteristic coordinates. Every positive support displacement in provider coordinates is $C_i+o-v$, $o\in\{0,\ldots,15\}^4$. Its planar displacement is Moore, with an additional $\mathbf e_x$ on the final step to actual data. Characteristic links have zero provider and planar characteristic displacement and bin difference at most two.

Use potentials $\psi_0=5\sum \chi_{P,i}$, $\psi_i=-5\chi_{P,i}$. Their good and worst gains are $85$ and $-715$. Terminal leaves with one source label have exactly equal potentials, even at unequal physical times. \Cref{lem:pole} gives at most $48(g-1)$ positive edges. Suppress each maximal source-free characteristic corridor to a corridor link and its entry link. By \cref{lem:decomposition}, the next nonterminal reached has a positive support step; charge the two links to that step, once along each copied arm. Thus there are at most two unary entry/corridor links per positive support edge, plus at most $5g$ terminal source-role links. The total is less than $149g<256g$.

Project the selected witness to characteristic coordinates, bin, virtual level and finite role, and take the union of repeated edges. It remains connected and keeps $g$ distinct marked source labels. Its degree is bounded by the $B$-independent $\Delta$ in \cref{lem:counts}, because the fixed helper and gate row shifts cancel. Doubled-edge traversals and source subsets give the claimed count there. The charge bound was applied before projection, so projected cycles cause no issue.
\end{proof}

\begin{theorem}[Two stationary laws]\label{thm:memory}
\proved\ If $U(n)$ does not halt, then for every $\eps\in I_n$ the perturbation has two distinct spatially translation-invariant stationary laws.
\end{theorem}
\begin{proof}
Condition on the correct chart. \Cref{lem:joint,lem:memorycount} bound every data or second-coordinate error, uniformly in time, by
\[
 \sum_{g\ge1}(K_{\rm mem}\gamma)^g=e_{\rm base}<\frac3{125},
\]
where the inequality follows from the parameter choices in \cref{sec:parameters}. Integrate over the invariant chart/helper initial law. In the bit-$b$ initialization the probability of the live data cylinder with bit $b$ is greater than $122/125$ at every site and time. The two bit cylinders are disjoint.

The configuration space is compact and the kernel is Feller, since updates are local. A weak limit of temporal Cesaro averages is stationary: the difference of an average and its one-step image is the initial law minus the final law, divided by the averaging length. The initial law and dynamics are spatially invariant, so each limit is invariant under all six spatial coordinates. The one-site cylinder inequalities persist in the limit. The limits from bits zero and one are distinct.
\end{proof}

\section{Forgetting when the computation halts}\label{sec:forget}

\subsection{A missed erasure in a complete history}

Fix $n$ with $U(n)$ halting. For this fixed instance the finite family of residues
\[
 \phi=\tau\bmod4^{a+1}\quad\text{in both planar coordinates}
\]
determines all critical finite origins, ends, their true rides, and their lane phase. For one such $\phi$ define a flag/mark event $\mathcal X^\phi$ at each critical gate: at depth $a-1$ it is the intrinsic short bad event; at depth $a$ it is the long bad event defined by \eqref{eq:lowerlaw}. It is absent elsewhere. By \cref{lem:ownership,lem:tails}, for distinct points $S_0$,
\begin{equation}\label{eq:Qlaw}
 \Prob\{\mathcal X^\phi\text{ at all points of }S_0\}\le\rho_B^{|S_0|},
 \qquad\rho_B\le\widehat\eta_B.
\end{equation}

\begin{lemma}[Missed-erasure implication]\label{lem:miss}
\proved\ In every complete trajectory with an unmarked chart of phase $\phi$, an unmarked live data output at a wrong critical end requires $\mathcal X^\phi$ at that output or a wrong live deep query at one prescribed base-characteristic lag $0\le\ell\le W$.
\end{lemma}
\begin{proof}
Such an output has live ordinary $A_t$ and no erasure. Its $Q_t$ is live because bottom would erase by default. If this live driver is wrong it is the lag-zero alternative. Otherwise it reads the chart's correct critical end and wrong flag.

At a short wrong flag zero, inspect its true origin. In the absence of a wrong live deep instruction along the ride, the actual interval contains the ideal upper bit at every step: a correct origin resets zero; bottom $Q$ inserts zero by union; bottom $A$ gives UNKNOWN; correct nonorigins copy; a raw interval mark contains its upper bit. This argument removes dependence on helpers older than the origin. With no short bad event there are fewer than $s$ upper ones. Every actual singleton one must be one of these, so $l_w<s$ and the flag-zero erasure fires.

At a long wrong flag one, absence of its bad event gives at least $s$ last HIGH marks in the recent $W$ window. Each writes singleton one. Its true origin is outside the window. Correct nonorigin copying, bottom-$Q$ union and bottom-$A$ UNKNOWN all preserve the possibility of one. In the absence of a wrong live deep instruction these $s$ fields still allow one, giving $h_w\ge s$ and an erasure.

Both contradictions prove the implication. At lag $\ell$ the nominal query is at $(w-\ell\mathbf e_x,P+\ell v,t-\ell)$, with the fixed provider recentering in \eqref{eq:drivers}. Thus all lanes give the same query location, and there are at most $W+1$ lag choices. The proof is simultaneous over complete trajectories, with no stationarity or independence of their charts assumed.
\end{proof}

\subsection{Packing the historical queries}

\begin{lemma}[Excluding unmarked liveness]\label{lem:eliminate}
\proved\ Almost surely every live data or second-coordinate reading in every complete trajectory lies in the same raw set $\MM$.
\end{lemma}
\begin{proof}
Otherwise \cref{lem:coherence} gives a full chart. \Cref{prop:chart} gives a wrong critical end $w$. Its planar site is in the ancient union, so an unmarked live data witness occurs there at some time; a second-coordinate witness supplies a data witness one update earlier. An unmarked live data root has a checked centre parent outside its sole possibly marked cluster. Continue such parents indefinitely at the same planar $w$.

At a sampled vertex of this path apply \cref{lem:miss}. If it supplies $\mathcal X^\phi$, move one path step backwards. If it supplies a wrong deep query, attach its contour from \cref{lem:deep} and restart at the path vertex one time before its earliest flag. The selected time slabs are disjoint. Let $f$ be their distinct raw flags, $m_0$ their distinct intrinsic events, and $c$ their contours. Each contour has at least eleven flags, so $c\le f/11$.

Only the next sampled endpoint of a skipped path segment needs recording. Put $V=L+m-1=175$. A gap of $h_*$ ordinary centre steps has net provider displacement in $[0,Vh_*]^4$, with at most $(Vh_*+1)^4$ possibilities. For a contour with $f_i$ flags, its lag has at most $2W$ choices and its restart gap at most $52f_iW$. Hence lag, gap and net endpoint together cost at most
\begin{equation}\label{eq:compression}
 2[54(V+1)]^5f_i^5W^6.
\end{equation}
The recentering is a fixed translation. A $\mathcal X^\phi$ step costs at most $q$ choices. Contour traversals, compositions of $f$, interleaving and source annotations are bounded in \cref{lem:counts}; the total count is at most
\[
 K_{\rm hist}^{f+m_0}W^{6c}\le
 K_{\rm hist}^{f+m_0}W^{6f/11}.
\]

The two event families can share raw atoms. H\"older's inequality and \eqref{eq:Qlaw} give the valid joint weight $\eps^{3f/4}\rho_B^{m_0/4}$. Since $W\le C/\eps$, their counted sum has bases
\begin{equation}\label{eq:histbases}
 K_{\rm hist}C^{6/11}\eps^{9/44},\qquad
 K_{\rm hist}\rho_B^{1/4},\qquad
 \frac34-\frac6{11}=\frac9{44}>0.
\end{equation}
Both bases are below $1/64$ by \cref{sec:parameters}. A path segment of length $N$ is covered only if
\[
 N\le [50+(W+1)/11]f+m_0.
\]
For fixed $W$, its length tending to infinity therefore forces $f+m_0\to\infty$. The geometric tail excludes the infinite path. Fix $\phi$ for the whole argument and then take the finite union of these null events over phases. Countability handles the possible root site, time and provider. Thus the assertion holds simultaneously over all complete histories.
\end{proof}

\subsection{Rare births of live proposals}

Define the minimal scalar field, separately at every planar $u$, by increasing limits from zero in the finite past:
\[
 W^{\rm live}_t(u,P)=E_t(u,P)\ \mathord\vee\
 \Theta_h(W^{\rm live}_{t-1})(u,P),
\]
where $\Theta_h$ means at least $h$ of the $q$ centre providers. This definition has no planar Moore neighbours. Set $N^{\rm live}_t=\Theta_h(W^{\rm live}_t)$ and $M=h-K=3K=196608$.

\begin{lemma}\label{lem:rare}
\proved\ Actual liveness in every complete trajectory is contained in $W^{\rm live}$, and a live ordinary proposal at time $t$ is contained in $N^{\rm live}_{t-1}$. The latter is dominated by a minimal two-cluster field whose seeds require at least $M$ fresh flags. For $g$ distinct seed groups its joint seed probability is at most
\begin{equation}\label{eq:groupnoise}
 (N_0\eps^M)^{g/16},\qquad
 \lambda=L+m=176,\quad N_0=\binom{(2\lambda)^4}{M}.
\end{equation}
\end{lemma}
\begin{proof}
An unflagged live output has $h$ live data centre parents. By \cref{lem:eliminate} all nodes of its complete live ancestry are in $\MM$. An infinite branch would contradict \cref{lem:rawpaths}. Finite branching then makes the ancestry finite, with flag leaves; induction places its root in $W^{\rm live}$. The same parent count places a live $A_t$ in $N^{\rm live}_{t-1}$.

Since $W^{\rm live}_t=E_t\vee N^{\rm live}_{t-1}$, a birth with at most $M-1$ fresh provider flags has at least $h-(M-1)=K+1$ providers in $N^{\rm live}_{t-1}$, occupying two clusters. Its only alternative seed is $M$ fresh flags. This recursion holds from every finite empty past and hence in the minimal limit.

Partition provider space into cubes $Y=\lambda z+[0,\lambda-1]^4$. All flag centres for a seed rooted anywhere in $Y$ lie in $\lambda z+[0,2\lambda-2]^4$. Declare one group at $(u,Y,t)$ when that box contains at least $M$ flags, and seed the entire root cube. Its probability is at most $N_0\eps^M$ by a subset union. Equal-parity cubes have disjoint enlarged boxes. Different planar sites or times also have disjoint atoms. A largest one of the sixteen parity classes gives \eqref{eq:groupnoise}. The resulting two-cluster closure is centre-only. Its terminals are these group events.
\end{proof}

\begin{lemma}[Full-state discrepancies]\label{lem:fullcoupling}
\proved\ For the parameters in \cref{sec:parameters}, any two complete full-state trajectories driven by the same noise coincide almost surely.
\end{lemma}
\begin{proof}
A flagged output is common. An unflagged differing output requires a live $A_t(w-j\mathbf e_y)$ in one run for some $0\le j<B$: otherwise data and second coordinate are bottom and every helper proposal is the same UNKNOWN by \eqref{eq:helper}. Thus every discrepancy lies in a fixed finite union of planar translates of the birth field at time $t-1$. It also has a discrepant parent in the actual neighbourhood, since equal parents give equal deterministic outputs. A discrepancy would therefore generate an infinite backward path in this shifted union.

Attach grouped contours for that rare field greedily along the path, restarting before the earliest group time. Draw the discrepancy path one time unit lower, so that a path vertex for actual time $s$ sits at $s-1$; attach the birth contour at its actual time $s-1$ by a time-preserving planar edge, without shifting the contour. Connect a path root to its selected planar target by one literal offset edge. These edges have zero pole increment. Leaves in one group share planar site, time and provider cube, with total potential oscillation at most
\[
 \Omega=2\cdot5\cdot4\cdot175=7000.
\]
By \cref{lem:pole} a contour has $e\le H_0g-A_0$ support edges. Its path coverage, including the restart, is at most $e+1\le H_0g$, since $A_0>2$. The support count is also at most $H_0g$, where $H_0=A_0+7000/85=11000/85<130$. Group identifications cost at most $5g$ star edges and shifted roots at most $g$ connectors. Thus fewer than $282g<512g$ edges suffice. The selected time slabs make all group identities distinct.

The augmented graph has fixed degree $\Delta_{\rm rare}$ as derived in \cref{lem:counts}; it includes actual dependency offsets, group stars, and the $B$ literal planar shifts. Keeping these as edge labels gives at most $K_{\rm rare}^g$ witnesses. By \cref{lem:rare} each has weight at most $(N_0\eps^M)^{g/16}$. A path length tending to infinity forces $g\to\infty$, by the span bound. The geometric series has base below $1/64$ by \cref{sec:parameters} and excludes every such infinite path, simultaneously at every root. This proves full-state equality.
\end{proof}

\begin{lemma}[Pullback and attraction]\label{lem:pullback}
\proved\ If complete common-noise trajectories are almost surely unique, then every finite-window pullback image is eventually a singleton almost surely. The PCA has a unique stationary law attracting every initial law.
\end{lemma}
\begin{proof}
For a fixed good two-sided noise realization, the images
\[
 K_j=\Psi_{-j,0}(\Sigma^{\Z^6})
\]
are nested nonempty compact sets. Their intersection consists precisely of the time-zero states of complete trajectories. To see the nontrivial direction, choose histories from successively earlier starting times and use product compactness to preserve each finite update constraint; then extend forwards. Uniqueness makes the intersection a singleton. Projecting nested compact sets onto any finite window commutes with intersection. The finite projected sets must therefore become singletons after some finite random time.

The failure event is determined by a finite noise cone and decreases to probability zero. Time homogeneity makes its probability the same all-start forward coupling bound. Existence of a stationary law follows from Feller compactness as in \cref{thm:memory}. Coupling that stationary start to any initial law bounds their finite-window total variation distance by the failure probability. This proves uniqueness and attraction.
\end{proof}

\section{Effective constants and the main theorem}\label{sec:parameters}

All searches below are over integers or rationals with decidable tests and specify finite constants.

Put
\begin{align}\label{eq:countconstants}
 \Delta&=10^6(T+2)(9q+10),&
 K_{\rm mem}&=2^{20}(2\Delta^2)^{257},\nonumber\\
 V&=175,&
 K_{\rm hist}&=2^{40}[100(V+1)]^{10}(2\Delta^2)^{100},\\
 C&=(128K_{\rm mem})^{-2}.&&\nonumber
\end{align}
Choose the least $k\ge1$ with $B=4^k\ge20/C$, $4z_B<1$, and
\begin{equation}\label{eq:Bsearch}
 \widehat\eta_B<\min\{(64K_{\rm mem})^{-6},(64K_{\rm hist})^{-4}\}.
\end{equation}
This terminates by \cref{lem:tails}, since its three exponential factors tend to zero. Set $s=\lceil3CB/4\rceil$, take $\kappa$ from \cref{prop:local}, and put
\begin{align}\label{eq:rareconstants}
 \lambda&=176,\quad M=196608,\quad N_0=\binom{(2\lambda)^4}{M},\nonumber\\
 \Delta_{\rm rare}&=10^6[\kappa+B+(2\lambda)^4+q+T+1],\\
 K_{\rm rare}&=2^{20}(2\Delta_{\rm rare}^2)^{513}.\nonumber
\end{align}

\begin{lemma}[Literal encoding bounds]\label{lem:counts}
\proved\ The constants in \eqref{eq:countconstants} and \eqref{eq:rareconstants} dominate the memory, raw-path, historical and rare-group witness counts used above.
\end{lemma}
\begin{proof}
The projected memory graph records actual data/second-coordinate roles, $T$ virtual levels, three source-hub roles and two entry/corridor roles. With neutral or one of five charge labels there are at most $6(T+7)\le32(T+2)$ roles. Support offsets have at most $9q$ Moore/provider choices for ordinary or final virtual edges and another $9q$ for nonfinal virtual edges. Entry, corridor and hub edges have zero spatial characteristic displacement. Allow ten such zero-offset types and five bin differences. Allow every destination role and both orientations. The undirected degree is at most
\[
 2\cdot32(T+2)\cdot5(18q+10)
 \le640(T+2)(9q+10)<\Delta.
\]

A connected embedded edge set with $e$ edges can be encoded by a deterministic traversal of its doubled edges, of length at most $2e$, and a marked source subset. Its count is bounded by $\Delta^{2e}2^{e+1}$. Summing $e\le256g$ gives
\[
 2\sum_{e=0}^{256g}(2\Delta^2)^e
 \le4(2\Delta^2)^{256g}\le K_{\rm mem}^g.
\]
The slack factor covers root and role conventions. This proves \cref{lem:memorycount}; the $101f$ raw-path encoding in \cref{lem:rawpaths} fits the same bound.

For a raw contour with $f$ flags, traversal, flag subset and root role cost at most $X^{51f}$ with $X=2\Delta^2$. Its duration is at most $50(f-1)$. The historical lag and restart endpoint envelope is \eqref{eq:compression}. Since $f^5\le2^{5f}$, and $H_*=100(V+1)$,
\[
 2[54(V+1)]^5f^5\le2^{6f}H_*^{5f}.
\]
Thus one attached contour costs at most $(2^6H_*^5X^{51})^fW^6$. Size compositions cost at most $2^f$; interleaving with $m_0$ intrinsic restarts at most $2^{f+m_0}$; binary endpoint/source annotations at most $2^{4(f+m_0)}$. Each intrinsic restart has $q\le\Delta\le X$ choices. Their product is bounded by $(2^{10}X)^{f+m_0}$, leaving the displayed $K_{\rm hist}^{f+m_0}W^{6c}$ as a majorant. The phase is fixed for the whole witness.

For rare groups the support and covered path counts are each at most $H_0g$; group stars and planar connectors add at most $6g$. Hence $E<(2H_0+6)g<282g<512g$. A finite role set for support, path, root, group, port, source and charges fits in 64 roles. Actual path edges have degree at most $\kappa$; ordinary supports at most $18q$ including Moore decorations; group-member edges at most $2(2\lambda)^4$; shifted-root edges at most $2B$. A supergraph connecting all roles on all these offsets has degree at most
\[
 64[2\kappa+2B+2(2\lambda)^4+18q+2(T+1)]<\Delta_{\rm rare}.
\]
Use the same doubled-edge traversal and source-subset count with $512g$ in place of $256g$. It is at most $K_{\rm rare}^g$. Every offset is a literal edge label, including group and planar shifts.
\end{proof}

Define the positive rational threshold
\begin{equation}\label{eq:epsstar}
 \begin{split}
 \eps_* =\min\{&C/100,\ (64K_{\rm mem})^{-4},\ (64K_{\rm hist})^{-44},\\
 &[64K_{\rm rare}(N_0+1)]^{-16}\}.
 \end{split}
\end{equation}
Finally choose the least $n_0\ge k+2$ such that
\begin{equation}\label{eq:n0}
 \frac{21C}{10(2\cdot4^{n_0+1}-1)}<\eps_*.
\end{equation}
Only after this choice construct $P$ and $R$.

\begin{proposition}[Parameter verification]\label{prop:parameters}
\proved\ The construction is noncircular, the searches terminate, all intervals \eqref{eq:windows} lie in $(0,\eps_*)$, and every convergent-series bound used in \cref{sec:memory,sec:forget} holds throughout those intervals.
\end{proposition}
\begin{proof}
The constants preceding $B$ depend only on the fixed geometry and $T$. The search for $B$ terminates as shown after \eqref{eq:Bsearch}. The dependency counts depend on $B$, geometry and $T$, with no palette-size dependence. Thus $\eps_*$ is determined before $n_0$ or the machine states. Search \eqref{eq:n0} terminates since its left side tends to zero. The phase decoder terminates by the compactness argument following \eqref{eq:phase}. Each step is a finite effective construction.

The memory bound follows from
\[
 K_{\rm mem}\sqrt C=1/128,\qquad
 K_{\rm mem}\widehat\eta_B^{1/6}<1/64,
\]
so $\theta:=K_{\rm mem}\gamma<3/128$ and
\begin{equation}\label{eq:baselineerror}
 e_{\rm base}:=\frac{\theta}{1-\theta}<\frac3{125}.
\end{equation}
This is the original-kernel error bound; $1/31$ below is the perturbed-kernel bound. The raw-path base $K_{\rm mem}\eps$ is small under \eqref{eq:epsstar}. The historical statistical base is below $1/64$ by \eqref{eq:Bsearch}. Its other base is at most $K_{\rm hist}C^{6/11}(64K_{\rm hist})^{-9}<1/64$. Finally, since $\eps\le1$ and $M\ge1$,
\[
 K_{\rm rare}(N_0\eps^M)^{1/16}
 \le \frac{N_0^{1/16}}{64(N_0+1)}<1/64.
\]
This also ensures that the seed weight is below one.

For $n\ge1$, \eqref{eq:n0} places $I_n$ below $\eps_*$. Also $a\ge k+3\ge4$, which supplies the $a\ge3$ premise in \cref{lem:tails}. The lower endpoint of $I_n$ gives $\ell_a>2C/\eps\ge2W$. Its upper endpoint gives $C/\eps>10\ell_a/21$. As $\ell_a=4\ell_{a-1}+3$, this last number exceeds $\ell_{a-1}+1$, hence $W>\ell_{a-1}$. Thus \eqref{eq:rulermargins} holds with integer rounding.
\end{proof}

\begin{theorem}\label{thm:main}
\proved\ For the single automaton \eqref{eq:proposal} and the rational intervals \eqref{eq:windows}, for every $n\ge1$ and every $\eps\in I_n$,
\[
 F_\eps\text{ is ergodic}\quad\Longleftrightarrow\quad U(n)\text{ halts}.
\]
If $U(n)$ does not halt there are at least two translation-invariant stationary laws. If it halts the unique stationary law attracts every initial law, and every finite-window common-noise pullback image is almost surely eventually a singleton.
\end{theorem}
\begin{proof}
The rule is fixed and computably specified by \cref{prop:local,prop:parameters}. Nonhalting gives \cref{thm:memory}. Halting gives \cref{lem:fullcoupling,lem:pullback}. The intervals are nonempty rational open intervals computable from $n$. Since $\ell_{a+1}=4\ell_a+3$, the upper endpoint for $a+1$ lies below the lower endpoint for $a$; their upper endpoints tend to zero. Thus they are pairwise disjoint and prove Theorem A.
\end{proof}

\begin{corollary}\label{cor:undecidable}
\proved\ For this fixed automaton, deciding ergodicity at a rational noise rate is undecidable, even under the promise that the rate lies in one of the intervals $I_n$.
\end{corollary}
\begin{proof}
The midpoint $\eps_n=41C/(20\ell_{n+n_0})$ is computable and rational. An algorithm deciding the promised question at $\eps_n$ would decide whether $U(n)$ halts by \cref{thm:main}.
\end{proof}

\subsection{Exponential coalescence and extremal states}

For the specified replacement maps let
\begin{equation}\label{eq:qt}
 q_t(\eps)=\Prob\{x\mapsto X_t^x(0)\text{ is not constant on }\Sigma^{\Z^6}\}.
\end{equation}
This tests all starts against one realization of the maps. Put
$W_n^*=\lfloor\ell_{n+n_0}/2\rfloor$.

\begin{theorem}[Exponential all-input coalescence]\label{thm:exponential}
\proved\ If $U(n)$ halts, there is $C_n>0$ such that, uniformly for $\eps\in I_n$ and $t\ge0$,
\[
 q_t(\eps)\le C_ne^{-c_nt},\qquad
 c_n=\frac1{4(4V+1)(W_n^*+52)}.
\]
The rate $c_n$ is total computable from $n$; a procedure computes $C_n$ and terminates on every halting input. For any finite target set $D_*$, the forward or pullback noncoalescence probability is at most $|D_*|C_ne^{-c_nt}$.
\end{theorem}
\begin{proof}
The finite motif search of \cref{lem:motif} terminates on halting inputs. The finite-history estimates \cref{lem:finitecoherence,lem:unmarkedfinite,lem:bulk,lem:finitepacking} prove the displayed bound, with the stated rate, for the existence of any two starts with different full output symbols. Their prefactors are computed from the finite motif radius and the constants of \cref{sec:parameters}. Translation and a union bound give the finite-set assertion. A forward and a pullback block have the same law. Coupling to the stationary law also gives exponential attraction on every finite set.
\end{proof}

\begin{theorem}[Continuum many extremal states]\label{thm:continuum}
\proved\ If $U(n)$ does not halt and $\eps\in I_n$, the set of all stationary laws has continuum many extreme points. The normalized space--time interaction has continuum many extremal Gibbs measures. No translation invariance is asserted for these extremes; the two translation-invariant laws of \cref{thm:main} are a separate conclusion.
\end{theorem}
\begin{proof}
The chart law of \cref{prop:chart} has a nonatomic Haar phase marginal. Unique phase determination gives uncountably many distinct correct typical charts $c$, with a fixed global bit. Condition on one such $c$, clone it over providers, and initialize helpers by its ideal negative-time process. The proof of \cref{thm:memory}, before spatial averaging, bounds a data error at every site and time by $e_{\rm base}<3/125$. A temporal Cesaro limit $\mu_c$ is stationary by the Feller property. All clopen cylinder bounds pass along one subsequence, so
\[
 \mu_c\{D(u,P)=c(u)\}\ge1-e_{\rm base}
 \quad\text{for every }(u,P).
\]
Different charts differ at some $u$. At that cylinder their laws differ by at least $1-2e_{\rm base}>119/125$. Their stationary trajectory laws are DLR states by \cite[Lemma~2.4]{DMst}, which assumes no spatial invariance. Their time-zero marginals give the same separation.

For either the stationary-law set or the DLR set, let $\mathcal K$ be the full compact convex set of states. It is metrizable; compactness follows from the Feller property in the first case and continuity of the finite-range specification in the second. If its extreme set were countable, the metrizable Choquet theorem \cite{Choquet} would express every state as a countable mixture of these extremes. Truncation, putting the tail mass on one fixed extreme, changes total variation by at most the tail mass. Rationalizing the remaining weights would give a countable total-variation-dense subset of $\mathcal K$. A separable metric space cannot contain an uncountable uniformly separated family, a contradiction.

For a compatible weak metric $d$, the nonextreme set equals
\[
 \bigcup_{j\ge1}
 \{(\nu+\rho)/2:\nu,\rho\in\mathcal K,\ d(\nu,\rho)\ge1/j\}.
\]
Each set in the union is compact. Every nonextreme point is the midpoint of a shorter nontrivial segment through it, proving the equality. Thus the extreme set is $G_\delta$ in a compact metrizable space. The perfect-set theorem \cite{Kechris} gives continuum cardinality for this uncountable Polish space. The ambient probability space has cardinality at most continuum. The stationary limits $\mu_c$ themselves need not be extreme. The stacked translation-invariant continuum of \cite[Corollary~8.17]{DMforget} is a precursor with a different invariance scope.
\end{proof}
\subsection{The fixed interaction}

\begin{gibbsB}
\proved\ There is one computably specified translation-invariant finite-range interaction on $\Z^7$, with a finite alphabet and energies in $\{0,1\}$, and computable pairwise disjoint rational open intervals $\mathcal J_n$ of inverse temperature, such that for every $\beta\in\mathcal J_n$: if $U(n)$ halts the interaction has exactly one Gibbs measure, with no invariance assumed, and its specification has arbitrary-volume exponential weak spatial mixing; if $U(n)$ does not halt it has at least two translation-invariant Gibbs measures and continuum many extremal Gibbs measures.
\end{gibbsB}
\begin{proof}
Use the automaton of \cref{thm:main}, write $N_\Sigma=|\Sigma|$, and let $N_F$ be its finite neighbourhood. At $(w,t)\in\Z^6\times\Z$ put
\begin{gather*}
 S_{w,t}=\{(w,t)\}\cup\{(w+z,t-1):z\in N_F\},\\
 \Phi_{w,t}(\sigma)=\ind\{\sigma(w,t)\ne f((\sigma(w+z,t-1))_{z\in N_F})\}.
\end{gather*}
The specification uses every factor with $S_{w,t}\cap\Lambda\ne\varnothing$, including factors of exterior children. Put $\eps(\beta)=N_\Sigma/(e^\beta+N_\Sigma-1)$. For every parent word $\eta$ and symbol $c$,
\[
 \frac{e^{-\beta\ind\{c\ne f(\eta)\}}}{1+(N_\Sigma-1)e^{-\beta}}
 =(1-\eps(\beta))\ind\{c=f(\eta)\}+\frac{\eps(\beta)}{N_\Sigma}.
\]
The parent-independent denominator cancels from every specification, giving precisely the normalized space--time interaction.

We apply \cite[Corollary~13.7]{DMst}: a PCA on $\Z^d$, $d\ge1$, must have a finite alphabet, a finite neighbourhood containing $0$, a strictly positive table, and an i.i.d.\ local-map representation whose all-input common-map coefficient satisfies $q_t\le Ce^{-\alpha t}$ for all $t\ge0$ with $C,\alpha>0$. Here $d=6$, $\Sigma$ is finite, and $N_F$ contains $0$ because $\zeta=0$, $s=0$ occurs in the ordinary driver. The table is at least $\eps(\beta)/N_\Sigma>0$. The independent $(E,Z)$ updates are a local-map representation, and \cref{thm:exponential} supplies exactly its common-map coefficient. All hypotheses hold on halting windows. The conclusion is a single DLR state $\bar\mu$ and constants $C_*,c_*>0$ such that
\begin{align}
 d_{\rm TV}((\gamma^\tau_\Lambda)_E,\bar\mu_E)
 &\le C_*\sum_{z\in E}e^{-c_*\operatorname{dist}_\infty(z,\Lambda^c)},\label{eq:gibbsconv}\\
 d_{\rm TV}((\gamma^\tau_\Lambda)_E,(\gamma^{\tau'}_\Lambda)_E)
 &\le C_*\sum_{z\in E}e^{-c_*\operatorname{dist}_\infty(z,\Lambda^c)}.\label{eq:gibbswsm}
\end{align}
After enlarging $C_*$, both hold for every finite $\Lambda$, every $E\subseteq\Lambda$ and all exteriors, including disconnected volumes and volumes with holes.

On nonhalting windows, \cref{thm:continuum} gives continuum many extremal DLR states. The two translation-invariant stationary laws of \cref{thm:main} give two fully space--time translation-invariant Gibbs measures by \cite[Lemma~2.4 and Corollary~10.4]{DMst}, the correspondence of \cite{GKLM89}. Their time-zero marginals distinguish them.

The decreasing computable map $\eps\mapsto\log(N_\Sigma/\eps-N_\Sigma+1)$ sends $I_n=(\eps_n^-,\eps_n^+)$ to an open interval with computable endpoints. Approximate these endpoints until two strictly interior rationals are certified and take the interval between them as $\mathcal J_n$. These choices are uniform in $n$ and disjoint. They give every asserted conclusion. The direct all-Gibbs precursor for varying rules is \cite[Corollary~8.8 and Remark~8.14]{DMforget}.
\end{proof}
\subsection{Field response and effective consequences}

Write a symbol as $\sigma=(D,H,I_0,\ldots,I_{B-1})$ and fix
\[
 m_{\rm bit}(\sigma)=
 \ind\{D\text{ is live with bit }1\}-
 \ind\{D\text{ is live with bit }0\}.
\]
Use counting a priori measure and a dimensionless field $h_{\rm f}$, with weights
$\exp[-\beta\sum_v\Phi_v+h_{\rm f}\sum_{v\in\Lambda}m_{\rm bit}(\sigma_v)]$.
Let $P(\beta,h_{\rm f})$ be the pressure. The subscript distinguishes this field from the checker threshold $h$.

\begin{theorem}[Field response for one fixed interaction]\label{thm:field}
\proved\ The pressure is computable for computable $\beta>0,h_{\rm f}$, is even in $h_{\rm f}$, and
\[
 P(\beta,0)=\log(1+(N_\Sigma-1)e^{-\beta})
\]
is analytic for $\beta>0$. For $\beta\in\mathcal J_n$, its field derivative at zero exists exactly when $U(n)$ halts, and then equals zero. Otherwise
\[
 \partial_{h_{\rm f}}^+P(\beta,0)>119/125,\qquad
 \partial_{h_{\rm f}}^-P(\beta,0)<-119/125,
\]
so the jump exceeds $238/125$. This conclusion uses only translation-invariant Gibbs facts.
\end{theorem}
\begin{proof}
Let $\iota$ flip the global bit in each live $D,H$ letter, fix bottom and leave helpers unchanged. Compatibility and the threshold checker are equivariant, and all chart predicates ignore that bit. Thus $F\iota=\iota F$, uniform replacement preserves the involution, $\Phi\circ\iota=\Phi$ and $m_{\rm bit}\circ\iota=-m_{\rm bit}$. Applying this bijection to free finite-volume partition functions proves evenness.

For a box $\Lambda_L$ of side $L$ in $\Z^7$, fix its exterior and first keep only factors anchored at children inside $\Lambda_L$. Each sum over a child is
\[
 \sum_c e^{-\beta\ind\{c\ne f(\eta)\}}
 =a_\beta:=1+(N_\Sigma-1)e^{-\beta},
\]
independent of its parents. Sum the latest row first, then each earlier row. Children in one row occur only in their own remaining factors, so this partition sum is exactly $a_\beta^{|\Lambda_L|}$. The DLR convention adds only $O_{R_\Phi}(L^6)$ exterior-child factors, where $R_\Phi$ bounds the support distance from each anchor. Each factor is between $e^{-\beta}$ and $1$. Their logarithmic contribution divided by $L^7$ tends uniformly to zero, proving the identity and analyticity. This is the binary-energy form of the normalized zero-pressure fact \cite[Remark~10.13]{DMst}.

For computability, use free cubes of side $L$ and put
$p_L=L^{-7}\log Z_{\Lambda_L}^{\rm free}(\beta,h_{\rm f})$. Tiling a large cube by these cubes removes crossing anchors at density at most $14R_\Phi/L$. Field terms are additive. Since $0\le\Phi\le1$,
\[
 p_L-14\beta R_\Phi/L\le P(\beta,h_{\rm f})\le p_L.
\]
Leftover boundary sites have vanishing density in the large-cube limit; this argument also gives existence of the pressure. Each $p_L$ is a finite computable sum of exponentials, and an effective upper bound on $\beta$ makes the displayed error a computable modulus.

We recall the finite-range Gibbs variational principle and Lanford--Ruelle identification \cite{LanfordRuelle69}, also used in \cite[Corollary~10.7]{DMthermal}. On the compact set $\mathcal I$ of translation-invariant laws, entropy density $\mathfrak s$ is upper semicontinuous and
\[
 P(\beta,h_{\rm f})=\max_{\nu\in\mathcal I}
 \{\mathfrak s(\nu)-\beta\nu(\Phi_0)+h_{\rm f}\nu(m_{\rm bit})\}.
\]
Maximizers are precisely the translation-invariant Gibbs states. Let $\mathcal G_{\beta,0}^{\rm inv}$ be the zero-field maximizers. For $h_{\rm f}>0$, substitution of any zero-field maximizer bounds the secant slope below by its magnetization. Conversely a maximizer $\nu_{h_{\rm f}}$ at $h_{\rm f}$ has zero-field objective at most $P(\beta,0)$, so the slope is at most $\nu_{h_{\rm f}}(m_{\rm bit})$. Pressure is $1$-Lipschitz in the field, and compactness and upper semicontinuity put every limit of these maximizers as $h_{\rm f}\downarrow0$ in $\mathcal G_{\beta,0}^{\rm inv}$. Applying the argument also to negative fields gives
\[
 \partial_{h_{\rm f}}^+P(\beta,0)
 =\max_{\nu\in\mathcal G_{\beta,0}^{\rm inv}}\nu(m_{\rm bit}),\qquad
 \partial_{h_{\rm f}}^-P(\beta,0)
 =\min_{\nu\in\mathcal G_{\beta,0}^{\rm inv}}\nu(m_{\rm bit}).
\]
Halting gives one translation-invariant Gibbs state by \cref{thm:main} and \cite[Corollary~10.4]{DMst}. Symmetry fixes it and its magnetization is zero. This uses neither \cref{thm:exponential} nor unrestricted Gibbs uniqueness. Nonhalting gives the two invariant bit states of \cref{thm:memory}, with error $e_{\rm base}<3/125$. Their trajectory laws have magnetizations at least $1-2e_{\rm base}>119/125$ and at most $-1+2e_{\rm base}<-119/125$. The slope formula proves the asserted gap. Rational midpoints of $\mathcal J_n$ reduce halting to differentiability.
\end{proof}

\begin{corollary}[Computable pressure and noncomputable response scales]\label{cor:response}
\proved\ At the rational midpoint $\beta_n$ of $\mathcal J_n$, put
$M_n=\partial_{h_{\rm f}}^+P(\beta_n,0)$. It is uniformly upper semicomputable but has no uniform computable approximation to error $1/4$. There is no total computable $g$ such that on every halting input
\[
 2^{g(n)}[P(\beta_n,2^{-g(n)})-P(\beta_n,0)]<1/2.
\]
\end{corollary}
\begin{proof}
Convexity gives $M_n=\inf_{k\ge0}2^k[P(\beta_n,2^{-k})-P(\beta_n,0)]$, a decreasing infimum of computable reals. Rational upper approximations and running minima give upper semicomputability. An error-$1/4$ approximation separates $M_n=0$ from $M_n>119/125$, deciding halting. A proposed $g$ would instead give a computable secant below $1/2$ on halting inputs and above $119/125$ on nonhalting inputs. Approximation to error $1/10$ separates these cases, again deciding halting. The unbounded-scale principle is the one in \cite[Theorem~11.5]{DMthermal}; here the variable is a field secant for one fixed interaction.
\end{proof}

\begin{corollary}[Finite common-map certificates]\label{cor:certificate}
\proved\ Let $r_F\ge1$ bound the sup norm of every offset in $N_F$, and put $M_H=(2r_FH+1)^6$. Exponential coalescence of the specified maps is equivalent to $M_Hq_H(\eps)<1$ for some $H$. These strict tests form an effectively open set of rates, and select exactly the halting windows on the promised domain.
\end{corollary}
\begin{proof}
For the last $H$ maps to transmit initial dependence, their composite must be nonconstant and at least one of their at most $M_H$ time-$t$ inputs must depend on the initial row. Independence of the two map blocks and the union bound give
$q_{t+H}\le M_Hq_Hq_t$. Iteration and monotonicity of pullback images give exponential decay if $M_Hq_H<1$. Conversely an exponential beats the polynomial $M_H$ for large $H$. Finite-cone enumeration makes $q_H(\eps)$ an effectively computable polynomial in $\eps$, so strict positivity margins give effective openness. On halting windows use \cref{thm:exponential}; on nonhalting windows coalescence would contradict the two stationary laws. This is the $\Z^6$ form of \cite[Propositions~4.2--4.3]{DMst}.
\end{proof}

\begin{corollary}[No computable bound on forgetting times]\label{cor:times}
\proved\ At the rational midpoint $\eps_n$ of $I_n$, let
\[
 \Delta_n(t)=\max_{x,y}d_{\rm TV}((\delta_xP_{\eps_n}^t)_{\{0\}},
                                  (\delta_yP_{\eps_n}^t)_{\{0\}}).
\]
The first times $\Delta_n(t)<1/2$ and $q_t(\eps_n)<1/2$ are partial computable functions with domain exactly $\{n:U(n)\text{ halts}\}$. Neither has a total computable upper bound on its domain. The mean all-input CFTP depth is finite on each halting input and has no total computable bound over these inputs.
\end{corollary}
\begin{proof}
The finite cone makes $\Delta_n(t)$ and $q_t(\eps_n)$ exact computable rationals. Both are nonincreasing: an extra initial step mixes deterministic-start laws for $\Delta_n$, and nested pullback images give monotonicity for $q_t$. Coupling gives $\Delta_n(t)\le q_t(\eps_n)$. On halting inputs they tend to zero by \cref{thm:exponential}. On nonhalting inputs the stationary bit marginals differ by more than $119/125$. These marginals are mixtures of the deterministic-start marginals at every $t$, so their diameter gives
$\Delta_n(t)>119/125$ for every $t$.

A total bound for either first time would decide halting by the corresponding exact test at that bound. If $T_n$ is the actual origin pullback depth, nestedness gives $\Prob(T_n>t)=q_t(\eps_n)$, so its expectation is finite by \cref{thm:exponential} and at least half the first time $q_t<1/2$. A total bound on these expectations would therefore bound those first times. This specializes \cite[Corollary~8.7]{DMforget} to a fixed rule.
\end{proof}

\begin{corollary}[Finitary coding and exact finite-window sampling]\label{cor:coding}
\proved\ On halting windows the stationary space--time law is a finitary, translation-equivariant factor of a finite-alphabet i.i.d.\ field on $\Z^7$, with exponential tails for coding depth, and is mixing of every order. At rational rates it admits exact finite-window coupling from the past. At rational rates its finite marginals are computable by a partial algorithm that terminates on halting inputs. No operation-count bound for exhaustive enumeration is asserted.
\end{corollary}
\begin{proof}
At each update use one innovation taking the deterministic-map value or one of $N_\Sigma$ replacement values, with probabilities $1-\eps$ and $\eps/N_\Sigma$. The nested pullback images eventually agree at every target by \cref{thm:exponential}. Their common value defines a translation-equivariant factor, depending on a finite random past cone. Its depth tail is bounded by $C_ne^{-c_nt}$; for a finite set use the union bound. The spatial radius is at most $r_F$ times the depth.

At a rational rate, sample the innovations in nested finite cones, enumerate all finite initial words and stop when the target outputs agree. This almost surely terminates and gives their exact joint law, the Propp--Wilson construction \cite{PW96}; see the finitary-coding context of \cite{vdBS99}. A deterministic partial algorithm instead enumerates $q_t$ until the finite-window failure bound is below a desired tolerance and uses any fixed-start law at that depth. All these finite laws are computable.

For mixing of every order, truncate finitely many local coded observables at a fixed depth. At sufficiently large pairwise separation their finite innovation cones are disjoint, so the truncations are independent. The sum of coding-failure probabilities controls the difference from the true joint law. Let separation and then truncation depth tend to infinity. This proves the assertion for the full space--time translation action. These coding-depth bounds do not control the time needed to enumerate all initial words.
\end{proof}

\begin{corollary}[Complexity on and beyond the promise]\label{cor:complexity}
\proved\ If $U$ is effectively many-one universal, the rational ergodicity set restricted to $\bigcup_n I_n$ is c.e.-complete, and its relative complement in that promise is co-c.e.-complete. For any fixed $U$ used above, the full positive-rational ergodicity set is $\Pi^0_2$.
\end{corollary}
\begin{proof}
Membership of a positive rational in the union of windows is decidable: their rational endpoints decrease to zero, so only finitely many can contain it, and endpoints are compared exactly. Within the promise, find its unique $n$ and simulate $U(n)$; \cref{thm:main} gives a semidecision for ergodicity. The computable midpoint map reduces the halting set to this set. Under the stated universality hypothesis that halting set is c.e.-complete. Taking relative complements proves the other completeness assertion.

For the full set, the table is rational and strictly positive at each positive rational rate. By Halawi's uniform-ergodicity result, as reproduced in \cite[Proposition~10.7]{DMst}, and \cite[Lemma~10.8]{DMst}, ergodicity is equivalent to $\forall I\ \forall k\ \exists t\ \Delta(I,t)<1/k$, where $\Delta(I,t)=\max_{x,y}d_{\rm TV}((\delta_xP_\eps^t)_I,(\delta_yP_\eps^t)_I)$. Each finite-window diameter is a computable rational by finite-cone enumeration. This is the recalled $\Pi^0_2$ placement of \cite[Remark~10.9]{DMst}; compare \cite[Remark~8.9]{DMforget}. The extra universality assumption is confined to this corollary and is not an assumption of \cref{cor:undecidable}.
\end{proof}

\paragraph{Scale for the six-dimensional construction.} Every window lies below $10^{-5\cdot10^8}$, and the alphabet has more than $10^{10^{31539}}$ symbols. With the constants above, $k=52386$, $\log_{10}B=31539.51470\ldots$ and $n_0=888653587$. The rare-birth term in \eqref{eq:epsstar} is the minimum. It and the largest window give
\begin{align*}
 \log_{10}\eps_*&=-535038537.39331\ldots,\\
 \log_{10}\sup I_1&=-535038537.61745\ldots.
\end{align*} Indeed the helper alphabet alone has $100^B$ symbols, and $\log_{10}(2B)=31539.81573\ldots>31539$. 

\begin{remark}[Chernoff alternative for the scales]\label{rem:chernoff}
\proved\ The original choices $w_B=2^{-\lfloor BC^2/50\rfloor}$, $k=52386$ and $n_0=888653587$ are retained in this construction because a Lean formalization of Theorem~A (Theorem~9.3) and Corollary~9.4 uses them (\cref{app:lean}). A smaller alternative follows from a standard multiplicative Chernoff bound \cite{Chernoff52}. For $X\sim\operatorname{Bin}(B,p_W)$, put $\vartheta=.8C$. We have $s\le .8CB$ since $B\ge20/C$, and $p_W>.9C$, so
\[
 \Prob(X<s)\le
 \exp\!\left[-B\frac{(p_W-\vartheta)^2}{2p_W}\right]
 \le e^{-BC/180}.
\]
Indeed the derivative of $(p-\vartheta)^2/(2p)$ is
$(p-\vartheta)(p+\vartheta)/(2p^2)>0$ for $p>\vartheta$, and its value at $.9C$ is $C/180$. Thus one could replace only the definition of $w_B$ by $2^{-\lfloor BC/180\rfloor}$ and recompute all later constants in their original order. The recomputed values are
\[
 k=26200,\qquad n_0=458784211,\qquad
 \log_{10}B=15773.97177279\ldots,
\]
\[
 \log_{10}\eps_*=-276231384.60659605\ldots,\qquad
 \log_{10}\sup I_1=-276231384.83074221\ldots.
\]
The rare-birth term remains the minimum. Also
$\log_{10}(2B)=15774.27280278\ldots$. The integer searches and these logarithms were independently recomputed with a second arithmetic method in the verification record. These alternative values are not substituted into the fixed rule of this paper.
\end{remark}

\section{Lower dimension}\label{sec:lower}

All constants and alphabets in this section belong to a second fixed construction. They do not replace those of Sections~2--9 or Appendices~A--C. We write $C_4=C$ and $n_4=n_0$ when comparing with the introduction. The planar chart, finite machine and interval controller remain those of Sections~2--4, with the new fixed $B,n_0$ specified below. Only the transverse geometry and the stochastic proof change. Theorems~9.6--9.7, Corollaries~9.8--9.12 and Theorem~B.1 keep their original dimensions; no lower-dimensional version of their extremal, field-response, effective or robustness conclusions is asserted. The marked-parent mechanism is that of \cite[Section~12, Lemmas~12.3--12.4]{DMforgetLower}; its complete-history coherence argument must also be made quantitative for finite histories here.

\subsection{One fixed rule and terminating parameter choices}

Use $(u,P)\in\Z^2\times\Z^2$, the planar Moore set $\Mo$, and
\[
 \begin{gathered}
 m=16,\quad L=160,\quad K=256,\quad q=768,\quad h=512,\\
 S_0=\{0,\ldots,15\}^2,\quad S_1=160e_1+S_0,\quad
 S_2=160e_2+S_0,\quad S=S_0\cup S_1\cup S_2,\\
 T=500,\qquad v=(58,58),\qquad V=175,\qquad
 D=\operatorname{diam}_{\ell^1}(S)=350.
 \end{gathered}
\]
The live alphabet is the same $3\times3$ block presentation $\mathcal L_R$ with one global bit; $\mathcal A_R=\mathcal L_R\sqcup\{\botv\}$. Define $F_R$ by the same eight-neighbour planar check, retaining a live letter when at least $h$ of these $q$ checked providers retain it. Since $2h>q$, the result is unique. Keep the 100-letter helper alphabet, with LOW, UNKNOWN and 98 HIGH aliases, and put $\Sigma_4=\mathcal A_R\times\mathcal A_R\times\HH^B$.

For clarity, the entire controller prescription is as follows. At proposal time $t$ use
\[
 A_t=F_R(D_{t-1}),\qquad
 Q_t(u,P)=F_R^T(D_{t-1})(u,P-(T-1)v).
\]
Field $j$ at $w$ uses base $z=w-j\mathbf e_y$ and predecessor $w-\mathbf e_x+v$. Its provisional interval $J_{j,t}(w)$ is UNKNOWN if $A_t(z)=\botv$; it is $\{0\}\cup I_{j,t-1}(w-\mathbf e_x+v)$ if $A_t(z)$ is live but $Q_t(z)=\botv$; it is $\{0\}$ if both are live and $\Orig(Q_t(z))$; otherwise it is the predecessor interval. At base $w$ let $h_w$ count the $J_{j,t}(w+j\mathbf e_y)$ containing one, and $l_w$ those equal to $\{1\}$. Erase the data when $Q_t(w)=\botv$, or when $Q_t(w)$ is a phase-zero end and either its row flag is one with $h_w\ge s$, or its row flag is zero with $l_w<s$. Propose the erased or retained $A_t$ as data, $A_t$ as the second coordinate, and the fixed representative of each $J_{j,t}$ as its helper output. Then perform independent uniform replacement of the whole product symbol with probability $\eps$. This defines the rule on every configuration, including malformed symbols and exceptional walls. No driver reads an old helper, and the rule receives neither $n$ nor $\eps$.

Here is a complete rational specification of its constants. Put
\[
 \begin{gathered}
 \Delta=10^6(T+2)(9q+10),\quad X=2\Delta^2,\quad
 K_{\rm mem}=2^{20}X^{257},\\
 K_{\rm hist}=2^{40}[100(V+1)]^{10}X^{100},\qquad
 C=(128K_{\rm mem})^{-2}.
 \end{gathered}
\]
For $B=4^k$ define
\[
 \begin{gathered}
 u_B=2^{-\lfloor3BC/100\rfloor},\quad
 z_B=2^{-\lfloor BC/3\rfloor},\quad
 w_B=2^{-\lfloor BC^2/50\rfloor},\\
 \eta_B=\frac{3u_B}{1-4z_B}+2w_B.
 \end{gathered}
\]
Choose the least integer $k\ge1$ for which $B\ge20/C$, $4z_B<1$, and
\[
 \eta_B<\min\{(64K_{\rm mem})^{-6},(64K_{\rm hist})^{-4}\}.
\]
Set $s=\lceil3CB/4\rceil$ and
\[
 \begin{gathered}
 n_A=9q,\quad n_Q=(2T+1)^2[T(L+m-1)+1]^2,\quad
 \kappa=2B(n_A+n_Q+2),\\
 \lambda=176,\quad M=h-K=256,\quad
 N_0=\binom{(2\lambda)^2}{M},\\
 \Delta_{\rm rare}=10^6[\kappa+B+(2\lambda)^2+q+T+1],\quad
 K_{\rm rare}=2^{20}(2\Delta_{\rm rare}^2)^{513}.
 \end{gathered}
\]
The dependency count is proved exactly as in \cref{prop:local}: the $B$ driver bases, their two-dimensional virtual cones and the $B$ helper predecessors use at most $B(n_A+n_Q+2)$ parents; reverse degree is the same. Recentring translates the virtual support and does not enlarge its cardinality. Define
\[
 \begin{split}
 B_0&=2(18q)^2=382205952,\\
 \eps_*&=\min\{C/100,\ (64K_{\rm mem})^{-4},\ (64K_{\rm hist})^{-44},\\
 &\hspace{31mm}[64K_{\rm rare}(N_0+1)]^{-16},\ B_0^{-33720}\}.
 \end{split}
\]
Choose the least $n_0\ge k+2$ with
\[
 \frac{21C}{10(2\cdot4^{n_0+1}-1)}<\eps_*.
\]
Only now construct the finite machine and chart palette of Sections~3--4, using these $B,n_0$. All counts above are independent of the palette size. The $B$ search terminates because its three exponential factors tend to zero; the $n_0$ search terminates because its left side tends to zero. The finite-approximant phase decoder of Section~3 also terminates. Thus no circular parameter choice is involved. For this new recipe, independent exact-integer/rational and certified fixed-point interval implementations verify the least values $k=43406$ and $n_0=713161615$, and that the rare-birth term is the minimum defining $\eps_*$. These are the new scale values; the six-dimensional values preceding this section are unchanged. The exact parameter recipe, rather than an approximate decimal scale, specifies the rule.

For every $n\ge1$ put
\[
 \begin{gathered}
 a=n+n_0,\quad \ell_d=2\cdot4^d-1,\quad
 I_n^{(4)}=\left(\frac{2C}{\ell_a},\frac{21C}{10\ell_a}\right),\\
 W=\lfloor C/\eps\rfloor\quad(\eps\in I_n^{(4)}),\qquad
 W_n^*=\lfloor\ell_a/2\rfloor=4^a-1.
 \end{gathered}
\]
The intervals are nonempty rational open intervals, uniformly computable and pairwise disjoint: $\ell_{a+1}=4\ell_a+3$ gives $21/(10\ell_{a+1})<2/\ell_a$. They lie below $\eps_*$, and the same integer rounding as in \cref{prop:parameters} gives
\[
 \ell_{a-1}<W,\qquad 2W\le\ell_a,\qquad W\le W_n^*.
\]
In particular both critical depths are phase-zero rows. These are analysis windows for one already fixed alphabet, neighbourhood and rule.

\begin{theorem}[The four-dimensional noise diagram]\label{thm:lower}
For the single rule just specified, for every $n\ge1$ and every $\eps\in I_n^{(4)}$, ergodicity is equivalent to halting of $U(n)$. Nonhalting gives two stationary laws invariant under every translation of $\Z^4$. If $U(n)$ halts, an algorithm terminating on that input finds an integer $C_n\ge1$ such that for every integer $t\ge0$, every $\eps\in I_n^{(4)}$ and every finite $F\subset\Z^4$,
\[
 \Prob\{\text{the shared depth-$t$ full-symbol map is nonconstant on }F\}
 \le |F|C_n\exp\!\left[-\frac{t}{4\cdot701(W_n^*+52)}\right].
\]
The same bound holds for a forward block and a pullback block, over all initial configurations simultaneously. Consequently every finite-window pullback image is almost surely eventually a singleton, and the unique stationary law attracts every initial law. The decay rate is total computable from $n$; the prefactor algorithm is only asserted to terminate on halting inputs.
\end{theorem}

\subsection{Marked-parent density and finite-history coherence}

Take two-sided independent replacement flags and their least finitary closure $\MM$ under the rule of two marked ordinary parents in distinct clusters. Ordinary parents have offsets $\Mo\times S$ and lower physical time by one. Write $\mathsf H(z)$ when $z$ has some ordinary parent in $\MM$. This auxiliary notation is distinct from the physical second coordinate $H$. Our raw set is exactly the marked set of \cite[Section~12]{DMforgetLower} with logical dimension $r=2$ and logical parent set $\{-1,0,1\}^2$. Set $\mathcal D_0=\varnothing$ and, for $N\ge1$, let $\mathcal D_N(z)$ mean that some ordinary backward $N$-edge path from $z$ has $\mathsf H$ at least $N/1404$ times in its first $N$ vertices.

The hypotheses of \cite[Lemmas~12.3 and 12.4]{DMforgetLower} apply as stated: the logical parent set is $\Mo$ of size nine, the three provider boxes and quorum are exactly theirs, and $\eps\le B_0^{-33720}$. In their notation take $\delta=1/1404$ and $C_\delta=16860$, reserving $B$ here for the helper count. They give finite contours with at most $12(f-1)$ edges and distinct physical flag terminals, the one-cluster exclusion budget $K$, and
\[
 \Prob(\mathcal D_N(z))\le
 4\sum_{f\ge\lceil N/16848\rceil}(B_0^{16860}\eps)^f
 \le\frac{16}{3}e^{-N}.                                      \tag{L.1}\label{eq:lower-density}
\]
Indeed $B_0^{16860}\eps\le B_0^{-16860}<1/4$, and
$16860\log B_0/16848>1$. The event holds uniformly over all paths, including paths chosen after seeing all innovations or initial configurations. Almost surely every sufficiently long path from any root has fewer than $N/1404$ marked-parent vertices, and no infinite all-$\MM$ ordinary path exists. We use these published lemmas directly rather than reproving their physical contour and density estimates.

\begin{lemma}[Finite-history coherence]\label{lem:lower-coherence}
Fix arbitrary replacement values and a trajectory starting with an arbitrary full row at time zero. Two unmarked live data or second-coordinate readings at time $s\ge1$, at equal or Moore-neighbouring planar sites and transverse distance $d$, agree or are compatible, respectively, if $s\ge2d+1$ and neither root has $\mathcal D_{s-1}$. The assertion is simultaneous over every initial row. In complete histories the corresponding coherence and compatibility hold almost surely without an age restriction; any nonempty unmarked live domain then determines one full legal time-constant planar chart.
\end{lemma}
\begin{proof}
At transverse distance zero, live data and second-coordinate outputs at positive time equal their ordinary proposal off flags. At neighbouring planar sites the two retained-provider sets overlap in at least $2h-q=256$ offsets; a check at a common offset verifies compatibility, including diagonal neighbours. No time-zero hypothesis is used.

Suppose a pair conflicts at positive distance. An unmarked live root has $h$ successful providers and at most $K$ contaminated offsets, so at least $h-K=256$ clean successes. Each supplies an unmarked live data centre with the same letter and a clean checked Moore star. If neither current root has $\mathsf H$, every provider is clean. Pair offsets within each box by a signed unit shift reducing the transverse distance. There are $3\cdot15\cdot16=720$ injective pairs. The two roots exclude at most $2(q-h)=512$ failed pairs, so a pair of successful clean centres remains, preserving the conflicting labels and reducing distance by one. At a step with $\mathsf H$ on either side, choose any clean success at each root; distance increases by at most $D=350$. Planar sites and their displacement remain fixed throughout.

Continue for $N=s-1$ steps, stopping if distance becomes zero. It cannot become zero while the remaining age is at least one, by the base case. Otherwise, with $b_N$ counting steps having $\mathsf H$ on either path,
\[
 1\le d_N\le d-N+(D+1)b_N.
\]
Outside the two $\mathcal D_N$ events, $b_N<2N/1404$, so $d_N<d-N/2\le0$ for $N\ge2d$, a contradiction. Every exceptional event depends only on flags, establishing the all-start assertion. In a complete history continue the paired paths arbitrarily far and use the almost-sure all-path density conclusion instead.

Let $V_t$ be the planar sites admitting unmarked live readings in a complete history. The just-proved coherence labels them uniquely; clean stars give $V_t+\Mo\subset V_{t-1}$ and preserve the centre label. A nonempty domain expands backwards to cover the plane. Nested domains and persistent labels define a full chart, with each Moore compatibility checked at a common earlier time. Every later unmarked live reading has its label. This is the block-chart version of \cite[Lemmas~12.6--12.8]{DMforgetLower}; unlike those raw-tile statements, the finite-history assertion above includes our retained second coordinate and all eight Moore directions.
\end{proof}

\begin{lemma}[Characteristic contours]\label{lem:lower-poles}
For the new geometry the grouped charge-arm estimate is
\[
 e\le A_0(g-1)+g\Omega/129,\qquad A_0=1440/129<12.
\]
Here $g$ counts distinct terminal groups, $\Omega$ bounds their total three-potential oscillation, and unary zero-potential links are excluded from $e$. With time-decreasing support edges and time-preserving identifications the span is at most $e$. A wrong live depth-$500$ forecast whose wrong actual leaves are marked has at least eleven distinct flags, at most $50(f-1)$ edges, and, at characteristic lag $\ell\le W$ from time $b$, earliest flag time at least $b-\ell-1-50(f-1)$.
\end{lemma}
\begin{proof}
For physical time use balanced potentials
\[
 \psi_0=3(P_1+P_2)+348t,\qquad
 \psi_1=-3P_1-174t,\qquad\psi_2=-3P_2-174t.
\]
A nonzero-cluster support gives charge zero gain at least $480-348=132$; a support outside cluster $i$ gives charge $i$ gain at least $174-45=129$. Every support gain is at least $-351$. In characteristic coordinates $\chi=P+vt$ the same potentials are $3(\chi_1+\chi_2),-3\chi_1,-3\chi_2$ and support displacements are $s-v$. Thus the characteristic gains are $129,-351$, rather than the $130,-350$ of the separate physical density proof in \cite[Lemma~12.3]{DMforgetLower}.

The frontier proof of \cref{lem:pole} applies with three charges: frontier components are joined by overlapping descendant terminal-group names; a binary exposure splits a component into at most two and cannot join an outside component. At a split add one three-charge source, placing each new endpoint in the daughter lacking its old endpoint, then follow good supports. Only the new source's first edge can be bad. At the terminal frontier, each charge bijects the sources and distinct selected groups in the root incidence component. There are $g$ sources and at most $3(g-1)$ bad edges, also when branches merge. Summing increments and comparing each group with one reference leaf gives
\[
 129e-480\cdot3(g-1)\le g\Omega,
\]
which proves the bound. Connectedness proves the clock claim when group and unary links preserve time; for unequal-time equal-potential groups only the edge bound is used.

For a wrong live forecast, every retained centre has its wrong letter, hence supplies wrong supports in two clusters at every virtual level. Append finite marked derivations to the wrong actual leaves. A virtual $F_R^j(D_s)$ uses characteristic $P+vs+jv$ and formal clock $s+j$; actual data use $P+vs$ and physical clock $s$. Recentring puts the root at the output characteristic. Every arm crosses $T=500$ virtual levels, so $500\le e<12(f-1)$; in particular $f\ge11$, and the conservative $50(f-1)$ bound holds. All actual flags are no later than the queried input time. Each support lowers the formal clock by one, giving the asserted conservative earliest-time bound. Auxiliary negative-time flags are genuine paid-for flags, never substitutes for arbitrary initial letters.
\end{proof}

\subsection{Planar computation, sensors and nonhalting memory}

The planar specification in Sections~3--4 is unchanged as a finite algorithm, now instantiated with the new $B,n_0$. In particular \cref{prop:scaffold,lem:cutoff,lem:allphase,lem:machine,lem:input,prop:chart} use no transverse coordinate. The finite-rank wall hypotheses of \cite[Lemmas~A.6--A.7]{DMhard} apply inside each closed frame even when the other phase coordinate is exceptional: the frame residues lie in $[4^d+1,3\cdot4^d]$ modulo $4^{d+1}$ and exclude $0,1$, forcing finite parent rank at most $2d$. The laminar routing argument therefore still forces every $N_d\times N_d$ board, $N_d=2^{d+2}-2$, with precisely its initialized one-head computation.

More explicitly, a cutoff at rank $2a-2$ or $2a-1$ makes the first odd-rank zero on a depth-$d>a$ board have index $2^{a+1}-2$ and rank $2a-1$. The machine marks that many cells plus two, halves until one remains, obtains $a+1$ tokens, skips $n_0+1$, and runs $U(a-n_0)=U(n)$. Smaller boards fail, have no zero, or meet an even-rank zero first. Thus nonhalting supplies a legal typical-phase cutoff chart with either bit. Conversely, if every depth-$(a-1)$ flag were one and every depth-$a$ flag zero, the same rank constraints would force $U(n)$ on every sufficiently large board in every phase. A board exceeding both the halting time and visited tape space would contain the forbidden halting state. This proves the same all-phase wrong-gate alternative for the new palette.

\begin{lemma}[Sensor law and source counts]\label{lem:lower-sensors}
For every fixed reference-chart phase and every $\eps\in I_n^{(4)}$, the intrinsic gate experiments have independent lanes, bin failure probability at most $\eta_B$, and critical-gate failure probability at most $\eta_B$. For $g$ distinct characteristic/bin labels, the joint probability of a raw or statistical source at every label is at most
\[
 \gamma^g,\qquad \gamma=\sqrt C+\eta_B^{1/6}.
\]
Memory witnesses with $g$ distinct labels have at most $K_{\rm mem}^g$ encodings. Historical witnesses with $f$ distinct flags and $m_0$ intrinsic restarts have at most $K_{\rm hist}^{f+m_0}W^{6f/11}$ encodings. These counts are independent of $B,W,n,n_0$ and palette size except for the explicitly displayed factors.
\end{lemma}
\begin{proof}
The invariant coordinates of backward copying are $(x-t,y,P+vt)$. A phase-zero row owns exactly its $B$ physical lane rows; those row blocks partition the plane. A finite ruler owns the source interval $[L_d,R_d-1]$ along each characteristic. Different rulers' intervals are disjoint, so each physical atom has one owner, and the $B$ lanes use disjoint atoms. The integer drift $58$ changes no ownership identity. The uniform product replacement gives upper success probability $.99[1-(1-\eps)^{\ell_d}]$ and recent lower success probability $.98[1-(1-\eps)^W]$.

For a short ruler $d=a-1-r$, $\eps(\ell_d+1)<.53C4^{-r}$. The exponential-moment bound for Bernoulli sums at threshold $.75CB$ has exponent exceeding $.03BC$ at $r=0$ and increasing by at least $BC/3$ per rank. At most $3\cdot4^r$ such ends lie in a $W$-bin, so their summed tail is at most $3u_B/(1-4z_B)$, exactly as in \cref{lem:tails}. For long rulers, $\eps W\ge.99C$ and $C<.01$ give success probability greater than $.9C$, whereas $s/B\le.8C$. The independent-lane lower tail is at most $e^{-BC^2/50}$; at most two long ends lie in a bin. This gives $\eta_B$. The bounds use only ruler lengths, independent lanes and the window margins, so the numerical proof of \cref{lem:tails} applies without a dimensional change.

A raw bin has probability at most $\eps W\le C$. A statistical bin uses its own and preceding bins on its private lanes. Partition statistical labels by time-bin residue modulo three; a largest class is independent and bounds $g_S$ simultaneous failures by $\eta_B^{g_S/3}$. Distinct raw bins are independent. Cauchy--Schwarz bounds mixed tests by $C^{g_R/2}\eta_B^{g_S/6}$, even if raw and statistical supports overlap. Summing the two choices at each label proves the stated joint source law. All these arguments first fix a phase and are uniform in it.

Here are the literal counts, to specify their dependence. The projected graph uses actual data/second-coordinate roles, $T$ virtual levels, source-hub and entry/corridor roles. Including neutral and charge decorations there are at most $32(T+2)$ roles. Support offsets have at most $18q$ choices; ten zero-offset types, five bin differences and both orientations give degree at most
\[
 640(T+2)(9q+10)<\Delta.
\]
A connected $e$-edge witness and its terminal subset have at most $2(2\Delta^2)^e$ doubled-traversal encodings. In a finite memory derivation all leaves with one source label have identical characteristic potentials. \Cref{lem:lower-poles} gives fewer than $12(g-1)$ positive edges. Each source-free characteristic corridor and its entry cost at most two links per positive edge, plus three terminal-role links per group, for fewer than $39g<256g$ edges. This yields
\[
 2\sum_{e=0}^{256g}X^e\le4X^{256g}\le K_{\rm mem}^g.
\]
The arm estimate precedes projection; projected cycles or merged branches do not change it. Characteristic links have zero spatial displacement and bin difference at most two, so no $B$ or $W$ choices are hidden in the graph degree.

For a historical contour with $f_i$ flags, use the conservative $50(f_i-1)$ duration. Its lag has at most $2W$ choices, its restart gap at most $52f_iW$, and its net transverse endpoint at most $(V\,\mathrm{gap}+1)^2$ choices. Keeping the larger four-coordinate envelope bounds their product by
\[
 2[54(V+1)]^5 f_i^5W^6.
\]
Traversal, terminal subset and root role cost at most $X^{51f_i}$. Since $f_i^5\le2^{5f_i}$, size compositions cost $2^f$, interleavings $2^{f+m_0}$, and endpoint/source/stop annotations at most $2^{4(f+m_0)}$, their product, including $q\le X$ choices per intrinsic restart, is bounded by $K_{\rm hist}^{f+m_0}W^{6c}$. Each contour has eleven flags, so $c\le f/11$. This proves the historical count, including a last whole contour extending below the scanned prefix.
\end{proof}

\begin{proof}[Nonhalting direction of \cref{thm:lower}]
Start data and second coordinate in a correct cutoff chart of bit $b$, cloned transversely. Initialize helpers by ideal negative-time chart instructions and independent negative-time innovations; almost surely an earlier reset or replacement determines each helper. Average the joint law spatially, so the phase is Haar and typical almost surely. This initialization is covariant under every planar and transverse translation, and independent of positive-time noise.

A bad ordinary forecast has at least $q-h+1=257$ providers not returning the reference letter, each with a bad Moore input, occupying two clusters. The same holds at every level of a bad deep forecast. At an unflagged data error, either $A$ is bad, $Q$ is bad, or a correct critical gate falsely erased. On a short correct gate its true origin is within $W$ steps, and correct drivers reproduce the upper experiment. On a long correct gate the last HIGH marks in the recent $W$ steps supply the necessary singleton fields independently of older helpers. Thus a false erasure requires its intrinsic bad event or a driver error along one base characteristic at lag at most $W$. A bottom driver counts as an error. Every source-free characteristic link lands at a driver with a positive support step. This is a finite acyclic derivation: recursion lowers physical time or a finite virtual level, and the pre-zero drivers are the perfect chart of the initialization.

Apply the memory count and source law of \cref{lem:lower-sensors} before integrating over the chart. Uniformly in site, time and rate, its error probability is at most
\[
 \sum_{g\ge1}(K_{\rm mem}\gamma)^g<3/125,
\]
using $K_{\rm mem}\sqrt C=1/128$ and $K_{\rm mem}\eta_B^{1/6}<1/64$. The live data bit-$b$ cylinder therefore has probability greater than $122/125$ at all times. The kernel is Feller on a compact configuration space. Temporal Ces\`aro subsequential limits are stationary by telescoping the one-step difference, and preserve invariance under all four spatial translations. The two disjoint bit cylinders separate the two limiting laws. This proves nonergodicity with the claimed invariance scope.
\end{proof}

\subsection{Halting: a finite motif and a simultaneous noise event}

\begin{lemma}[Finite motif certificate]\label{lem:lower-motif}
On every halting input a terminating finite search finds $r\ge1$ such that every Moore-compatible live patch on $Q_r=[-r,r]^2$ contains a type S or type L motif. Type S is a segment of $\ell_{a-1}+1$ letters with $\Orig$ only at its left endpoint, and $\End$, phase zero and row flag zero at its right endpoint. Type L has $\ell_a+1$ letters and right flag one.
\end{lemma}
\begin{proof}
The planar wrong-gate alternative just proved holds in every phase. Its finite ride has only its endpoint walls: a bottom row of depth $d$ lies outside all smaller row frames, same-depth column frames are disjoint, and strict nesting excludes deeper endpoints from the closed frame. The finite-rank classification applies there even in exceptional phases. Thus every legal full chart contains one of these motifs. Enumerate all live patterns on increasing $Q_r$, test internal Moore compatibility and reject both contained motifs. If a survivor existed for every $r$, product compactness would give a legal full chart with neither motif, a contradiction. Every finite test uses only the effective palette and its finite phase lookup. The first empty test terminates on halting inputs, with no extension oracle. This is precisely the planar search of \cref{lem:motif}, instantiated with the new fixed parameters.
\end{proof}

Fix such an $r$, put $R=r+T$, $A=2D+1=701$ and $J=W_n^*+52$. Choose a computable integer radius $r_F\ge1$ bounding every actual read offset. The following stronger-than-needed burn-in will avoid all floor and boundary cases:
\[
 t\ge t_0:=4[A(R+W_n^*+1)+2D(T+1)+4].              \tag{L.2}\label{eq:lower-burnin}
\]
From an unmarked live root $(u_0,P_0,t)$, propagate $R$ clean checked stars backwards. They supply unmarked live data at every site of $u_0+Q_R$ at time $t-R$, with providers in $P_0+[0,VR]^2$. At each such site choose a clean centre path further backwards. The supplied letter persists and its provider at age $\sigma$ lies in $P_0+[0,V(t-\sigma)]^2$.

Follow a motif endpoint's path for $N=\lfloor t/(2A)\rfloor$ vertices starting at $t-R$. At sampled time $b$ and lag $\ell\le W$, a depth-$T$ query has actual input age $\sigma=b-\ell-1$ and leaf providers
\[
 P_b+\ell v-(T-1)v+\sum_{i=1}^T s_i,\qquad s_i\in S.
\]
Relative to $P_0$ they lie coordinatewise in $[-VT,V(t-\sigma+T)]$; their distance from the reference providers is at most $D(t-\sigma+T)$. Every planar query leaf lies inside $Q_R$. Since $t-\sigma\le R+N+W_n^*$, \eqref{eq:lower-burnin} implies
\[
 \sigma\ge2D(t-\sigma+T)+1,\qquad \sigma\ge\lceil t/2\rceil+1,
\]
and the analogous age condition for the reference patch at $t-R$.

To make these comparisons simultaneous, take
\[
 R_n=10[r_F+B+V(T+W_n^*+r+1)+r+T+W_n^*+1]
\]
and let $\mathcal G_t$ exclude $\mathcal D_{s-1}(z)$ at every deterministic root with
\[
 \lceil t/2\rceil+1\le s\le t,\qquad
 \|z-(u_0,P_0)\|_\infty\le R_n(t+1).
\]
All patch, path, recentered query and reference roots used above lie in this set; if necessary its spatial radius also encloses their one-step ordinary parents. The displayed $R_n$ already dominates those fixed offsets. By \eqref{eq:lower-density}, for $t\ge t_0$,
\[
 \Prob(\mathcal G_t^c)\le
 16(t+1)[2R_n(t+1)+1]^4e^{-t/2}.                  \tag{L.3}\label{eq:lower-good}
\]
The harmless factor $16$ covers integer rounding. On $\mathcal G_t$, \cref{lem:lower-coherence} makes the reference patch compatible and identifies every unmarked live query leaf with its reference letter, simultaneously for every initial row. For Moore compatibility one can equivalently take a clean checked star and apply same-site agreement one physical age earlier. All those earlier roots lie in the same enlarged deterministic set; the burn-in also gives $t-R-1\ge2D(R+1)+1$. In particular, \eqref{eq:lower-good} also gives the weaker estimate, since $t\ge t_0\ge2$:
\[
 \Prob(\mathcal G_t^c)\le12(t+1)[2R_n(t+1)+1]^4e^{-t/4}.
\]
This event uses flags only. All remaining certificate estimates are unconditional; none is obtained by conditioning on $\mathcal G_t$ or on a selected patch.

\begin{lemma}[Uniform unmarked-liveness estimate]\label{lem:lower-unmarked}
An integer $C_U$ is effectively findable from $r$ and the fixed parameters such that for every $\eps\in I_n^{(4)}$, every integer $t\ge0$ and every fixed target, the probability that some initial row yields an unmarked live data or second-coordinate reading there at time $t$ is at most $C_Ue^{-t/(701J)}$.
\end{lemma}
\begin{proof}
For $t\ge t_0$ work on $\mathcal G_t$. A putative live root gives a compatible reference patch, a motif and its sampled centre path. A wrong live deep query has marked actual leaves, hence the contour of \cref{lem:lower-poles}, with at least eleven flags. At a sampled live unmarked data output, $A$ and current $Q$ are live and the erasure test did not fire. A wrong current $Q$ is a lag-zero contour. Otherwise it reads the correct wrong-gate motif. If no wrong live deep instruction occurs on its short ride or long recent window, the helper inclusions force erasure unless the intrinsic experiment failed. For a short flag-zero gate the origin starts with a set containing the ideal upper bit; UNKNOWN, union with zero, correct copies and replacement intervals preserve that inclusion. Fewer than $s$ upper ones forces $l_w<s$. For a long flag-one gate, $s$ last-HIGH marks preserve the possibility of one under these same operations, so $h_w\ge s$. Both contradict the surviving data. This argument holds for arbitrary initial helpers. Thus every sampled vertex supplies an intrinsic failure or a wrong live deep query at one lag at most $W$.

Attach whole deep contours greedily and restart one time before their earliest flag, retaining the last contour even when it passes the prefix end. Intrinsic failures advance one centre step. Distinct contour terminal slabs are disjoint; with $f$ flags, $c$ contours and $m_0$ intrinsic tests, $c\le f/11$. Their coverage obeys
\[
 N\le[50+(W+1)/11]f+m_0\le J(f+m_0).
\]
At one fixed planar endpoint, distinct sampled times have disjoint intrinsic atoms: an atom's planar coordinate is $g_x-\ell$ and time $b-\ell$, so equality of atoms forces the same lag and sampled time, regardless of the chosen provider path. For a fixed certificate the intrinsic tests therefore have probability at most $\eta_B^{m_0}$; raw tests have probability $\eps^f$. Their two families may overlap. H\"older with exponents $4/3$ and $4$ gives weight at most $\eps^{3f/4}\eta_B^{m_0/4}$.

Multiply by the historical count in \cref{lem:lower-sensors} and use $W\le C/\eps$. The resulting bases are
\[
 K_{\rm hist}C^{6/11}\eps^{9/44}<1/64,\qquad
 K_{\rm hist}\eta_B^{1/4}<1/64.
\]
The first follows from $\eps\le(64K_{\rm hist})^{-44}$ and the second from the $B$ search. In particular both are below $1/32$. Summing masses $k=f+m_0\ge\lceil N/J\rceil$ gives
\[
 \sum_{k\ge\lceil N/J\rceil}(k+1)32^{-k}
 \le(16/15)16^{-\lceil N/J\rceil}.
\]
There are at most $2|Q_r|(VR+1)^2$ endpoint/type candidates. The motif segment fixes the tested instructions and atoms; there is no count over arbitrary chart fillings. Each fixed certificate is a noise event, and the deterministic union covers existence of every possible start. Add \eqref{eq:lower-good} by an unconditional union bound. Since $N\ge t/(2A)-1$ and $\log16/2>1$, the historical tail has rate at least $1/(AJ)$. The polynomial in \eqref{eq:lower-good} decays faster. Increasing an effective integer constant absorbs this polynomial and all $t<t_0$, proving the assertion uniformly over the interval.
\end{proof}

\subsection{Rare births and full-symbol all-input coalescence}

Define the minimal scalar centre-only field at each planar site by finite-past increasing limits,
\[
 W^{\rm live}_t=E_t\vee\Theta_h(W^{\rm live}_{t-1}),\qquad
 N^{\rm live}_t=\Theta_h(W^{\rm live}_t).
\]
Here $\Theta_h$ tests the $q$ transverse centres and uses no planar Moore neighbours. A birth of $N^{\rm live}$ with at most $M-1=255$ fresh flags has at least $h-(M-1)=257=K+1$ old $N^{\rm live}$ providers, hence two distinct clusters. Its alternative seed requires $256$ fresh flags.

Partition transverse space into squares $Y=\lambda z+[0,\lambda-1]^2$. All fresh flags for a seed rooted anywhere in $Y$ lie in $\lambda z+[0,2\lambda-2]^2$. The group event that this box has at least $M$ flags has probability at most $N_0\eps^M$. Different planar sites and times have disjoint atoms; equal-parity squares have disjoint enlarged boxes. Among $g$ distinct groups a largest one of the four parity classes has at least $g/4$ members. Since the seed weight is below one, the weaker bound
\[
 \Prob\{\text{all }g\text{ group events}\}\le(N_0\eps^{256})^{g/16}
                                                               \tag{L.4}\label{eq:lower-seeds}
\]
also holds. The groups seed whole root squares in the minimal two-cluster closure.

Within one group, planar site and physical time agree and the total three-potential oscillation is at most $\Omega=2\cdot3\cdot2\cdot175=2100$. Hence \cref{lem:lower-poles} gives
\[
 e\le H_0g-A_0,\qquad H_0=(1440+2100)/129=3540/129<28,
 \qquad A_0>2.                                                \tag{L.5}\label{eq:lower-groups}
\]
These identifications preserve physical time.

\begin{lemma}[Late ancestry and discrepancy packing]\label{lem:lower-full}
The one-site full-symbol all-input coefficient is at most $C_ne^{-t/(4\cdot701J)}$, for an integer $C_n$ effectively findable on halting inputs and uniform over $I_n^{(4)}$.
\end{lemma}
\begin{proof}
First justify the minimal-field domination for finite histories. A prescribed marked ordinary path of $L_*$ vertices has probability at most
\[
 \sum_{f\ge\lceil L_*/50\rceil}(K_{\rm mem}\eps)^f.
\]
To see this, attach whole flag contours and restart before the earliest terminal; selected slabs have distinct flags, each contour covers at most $50f_i$ vertices, and the whole path plus contours costs at most $101f$ edges. The doubled-traversal bound used in \cref{lem:lower-sensors} applies. Its base is below $1/64$.

Put $m_t=\lfloor t/4\rfloor$. In a deterministic enlarged dependency cone of the time-$t$ target, union-bound \cref{lem:lower-unmarked} and this marked-path estimate with $L_*=m_t+1$ at all relevant roots of ages $\lceil t/2\rceil$ through $t$. The cone includes every fixed helper shift and ordinary ancestor used below and has size bounded by a computable polynomial in $t$, for example after enlarging $r_F$, $(t+1)[2r_F(t+1)+1]^4$. Except on a noise event of probability at most a computable polynomial times $e^{-t/(2\cdot701J)}$, simultaneously over all starts, every relevant live reading is marked and no marked branch of $m_t$ ordinary edges persists.

A live ordinary proposal along the top $m_t$ discrepancy vertices has its live data parents at time at least $t-m_t$. Follow their retained live centres, stopping at actual flags. Through $m_t$ generations they remain in the late cone, since $t-2m_t\ge\lceil t/2\rceil$. All nodes are marked and a depth-$m_t$ branch is excluded. Finite branching therefore makes this ancestry finite with flag leaves. Induction puts it in the minimal $W^{\rm live}$ field, and the proposal in $N^{\rm live}$ at its true input time. This step does not assign noise contours to arbitrary time-zero symbols.

If the full-symbol map at the target is nonconstant, choose a discrepant pair after fixing the innovations. A replacement output is common. A differing deterministic output has a discrepant actual parent. Moreover at least one of the $B$ ordinary proposals at bases $w-j\mathbf e_y$ is live in one run: if all were bottom, both data coordinates would be bottom and every helper would be the common UNKNOWN, irrespective of old helpers or deep instructions. Follow discrepant parents for $m_t$ vertices. On the good late-cone event, attach to its chosen live-proposal base a grouped contour at its actual time $s-1$, keeping every group time unchanged. The connector from discrepancy time $s$ to time $s-1$ lies outside the pole contour. Restart before its earliest group time, retaining the whole last contour. By \eqref{eq:lower-groups}, coverage including connector and restart is at most $e+2\le H_0g_i$. Distinct selected slabs give distinct groups, and $m_t\le H_0g$.

Support and covered-path edges cost at most $2H_0g$, group stars at most $3g$ and literal planar shifted-root connectors at most $g$; allowing the conservative extra connectors gives fewer than $(2H_0+6)g<62g<512g$ edges. A supergraph with 64 roles includes actual path offsets, ordinary supports, square stars and all $B$ literal shifts. Its degree is at most
\[
 64[2\kappa+2B+2(2\lambda)^2+18q+2(T+1)]<\Delta_{\rm rare}.
\]
Thus doubled traversals and source subsets give at most $K_{\rm rare}^g$ certificates. By \eqref{eq:lower-seeds} the weighted base is
\[
 K_{\rm rare}(N_0\eps^{256})^{1/16}<1/64,
\]
using the defining threshold and $\eps^{256}\le\eps$. Summing $g\ge\lceil m_t/H_0\rceil$ gives exponential rate at least $\log32/(4H_0)$. This and the marked-path rate $\log64/200$ exceed $1/(2\cdot701J)$. Add the late-cone exception and absorb its polynomial by halving that slowest exponent. All constants are effective from the motif radius, finite lookup radius and rational parameters. For example $y^ke^{-by}\le(k/b)^k$ for $k\ge1,b>0$, together with rational lower bounds on logarithms, gives an effective integer absorption constant. Enlarge it for finitely many small times. Every certificate is an unconditional noise event covering existence of any discrepant pair; no selected pair or history is presumed independent of the noise. This proves the claimed coefficient bound.
\end{proof}

\begin{proof}[Completion of the proof of \cref{thm:lower}]
Nonhalting was proved above. On halting inputs \cref{lem:lower-motif,lem:lower-full} give the stated rate and a prefactor search terminating on every such input. Translation and a union bound give the finite-set factor $|F|$. The ordered forward and pullback blocks have the same distribution. At a fixed finite window the pullback images of all starts are nested nonempty finite sets, and their nonconstancy probabilities tend to zero. Hence they are eventually singletons almost surely, simultaneously over the countably many finite windows. Compactness and the Feller property provide a stationary law. Coupling it to any initial law bounds their finite-window total variation distance by the all-input coefficient, proving attraction and uniqueness. The midpoint map $n\mapsto41C/(20\ell_{n+n_0})$ reduces halting to ergodicity at rational rates for this one fixed rule. Thus its noise diagram is uncomputable and Theorem~C follows.
\end{proof}

\subsection{The all-Gibbs transfer and a fixed interaction on five dimensions}

The exact transfer needed here is \cite[Corollary~13.7, equations~(7)--(8)]{DMst}, whose statement covers every spatial dimension $d\ge1$. We record its hypotheses and proof to distinguish all-DLR uniqueness from stationary uniqueness. For probability laws on a finite target set, $d_{\rm TV}$ denotes total variation distance. A basic local-map representation means independent identically distributed random maps at the update sites whose evaluation at each parent word has the specified transition law; the same map acts on every possible parent word.

\begin{proposition}[Exponential common maps imply all-DLR mixing]\label{prop:lower-transfer}
Let a translation-invariant PCA on $\Z^d$, $d\ge1$, have a finite alphabet with at least two symbols, finite neighbourhood $N$ containing zero and strictly positive transition table $p$. Suppose it has an i.i.d.\ basic local-map representation with
\[
 q_t=\Prob\{\text{the depth-$t$ common map at one site is nonconstant
 over all initial rows}\}\le C_0e^{-\alpha t}
\]
for every integer $t\ge0$, with $C_0,\alpha>0$. The space--time interaction $-\log p$ has exactly one DLR measure $\bar\mu$, without an invariance assumption. There are constants $A_*,c_*>0$ such that for every finite $\Lambda\subset\Z^{d+1}$, every $E\subset\Lambda$ and all exterior configurations $\tau,\tau'$,
\[
 \begin{aligned}
 d_{\rm TV}((\gamma_\Lambda^\tau)_E,\bar\mu_E)
 &\le A_*\sum_{z\in E}e^{-c_*\operatorname{dist}_\infty(z,\Lambda^c)},\\
 d_{\rm TV}((\gamma_\Lambda^\tau)_E,(\gamma_\Lambda^{\tau'})_E)
 &\le 2A_*\sum_{z\in E}e^{-c_*\operatorname{dist}_\infty(z,\Lambda^c)}.
 \end{aligned}
\]
Every factor meeting $\Lambda$, including each exterior-child factor, is included in $\gamma_\Lambda^\tau$. Holes and disconnected volumes are allowed. No lower bound on $\alpha$ in terms of the radius is required.
\end{proposition}
\begin{proof}
Put $r=\max(1,\max_{z\in N}\|z\|_\infty)$. Choose an integer slab height $H$ and partition spatial rows into cells of side $b=rH$. A slab marker $(j,k)$ is bad when its $H$-step output word on cell $k$ is nonconstant over all input rows. Its probability is at most
\[
 p_*=b^dC_0e^{-\alpha H}.
\]
It depends only on the slab's backward seed cone $U_{j,k}$. Within one slab cones at cell-index distance greater than two are disjoint, and different slabs are independent. One of the $3^d$ cell residue classes therefore supplies independent markers for any fixed list.

For a single root spin use the expanding future-only volume
\[
 V_n=\{(0,0)\}\cup
 \{(x,t):1\le t\le nH,\ \|x\|_\infty\le5b+3rt\}.
\]
Starting with the root cell $S_0$, let $S_j$ be the bad cells at index distance at most three from $S_{j-1}$. Every reached cell is within $3j$ cells of the root, and its seed cone is within $(3j+2)b$. This is inside the buffer, whose half-width at the beginning of slab $j$ is $5b+3r(j-1)H=(3j+2)b$. Thus the clipped and unrestricted common maps agree wherever this exploration uses them.

We need a deterministic splice before any probability conditioning. Fix original seeds $x$ and a reached prefix through level $j$, and replace seeds only in the union of its listed cones, obtaining seeds $y$. At each level $i\le j$ the clipped outputs for a fixed root value differ only within two cells of $S_i(x)$. Indeed, a cell whose computation reads a replaced seed is within two cells of the corresponding listed bad cell. Otherwise its map is unchanged; a different output can only come from a differing input within two cells of $S_{i-1}$, and finite speed adds one cell. It is then a distance-three candidate and its unchanged original map is nonconstant, so its original marker is bad and reached. This proves the induction without any assumption on the repaired marker status. The number of changed spins at level $i$ is at most $c_0|S_i|$, where $c_0=(5b)^d$. For identical seeds and different root values, the same induction confines differences to $S_i$ itself. Extinction therefore makes the top row common for every root value.

All exterior children of $V_n$ lie in its top cap, because the sides expand faster than the dependency radius. Let $h$ be the product of their transition factors. With root value $a$, clipped forward evolution $X^a$ and fixed parents, write
\[
 Z_a=\mathbb E h(X^a_{nH}),\qquad w_a=p(a\mid\text{fixed parents}).
\]
The finite-volume root law is $w_aZ_a/\sum_cw_cZ_c$. The cap cannot be discarded. Put
\[
 K_p=\max_{c,u,u'}\frac{p(c\mid u)}{p(c\mid u')},\qquad
 \zeta=|N|\log K_p.
\]
If two rows $z,z'$ at slab level $j$ differ in $k$ spins, their integrated future messages satisfy $g_j(z)\le e^{\zeta k}g_j(z')$. Pair identical future histories term by term: only immediate next-row factors change, at most $|N|k$ of them, each by a ratio at most $K_p$. Later factors and the cap agree. At the final level the same argument applies directly to cap factors. The ratio bound is independent of the remaining future duration and exterior values.

Fix a deterministic exact reached-prefix event $E_S$ through level $j$. Its mass is $M_j=\sum_{i=1}^j|S_i|$, excluding the root $S_0$, and its last width is $a_j=|S_j|$. Independently resample the deterministic union $U(S)$ of its listed seed cones. On $E_S$ the splice and message bounds give
\[
 g_j(X^a_j(x))\le e^{\zeta c_0a_j}g_j(X^a_j(y)).
\]
Only after using this inequality enlarge $E_S$ to the event $B_S$ that every listed nonroot marker is bad in $x$. This event is measurable in $x$ restricted to $U(S)$ and is independent of the repaired product seed field $y$. The residue-class bound gives $\Prob(B_S)\le p_*^{M_j/3^d}$. Consequently, with $b_0=\zeta c_0$,
\[
 \mathbb E[\ind_{E_S}h(X^a_{nH})]
 \le e^{b_0a_j}p_*^{M_j/3^d}Z_a.                 \tag{L.6}\label{eq:lower-prefix}
\]
This is an unconditional finite-product resampling inequality; it does not assert posterior independence under a Gibbs law.

If $S_n$ survives, each level is nonempty and $|S_i|\le(6i+1)^d$. Choose $L_0=\max(2,b_0)$. For all sufficiently large $n$, some $j\in[\lceil n/2\rceil,n]$ has $M_j\ge L_0a_j$. Otherwise $M_j>[L_0/(L_0-1)]M_{j-1}$ throughout that half, contradicting the polynomial bound $M_n\le\sum_{i\le n}(6i+1)^d$. The finite threshold depends only on $d,L_0$.

A prefix is a connected rooted vertex set in a level-cell graph of degree at most $\Delta_d=2\cdot7^d$. Its spanning-tree doubled traversal gives at most $\Delta_d^{2M}$ prefixes of mass $M$, jointly over all heights. At the selected late cut $e^{b_0a_j}\le e^{M_j}$. Choose $H$ so large that
\[
 \theta_*=e\Delta_d^2p_*^{1/3^d}<1.
\]
This is possible for every $\alpha>0$ since $H^de^{-\alpha H}\to0$. The possibly large $b_0$ changes only the finite late-cut threshold, not this inequality. Summing \eqref{eq:lower-prefix} over deterministic prefixes gives, for large $n$,
\[
 \mathbb E[\ind_{\{S_n\ne\varnothing\}}h(X^a_{nH})]
 \le\delta_nZ_a,\qquad
 \delta_n=\frac{\theta_*^{\lceil n/2\rceil}}{1-\theta_*}.
\]
On extinction the top cap row is common for all $a$, so $Z_a=Z^{\rm common}+R_a$ with $R_a\le\delta_nZ_a$. The root conditional is a mixture of $w$ and a residual law of mass at most $\delta_n$. Enlarging a constant for bounded $n$ gives root error at most $C_1s^n$, uniformly over every exterior, where $s=\sqrt{\theta_*}<1$.

For any DLR measure $\nu$, use the translated future-only volumes at a site $v$. The entire past and every other same-row spin are outside each such volume. The tower property and the uniform root error tending to zero imply
\[
 \nu(\sigma_v=a\mid\text{entire past and all other same-row spins})
 =p(a\mid\text{parents})
\]
almost surely. Factoring finite subsets of a row gives the product PCA transition given the past. Thus every DLR measure is a two-sided PCA trajectory law, with no invariance assumed. A stationary PCA law $\mu$ exists by Feller compactness; its trajectory law is DLR by the finite-factor cancellation, including exterior children, of \cite[Lemma~2.4]{DMst}. If $\nu_t$ is any DLR row marginal, then $\nu_t=\nu_{t-m}P^m$. Common-map coupling bounds its distance from $\mu$ on finite $F$ by $|F|q_m$. Letting $m\to\infty$ identifies every row marginal, and the product transitions identify the whole trajectory law. This proves all-DLR uniqueness.

Finally take any finite $\Lambda$ and exterior and sample its Gibbs law in increasing time, with a fixed order within each row. Whenever $w+V_n$ fits in $\Lambda$, previously exposed spins are outside that future-only volume; the tower property bounds the chain-rule conditional $K_w$ from the PCA update at its actual parents by $C_1s^n$. Draw a fresh basic map $M_w$ independently of the joint past and evaluate it at those parents, obtaining $Y_w$. Couple its transition law maximally to $K_w$, then sample the actual spin from that coupling conditionally on $Y_w$. The Gibbs chain rule is exact, and the basic maps retain their i.i.d.\ law. Draw one stationary reference base row independently below all target cones, evolve it with these maps and with fresh maps outside $\Lambda$, and use this common reference for all targets.

For a target $z$ at distance $\rho$ from $\Lambda^c$, take backward depth $k=\lfloor(\rho-1)/(2r)\rfloor$ and $n$ maximal with $5b+3rnH\le(\rho-1)/2$. For sufficiently large $\rho$ these are positive and all required comparison volumes fit. A union bound on the cone and the unconditional all-input map event give
\[
 \Prob(\text{actual}_z\ne\text{reference}_z)
 \le k(2rk+1)^d C_1s^n+C_0e^{-\alpha k}.
\]
No conditioning on absence of mismatches is used; the actual lower row may depend on the maps, but common-map constancy covers every input. Both depths grow linearly in $\rho$, so polynomial absorption bounds this by $A_*e^{-c_*\rho}$. For bounded $\rho$ enlarge $A_*$ using the probability bound one. The common coupling and a union bound over $E$ give the first estimate; the triangle inequality gives the second. Only distance to the complement was used, so disconnected volumes and holes require no additional hypothesis.
\end{proof}

\begin{corollary}[The five-dimensional binary-energy diagram]\label{cor:lower-gibbs}
There is one computably specified translation-invariant finite-range interaction on $\Z^5$, on a finite alphabet and with individual terms in $\{0,1\}$, and uniformly computable pairwise disjoint nonempty rational open inverse-temperature intervals $\mathcal J_n^{(5)}$, such that for every $n\ge1$ and $\beta\in\mathcal J_n^{(5)}$ uniqueness among all Gibbs measures is equivalent to halting of $U(n)$. On halting windows the specification has arbitrary-volume exponential weak spatial mixing with the two bounds in \cref{prop:lower-transfer}; their constants may depend on $\beta$, but not on the volume, target set or exteriors. On nonhalting windows there are at least two Gibbs measures invariant under every translation of $\Z^5$. Reciprocating endpoints also gives rational open temperature windows.
\end{corollary}
\begin{proof}
Let $f_4$ be the deterministic rule of \cref{thm:lower}, $N_F$ its finite neighbourhood and $N_\Sigma=|\Sigma_4|$. At $(w,t)\in\Z^4\times\Z$ use the factor
\[
 \Phi_{w,t}(\sigma)=
 \ind\{\sigma(w,t)\ne f_4((\sigma(w+z,t-1))_{z\in N_F})\}
\]
on the support consisting of its child and all parents. Every factor has value zero or one; the interaction is fixed, finite-range and translation-invariant. In every finite-volume specification include all factors whose supports meet the volume. The parent-independent normalizer gives
\[
 \begin{aligned}
 p_\beta(c\mid\eta)
 &=\frac{e^{-\beta\ind\{c\ne f_4(\eta)\}}}
 {1+(N_\Sigma-1)e^{-\beta}}\\
 &=(1-\eps(\beta))\ind\{c=f_4(\eta)\}+\eps(\beta)/N_\Sigma,
 \qquad \eps(\beta)=\frac{N_\Sigma}{e^\beta+N_\Sigma-1}.
 \end{aligned}
\]
Dividing every factor by this constant cancels from finite-volume weights, including the exterior-child factors. Thus this is exactly the replacement PCA space--time specification.

The neighbourhood contains zero, because the ordinary driver includes the zero Moore and provider offsets (adding a dummy parent would also be harmless). The alphabet has at least two symbols (already $\HH$ has 100). Every transition entry is at least $\eps(\beta)/N_\Sigma>0$. The independent whole-symbol replacement maps are basic i.i.d.\ local maps. On a halting window \cref{thm:lower} supplies the required exponential all-input coefficient with positive rate $1/[4\cdot701(W_n^*+52)]$. All hypotheses of \cref{prop:lower-transfer}, equivalently \cite[Corollary~13.7]{DMst}, therefore hold in $d=4$. Ordinary stationary uniqueness alone would not supply this conclusion.

On a nonhalting window the two stationary laws invariant under all four spatial translations have distinct time-zero marginals. Their stationary two-sided trajectory laws are DLR by \cite[Lemma~2.4]{DMst} and invariant under all five space--time translations; those time-zero marginals distinguish them. Normalization preserves the binary-energy specification.

If $I_n^{(4)}=(\eps_n^-,\eps_n^+)$, its inverse-temperature preimage is
\[
 \left(\log(N_\Sigma/\eps_n^+-N_\Sigma+1),\
       \log(N_\Sigma/\eps_n^--N_\Sigma+1)\right).
\]
The endpoints are computable reals and strictly ordered. Enumerate rational approximations until two ordered rationals are certified strictly inside, and take the open interval between them as $\mathcal J_n^{(5)}$. Openness guarantees termination uniformly in $n$, without a halting answer. Strict monotonicity of $\eps(\beta)$ inherits disjointness from the rate windows. The endpoints are positive; their reciprocals give the stated temperature windows. This proves Corollary~D and answers \cite[Open problem~9.2]{DMhard} in dimension five with finite energies.
\end{proof}


\begingroup
\edef\threeoldequation{\number\value{equation}}
\setcounter{equation}{0}
\renewcommand{\theequation}{\thesection.\arabic{equation}}
\renewcommand{\theHequation}{three.\arabic{equation}}
\section{Dimension three}\label{sec:three}

This section specifies a third rule; none of its alphabets or constants replaces an earlier one. Section~10 and Appendices~A--C retain their statements and numbers. The physical space is $(i,P)\in\Z\times\Z^2$, with chart coordinates
\begin{equation}\label{eq:three-chart-time}
 (x,y)=(i,t-i).
\end{equation}
WEST and SOUTH are the actual-time parents $(i-1,t-1)$ and $(i,t-1)$. We call this determinism from southwest inputs; \cite{LGO12} calls the resulting completion NE-deterministic. Only two-input uniqueness and a lift of one canonical configuration are needed from its cover theorem.

\subsection{Fixed parameters and the statement}

Fix
\[
 \begin{gathered}
 S_0=\{0,\ldots,15\}^2,\quad S_1=(160,0)+S_0,\quad
 S_2=(0,160)+S_0,\quad S=S_0\cup S_1\cup S_2,\\
 q=768,\quad h=512,\quad K=256,\quad T=500,\quad
 V=175,\quad v=(58,58),\quad C_*=2^{40}.
 \end{gathered}
\]
Here $K$ is the size of one provider box. Put
\begin{equation}\label{eq:three-memory-constants}
 \begin{gathered}
 \Delta_{\rm mem}=10^6(T+2)(9q+10),\quad X=2\Delta_{\rm mem}^2,\\
 K_{\rm mem}=2^{20}X^{257},\qquad C=(128K_{\rm mem})^{-2},\\
 \Delta_{\rm full}(B)=10^{30}(T+2)^{10}(q+1)^{10}(B+1)^4,\\
 K_{\rm full}(B)=2^{20}[2\Delta_{\rm full}(B)^2]^{2C_*}.
 \end{gathered}
\end{equation}
For $B=4^k$ define
\[
 u_B=2^{-\lfloor3BC/100\rfloor},\quad
 z_B=2^{-\lfloor BC/3\rfloor},\quad
 w_B=2^{-\lfloor BC^2/50\rfloor},\qquad
 \eta_B=\frac{3u_B}{1-4z_B}+2w_B.
\]
Choose the least $k\ge1$ with
\begin{equation}\label{eq:three-B-search}
 B\ge20/C,\quad4z_B<1,\quad
 \eta_B<\min\{(64K_{\rm mem})^{-6},(64K_{\rm full}(B))^{-4}\}.
\end{equation}
Set $s=\lceil3CB/4\rceil$ and, as a definite rational number, set
\begin{equation}\label{eq:three-epsstar}
 \eps_* =\min\{C/100,\ 382205952^{-33720},\
 (64K_{\rm mem})^{-4},\ (64K_{\rm full}(B))^{-44}\}.
\end{equation}
Choose the least $n_0\ge k+2$ with $21C/[10(2\cdot4^{n_0+1}-1)]<\eps_*$. Only then construct the chart and the machine of Section~4 with this fixed $n_0$. For $n\ge1$ put
\begin{equation}\label{eq:three-windows}
 \begin{gathered}
 a=n+n_0,\quad \ell_d=2\cdot4^d-1,\quad
 I_n^{(3)}=(2C/\ell_a,\ 21C/(10\ell_a)),\\
 W=\lfloor C/\eps\rfloor,\qquad W_n^*=4^a-1.
 \end{gathered}
\end{equation}
Throughout $I_n^{(3)}$, $\ell_{a-1}<W$, $2W\le\ell_a$ and $W\le W_n^*$. The windows are rational, nonempty and pairwise disjoint, since $\ell_{a+1}=4\ell_a+3$ and $21/(10\ell_{a+1})<2/\ell_a$. They lie below $\eps_*$. The rule reads neither $n$, $\eps$ nor $W$.

The $B$ search terminates: $C$ is fixed, $\eta_B$ decays exponentially in $B$, and $K_{\rm full}(B)$ is a fixed polynomial. All graph counts below are independent of palette size. Two independently written exact-arithmetic interval programs in \texttt{verification/z3-scale/} certify
\begin{equation}\label{eq:three-scale}
 \begin{gathered}
 k=43429,\qquad n_0=33644133613821753975,\\
 \log_{10}\eps_*\simeq-20255786791774346221.28,
 \end{gathered}
\end{equation}
with $(64K_{\rm full}(B))^{-44}$ binding. These are not the scale values of Section~10. Define
\begin{equation}\label{eq:three-rate}
 \begin{gathered}
 r_F=B+2(V+1)(T+1)+10=B+176362,\\
 K_0=10^8(V+1)(r_F+1)=17600000000(4^{43429}+176363),\\
 c=\frac1{2^{42}K_0}
   =\frac1{2^{42}\cdot17600000000(4^{43429}+176363)}.
 \end{gathered}
\end{equation}
Thus $\log_{10}(1/c)=26169.7521\ldots$. Exact recipes specify the rule.

\begin{theorem}[One noise diagram on $\Z^3$]\label{thm:three}
The single computably specified finite-alphabet finite-range automaton below has, for every $n\ge1$ and every $\eps\in I_n^{(3)}$, ergodicity under uniform whole-symbol replacement if and only if $U(n)$ halts. Nonhalting gives two stationary laws invariant under every spatial translation. On a halting input an algorithm terminates and returns an integer $C_n\ge1$ such that for every $\eps\in I_n^{(3)}$, integer $t\ge0$ and finite $F\subset\Z^3$,
\[
 \Prob\{\text{the shared depth-$t$ full-symbol map is nonconstant on }F\}
 \le |F|C_ne^{-ct}.
\]
This holds simultaneously over all starts, for forward and pullback blocks. Every finite-window pullback image is almost surely eventually a singleton; the unique stationary law attracts every initial law. The coefficient $c$ in \eqref{eq:three-rate} is independent of $n$, and the effective halting-dependent prefactor is uniform on the whole window.
\end{theorem}

\subsection{An automatic geometry with a deterministic cover}

\begin{lemma}\label{lem:three-cover}
There is an effective finite geometry alphabet, deterministic from its full west and south neighbours, projecting to the scaffold, orientation/event and lane-phase constraints of Section~3. It has a full-plane tiling, and its projected finite-board geometry holds in every represented phase.
\end{lemma}
\begin{proof}
Use the decorated 276-letter substitution $\sigma$ of \cite[Tables~4--5]{DMhard}. Its child $(1,1)$ has components $(b,F(\theta),\mathsf A,0,0)$; the latter's child $(0,0)$ has components $(a,F(\theta),\mathsf A,1,1)$. Thus the child $(2,2)$ map of $\sigma^2$ fixes
\[
 t_*=(a,\angle b,\mathsf A,1,1),
\]
since $F(a)=\angle b$. Applying that map twice is constant at $t_*$. Table~5 admits this letter (a thin-$a$ corner requires equal ruler bits).

Set $q(z)=\lfloor(z-1)/4\rfloor$ and $d(z)=(z-1)\bmod4$ componentwise. Every integer coordinate iterates to $-1$. Start with $t_*$ at $(-1,-1)$ and apply $\sigma^2$ with offset $(1,1)$. The nested legal patches exhaust the plane; their limit satisfies $A_0(z)=\sigma^2(A_0(q(z)))(d(z))$. Its phase solves $\varpi=1+4\varpi$, hence $\varpi=-1/3$ on both axes. Quartering with offset zero gives $\tau_0=(-2/3,-2/3)$.

This phase is not typical in \cite[Definition~3.15]{DMhard}. The three properties used are: all ranks are finite; every coordinate has a finite frame owner; and walls on each row and column are unbounded both ways. The first follows from nonintegrality, as in \cite[Lemma~3.14]{DMhard}. For the second,
\[
 0-\tau_{0,i}\pmod{4^{d+1}}=(4^{d+1}+2)/3
       \in[4^d+1,3\cdot4^d]\quad(d\ge1).
\]
Any fixed integer perturbation remains in this interval for sufficiently large $d$. If a row belongs to a depth-$e$ row frame, the periodically repeated depth-$e$ column endpoints give walls arbitrarily far to either side; transpose for columns.

Check (O1)--(E3) directly. On edge $(x,x+1)$ put the side of the nearest vertical wall at or left of $x$ in the minus orientation entry, and that of the nearest wall at or right of $x+1$ in the plus entry. At a non-wall the west/east pairs coincide; at a wall of side $a$ the west plus and east minus entries both equal $a$. Transpose vertically. Put a pair event tag exactly on edges with orientation $(\mathrm L,\mathrm R)$, or $(\mathrm B,\mathrm T)$ vertically, and $*$ otherwise. On a carrier edge its event entries are the cell/wall types of the nearest stops on either side, with stops the active cells and bounding walls. Between stops entries copy. An active site resets west plus and east minus to cell; inward left/right walls reset their inward entries to wall. Transpose for south/north resets. These verify the orientation, tag-admissibility, copy and reset cases, including empty corridors with pair (wall,wall). Finite laminar frames give the carriers and active sites by \cite[Lemma~3.13, Lemma~A.8 and Proposition~5.1(b,c)]{DMhard}. Its typical-only Proposition~5.7(b) is not used.

This canonical coloring is effectively automatic on signed coordinates: base-four subtraction has finite borrow states; compositions of the sixteen finite child maps form a finite set of maps. Sign extension eventually selects child $(2,2)$, whose two-step constancy makes the output stable. Reversal and determinization give a conventional digit-output automaton. Quartering and fixed collars preserve automaticity. Nearest-wall and nearest-stop predicates are first-order formulas in automatic predicates and integer order. For instance a nearest left wall satisfies
\[
 \operatorname{Wall}(w,y),\quad w\le x,\quad
 \nexists w'\,[\operatorname{Wall}(w',y),\ w<w'\le x].
\]
Projection, complement and finite products implement the formulas. Include $p=y-\tau_{0,2}-1\pmod B$ and the finite phase-decoder collar from Section~3 in the output alphabet $\mathcal C$.

Each finite canonical patch has finite rank, owner, wall and stop witnesses determined by a finite hierarchy prefix. Sufficiently large positive translations divisible by an appropriate power of four preserve that prefix and the residue modulo $B$. Thus every patch recurs arbitrarily far in the positive quadrant. Convert the output automaton to base two, respecting leading zeros. For its residual-state set $Q$ and transition $\delta$, the substitution $\rho(q)(e)=\delta(q,e)$ has a fixed state quadrant projecting to this coloring. Translated state quadrants reproducing increasingly large canonical patches have a whole-plane compact limit $S_\infty$. For each $m$, taking the finitely many block phases and preimages to a subsequence gives a whole-plane $\rho^m$ preimage up to shift.

The definition in \cite[Section~1]{LGO12} is
\[
 \Lambda_\rho=\bigcap_{m\ge0}\ \bigcup_{u\in\Z^2}
 u\cdot\rho^m(Q^{\Z^2}),
\]
equivalently the set admitting an infinite desubstitution history. Hence $S_\infty\in\Lambda_\rho$. Theorem~4 of that paper effectively gives a strongly deterministic tileset projecting onto this limit set. Apply it once to the state substitution, compose with the output map, and reject violations of the original local scaffold/routing/phase and collar-overlap constraints. Append projected edge colors and overlap strips to covering colors. Incoming full neighbours still determine the completion, because these are restrictions of the old covering choices. The canonical coloring has a lift and survives. Every surviving projection satisfies the local geometry used in Lemma~3.3's all-phase finite-board argument. This is one cover including routing, not a product of two covers.
\end{proof}

Add $H_{\rm row}$ from the west and $V$ from the south, imposing \eqref{eq:rankfilter} and quarter equalities, and a bit equal on all edges. These only reject completions. Lemma~3.2 gives the common nonincreasing sequence $b_r$. Its cutoff converse uses only finite ranks and the black-type rank comparisons, so also applies to the canonical seed.

\subsection{The causal L-frame computation}

Retain the finite machine $P$ of Section~4, with the fixed $n_0$ above, but replace its tableau. Its content alphabet is enlarged by computational WALL, distinct from erasure. Write tape address $j\ge0$ as $k=-j$, put WALL at $k>0$, and use the reflected radius-one rule $F(a,b,c)=f_P(c,b,a)$ on non-WALL centers, interpreting a neighbouring WALL as the tape end, preserving central WALL, and sending attempted exits to $\fail$. The malformed-triple convention remains total. Forbid only central $\halt$.

An active payload has role LEFT, DIAG or RIGHT, a symbol $b$, a center cache $c$, and a pair $L$ for LEFT, $R$ for RIGHT, both for DIAG. Unused boundary pairs are fixed arbitrarily. The following is the complete interior table; unlisted role pairs are rejected.
\[
\begin{array}{c|c|c}
 (w\text{ role},s\text{ role})&\text{role}&\text{output}\\\hline
 (\{\mathrm{DIAG},\mathrm{RIGHT}\},\mathrm{RIGHT})&\mathrm{RIGHT}&
 \begin{gathered}w.R=(u,v'),\ b=F(u,v',s.b),\\R=(v',s.b),\ c=w.c\end{gathered}\\[5pt]
 (\mathrm{LEFT},\{\mathrm{DIAG},\mathrm{LEFT}\})&\mathrm{LEFT}&
 \begin{gathered}s.L=(u,v'),\ b=F(w.b,u,v'),\\L=(w.b,u),\ c=s.c\end{gathered}\\[5pt]
 (\mathrm{LEFT},\mathrm{RIGHT})&\mathrm{DIAG}&
 \begin{gathered}w.c=s.c=c_{\rm old},\ b=c=F(w.b,c_{\rm old},s.b),\\
 R=(c_{\rm old},s.b),\ L=(w.b,c_{\rm old})\end{gathered}
\end{array}
\]
On the first active column above the corner put LEFT with $b=\init(H_{\rm row},\chi_2,0)$ and $c$ copied from the south predecessor, of role DIAG or LEFT. On the first active row right of the corner put RIGHT with $b=\mathrm{WALL}$ and $c$ copied from the west predecessor, of role DIAG or RIGHT. At the corner put DIAG with $b=c=\init(H_{\rm row},\chi_2,1)$. Incoming outer colors are empty. The flags $f_W,f_S$ recognize these cases; $f_E,f_N$ suppress outgoing colors, without suppressing the payload or its no-halt test. Horizontal-only carriers copy the whole west payload east, vertical-only carriers copy south north, and other colors are empty. Interior input tracks are retained by $F$ and are not compared with external row flags: those comparisons occur only on the initial left leg.

\begin{lemma}[The L-frame chart]\label{lem:three-chart}
The resulting finite chart alphabet $\mathcal R$ is deterministic from west and south. Each depth-$d$ board represents the initialized $P$ computation on $t+|k|<N_d$, where $N_d=2^{d+2}-2$. For $a=n+n_0$, nonhalting supplies legal charts with either bit and cutoff $b_r=\ind\{r<2a-2\}$. Halting forces in every full chart either a depth-$(a-1)$ finite bottom end with flag zero or a depth-$a$ finite bottom end with flag one, in every represented phase.
\end{lemma}
\begin{proof}
The two incoming letters determine geometry, flags, bit, and then each carrier or active case. Thus the completion is unique. The geometric proof of Lemma~3.3 forces the same $N_d\times N_d$ grid and flags. Sparse channels are causal for time $x+y$: a gap of length $g$ takes $g$ local steps.

At active position $(i,j)$ set $t=\min(i,j)$, $k=i-j$. Induction on $i+j$ gives $b(i,j)=a^t_k$ for the $F$ orbit from the initial legs. If $i>j$, the south value is $a^{j-1}_{k+1}$ and west pair is $(a^{j-1}_{k-1},a^{j-1}_k)$; RIGHT gives the update and shifted pair. LEFT is the transpose. On the diagonal west/south give $a^{t-1}_{-1},a^{t-1}_{1}$ and both centers equal $a^{t-1}_0$, yielding DIAG and its two caches. Boundary pairs are never read at interior updates: the first interior pair comes from a computed diagonal. Since $\max(i,j)=t+|k|$, the represented region is exactly the one stated. Any infinite extension of its finite initial legs gives the same represented states; every finite halt appears on sufficiently large boards.

Lemmas~3.2 and~4.2 give the unchanged input calculation. At cutoff $2a-2$, a depth-$d>a$ board has its first odd-parity zero at $j_*=2^{a+1}-2$, and $P$ decodes $n$. A smaller board either has an even first zero and fails, or has an all-one extension on which SCAN runs forever. On large boards use the initialized cutoff tape. If $U(n)$ never halts, none of these finite restrictions contains halt. Fill each grid by its restriction and each channel from its endpoints; empty corridors have empty colors. Finite ownership of the canonical geometry covers all active sites, and the direct routing check gives legality. Choose either bit. Spatial averages have weak limits supported on legal charts. Their phase laws are Haar on finite residue tori, hence typical almost surely; each finite-rank cutoff condition persists.

If neither wrong critical gate exists, the common-rank filter forces $b_{2a-3}=1$ and $b_{2a-1}=0$. Lemma~4.2 then forces input $n$ on all sufficiently large boards. A finite halting trace appears in one L-frame diamond and violates the no-halt filter. The finite-board argument does not assume typicality. Both critical depths are gated since $a-1\ge k$. All algorithms terminate before selecting $n$; no halting answer enters the alphabet.
\end{proof}

\subsection{History strips and the literal noisy rule}

A live augmented symbol is a vertically matching $T$-tuple $c=(c_0,\ldots,c_{T-1})$ of chart tiles. With erasure $\botv$, define the total bottom-spreading generator $g$ by
\[
 c_0=g_{\mathcal R}(w_0,s_0),\qquad c_j=s_{j-1}\quad(1\le j<T).
\]
Return bottom if either input is bottom, the central completion fails, internal vertical matching fails, or any $c_j$ fails its horizontal match to $w_j$. On a legal chart $\tau$, $c_j(i,t)=\tau(i,t-i-j)$ is an all-live recurrence. Conversely every all-live recurrence projects to locally legal chart tiles. Let $\mathcal A$ contain all live tuples and bottom.

For $D\in\mathcal A^{\Z^3}$ let $F(D)(i,P)$ retain the unique live value occurring at least $h$ times among
\[
 g(D(i-1,P+\alpha),D(i,P+\alpha)),\qquad\alpha\in S,
\]
or bottom if none. Use the helper alphabet of Section~2 and $\Sigma_3=\mathcal A\times\mathcal A\times\HH^B$. At time $t$ form
\[
 A_t=F(D_{t-1}),\qquad Q_t(i,P)=F^T(D_{t-1})(i,P-(T-1)v).
\]
Decode a live $Q_t$ through slot $T-1$. On the moving reference, its $T$ forecast steps and $T-1$ old slots cancel: this is exactly the current tile $\tau(i,t-i)$.

Field $j$ at $(w,P)$ has base $(w+j,P)$ and predecessor $(w-1,P+v)$ at time $t-1$. Its provisional interval is UNKNOWN if base $A_t$ is bottom; predecessor interval union $\{0\}$ if $A_t$ is live but $Q_t$ bottom; $\{0\}$ if both are live and the decoded instruction is $\Orig$; otherwise the predecessor interval. At base $(i,P)$ gather these from $(i-j,P)$, $0\le j<B$. Let $h_i$ count intervals containing one and $l_i$ those equal to $\{1\}$. Pin when $Q_t$ is bottom, or its instruction is a phase-zero $\End$ and either $H_{\rm row}=1,h_i\ge s$, or $H_{\rm row}=0,l_i<s$. Output pinned-or-retained $A_t$, unpinned $A_t$ as second coordinate, and canonical representatives of provisional intervals. Then replace the entire symbol uniformly with probability $\eps$. This specifies malformed inputs too. History is inside the voted alphabet and is erased with data.

Every read has radius at most $r_F$. A conservative full-graph dependency count, including both orientations, is
\begin{equation}\label{eq:three-degree}
 \kappa(B)=2B(9q+n_Q+2),\qquad n_Q=(2T+1)^2(175T+1)^2.
\end{equation}
This includes literal helper shifts and virtual supports; the two-input support only reduces the old nine-input overcount. It is independent of palette size.

\begin{lemma}[Relocated sensors and deep witnesses]\label{lem:three-sensors}
Lemmas~6.1--6.3 and~10.4 give the same intrinsic laws, $\eta_B$, and mixed-source bound $\gamma^g$ with $\gamma=\sqrt C+\eta_B^{1/6}$. A wrong live decoded depth-$T$ instruction whose wrong actual leaves are marked has a complete witness with $f\ge11$ distinct physical flags and at most $50(f-1)$ support edges. A helper lag at most $W$ is one zero-characteristic link.
\end{lemma}
\begin{proof}
A physical atom $(w,P,t)$ belongs to a phase-zero experiment for a unique $j$ modulo $B$, because its base row is $y=t-w-j$. Lane $j$ lies at old $(x-j,y+j)$ relative to the base. Copying preserves $(i-t,P+vt)$. A phase-zero row owns $B$ consecutive physical rows; its finite rulers own disjoint intervals $[L_d,R_d-1]$, with exactly $\ell_d$ updates ending before the gate replacement. Thus each atom has at most one owner and a gate's lanes are independent. The scalar upper and recent lower probabilities remain $.99[1-(1-\eps)^{\ell_d}]$ and $.98[1-(1-\eps)^W]$. Lemma~6.2's moment and lower-tail estimates give $\eta_B$. Raw bins cost at most $C$; every third statistical bin has disjoint support. Cauchy--Schwarz gives $\gamma^g$ without raw/statistical independence, exactly as in Lemma~6.3.

A wrong slot $j>0$ forces $h$ providers with wrong copied south slot $j-1$. At slot zero functionality gives a wrong actual input at west or south. No history copy bypasses a quorum level, and $h>K$ gives two clusters at every level. Attach complete raw-mark derivations at actual leaves. Lemma~10.3 applies with characteristic $\chi=P+vt$, poles $3(\chi_1+\chi_2),-3\chi_1,-3\chi_2$, gains $129,-351$, and coefficient $1440/129$. Logical offsets do not enter these poles. Crossing all $500$ virtual levels gives $f\ge11$ and at most $50(f-1)$ edges. Actual flags precede the queried input time. Recentring and copying preserve $\chi$, so helper lag adds only its one anchor.
\end{proof}

\begin{proof}[Nonhalting direction of Theorem~\ref{thm:three}]
Use either-bit cutoff charts from Lemma~\ref{lem:three-chart}, lift their strips and clone transversely. Initialize helpers by ideal negative-time instructions and innovations, spatially average, and keep positive-time innovations independent. A bad ordinary proposal has at least $257$ failing providers, each with a bad west or south input by functionality. Bottom $Q$ also counts as bad. The same recursion applies virtually. False erasure with correct drivers requires an intrinsic source or a driver error on one base characteristic within lag $W$, by the interval inclusions of Section~7.

The memory count of Lemma~10.4 transfers literally: all $B$ lanes of a base use the same $A/Q$ driver characteristic, so collapsing source-free corridors adds no $B$-valued role or shifted-driver choice. Logical degree $2q\le9q$ is within its bound. At most $32(T+2)$ roles, support choices, ten zero-offset types, five bin differences and both orientations give degree below $\Delta_{\rm mem}$. The characteristic arm estimate applied before projection gives fewer than $12(g-1)$ positive edges for $g$ source labels, and fewer than $39g<256g$ edges after corridors and entries. Doubled traversal and terminal subsets give
\[
 2\sum_{e=0}^{256g}X^e\le4X^{256g}\le K_{\rm mem}^g.
\]
Recursion lowers physical time or virtual level and ends at ideal pre-zero drivers. The source law bounds either data-coordinate error by
\[
 \sum_{g\ge1}(K_{\rm mem}\gamma)^g<3/125,
\]
since $K_{\rm mem}\sqrt C=1/128$ and $K_{\rm mem}\eta_B^{1/6}<1/64$. Its live bit-$b$ cylinder has probability above $122/125$ at every time. Feller compactness and temporal Ces\`aro limits give two stationary laws separated by these disjoint cylinders. Spatial invariance persists, proving the stated nonhalting conclusion.
\end{proof}

\subsection{Coherence down to time zero}

Let $M$ be the two-sided least finitary raw marked closure for ordinary parents $\{-1,0\}\times S$: a root is marked if flagged or if it has marked parents in two distinct provider boxes. Write $\mathsf H(z)$ if it has a marked ordinary parent, and $\mathcal D_N(z)$ if some $N$-edge backward path has $\mathsf H$ at least $N/1404$ times among its first $N$ vertices. Apply \cite[Lemmas~12.3--12.4]{DMforgetLower} with $r=1$, $J=\{-1,0\}$ (reflection of $\{0,1\}$), $|J|\le9$, and their constants
\[
 B_0=382205952,\quad\delta=1/1404,\quad C_\delta=16860,
 \quad e_0=B_0^{-33720}.
\]
For $\eps\le e_0$,
\begin{equation}\label{eq:three-density}
 \Prob(\mathcal D_N(z))\le(16/3)e^{-N}\le6e^{-N/4212}.
\end{equation}
These noise-only events control paths selected after innovations are fixed.

\begin{lemma}[Full-age pairing]\label{lem:three-coherence}
In one finite history of integer age $s\ge1$, two unmarked live readings at equal logical position/time and transverse $\ell^1$ distance $d$ agree if $s\ge2d$ and neither root has $\mathcal D_s$. This includes the ordinary-proposal coordinate and holds simultaneously over all starts. One further clean step gives local west/south chart compatibility.
\end{lemma}
\begin{proof}
At positive time, unmarked live data and stored ordinary-proposal readings at the same physical site both equal its ordinary proposal. After one backward step, both selected readings are of the first data coordinate. At distance zero these agree in one history, including at time zero; no equality between the two arbitrary time-zero coordinates is assumed. An unmarked live root has at least $h-K=256$ clean successes. Conflicting successful outputs have conflicting inputs at a common west/south index by functionality. If neither root has $\mathsf H$, signed unit pairing inside the three boxes gives $720$ pairs reducing distance by one; at most $2(q-h)=512$ fail. Otherwise choose any clean success at each root, and use functionality to select a common conflicting logical input index; distance increases by at most $350$. The two selected offsets need not coincide at these exceptional steps. This is the finite-history pairing of \cite[Theorem~12.12]{DMforgetLower}, now for the functional two-input tuple rule.

Run it for exactly the age $s$, down to time zero. If $b_s$ steps have $\mathsf H$ on either path,
\[
 d_s\le d-s+351b_s<d-s+351(2s/1404)=d-s/2\le0.
\]
Distance-zero conflict in the same history is impossible even at time zero. At a common clean provider $g$ checks west/south matching; same-site agreement one step earlier identifies those inputs with their references. Off flags a live second coordinate equals its proposal; a flagged root is already marked. The exceptions are noise-only, proving simultaneity over starts.
\end{proof}

\subsection{Packing a single wrong-gate boundary}

Fix a gate plane $\Gamma=\{(i_\Gamma,P,b):P\in\Z^2\}$, a critical type/reference phase and finite deterministic transverse footprint. Boundary exceptions have complete witnesses: (R) a raw contour rooted there with $f\ge1$ flags, at most $50(f-1)$ edges, flags at times at most $b$; or (D) a wrong-deep contour with $f\ge11$, the same edge bound and flags at times at most $b-1$, connected by a zero-$\chi$ lag in $\{0,\ldots,W\}$. Intrinsic singleton seeds $S_P$ have joint probability at most $\eta_B^m$ for any $m$ distinct fixed $P$. They may overlap raw flags.

\begin{lemma}[Boundary eroder on the full graph]\label{lem:three-boundary}
Enlarge boundary exceptions by the packed balls below and close by the finitary two-cluster rule, adding flags strictly above $b$; call this field $M_\Gamma$. For every fixed root, the probability that some $N$-vertex full-actual backward path entirely above $b$ has shifted $M_\Gamma$ readings at at least half its vertices is at most
\begin{equation}\label{eq:three-boundary-tail}
 16\exp[-N/(2C_*)].
\end{equation}
Shifts may be any literal helper shift or zero. The same estimate holds for two-sided raw-only $M$, with no boundary floor. Live-data inclusion on the boundary implies domination of all live data and ordinary proposals in a finite region whose chosen parent routes end in that footprint.
\end{lemma}
\begin{proof}
Choose one complete witness for each exceptional boundary point, minimizing flag count then a finite lexicographic code. Process them in decreasing size, keeping one exactly when its flags are disjoint from all previously kept flags. A discarded size-$f$ witness shares a flag with a kept size-$f'\ge f$ witness. Their gate roots are within
\[
 350\cdot50(f+f')\le35000f'<Rf',\qquad R=10^6,
\]
in characteristic distance, because long helper links have zero displacement. Balls of radius $Rf'$ at kept roots cover all exceptions. Only kept complete witnesses are encoded. If their total flag count is $f_B$ and $c_D$ are deep witnesses, then
\begin{equation}\label{eq:three-anchor-cost}
 c_D\le f_B/11.
\end{equation}
No eleven-flag claim is made for a pruned one-charge arm.

Treat a used ball as one terminal group, and intrinsic seeds and future flags as singleton groups. Let $f_+$ count future flags, $m$ intrinsic seeds, $f=f_++f_B$ and $k'=f+m$. There are $g\le k'$ groups. Lemma~10.3's printed gains (unchanged from version~1.3.0 \cite{DMnoise13}) $129,-351$ and coefficient $1440/129$ give the sum form
\[
 129E-1440(g-1)\le\sum_{j=1}^g\Omega_j\le9Rf_B.
\]
Its proof compares each group with its own reference leaf; each of three poles varies by at most $3Rf'$ on a ball. Hence
\begin{equation}\label{eq:three-boundary-edges}
 E<12g+69768f_B<2^{18}k'.
\end{equation}
Here $E$ counts positive edges above $b$. Connect each ball's at most three charge endpoints to its center by unit transverse paths, totaling at most $3Rf_B$ edges, and attach the complete kept witnesses with at most $50f_B$ edges. Neither ball area nor discarded-witness multiplicity is charged.

Inspect marked path occurrences from latest to earliest, attaching an entire grouped contour and skipping strictly earlier than its earliest ordinary terminal time. A contour covers at most $E+2$ occurrences. Keep the final contour whole even past the path prefix. A contour reaching $b$ is final, since every path vertex is above $b$; boundary groups and intrinsic labels thus appear in at most one contour. Future flags have disjoint strict terminal slabs and lie above all boundary flags. Summing gives $N\le2^{20}\sum k'$. The whole actual path, helper spurs, contours, group connectors and full witnesses use fewer than $2^{23}\sum k'$ underlying edges and seed incidences. Fewer than $64(T+2)<2^{15}$ role incidences per item leave fewer than $2^{38}\sum k'<C_*\sum k'$ role-graph edges.

For the prescribed graph \eqref{eq:three-degree}, degree is below $\Delta_{\rm full}(B)$: roles are actual path, shifted root, three charges, group/source center, raw/intrinsic terminal, unit connector and $T+1$ forecast levels. Choices are bounded by $\kappa(B)$, $B+1$ spurs, $q$ providers, two logical inputs and fixed orientations. Encode the connected embedded source test by doubled traversal, role masks, physical flag subset and intrinsic subset. A first deep-anchor traversal selects a lag in at most $W+1\le2W$ ways; known gate time, characteristic and virtual role determine its endpoint. Repeated traversals reuse that edge. Different witness partitions yielding the same annotated graph and source subsets are the same probability event and counted once. A valid decomposition proves its size and \eqref{eq:three-anchor-cost} before that decomposition is forgotten. Traversal, masks, compositions and stopping data therefore give, for fixed $f,m$,
\[
 \#\text{tests}\le K_{\rm full}(B)^{f+m}W^{6c_D}.
\]
The sixfold power overcounts lag orientations and endpoint conventions. Fixed polynomial annotations satisfy $f^j\le2^{jf}$ and fit the edge budget. There is no factorial assignment to group names, no free group-position factor and no lag charge for discarded witnesses. The polynomial $K_{\rm full}$ depends on neither palette, $W$ nor input.

For a fixed test, H\"older with weights $3/4,1/4$ bounds probability by $\eps^{3f/4}\eta_B^{m/4}$. Only the distinct raw flags and the fixed distinct intrinsic experiments require within-family independence. Since $W\le C/\eps$ and $c_D\le f_B/11$, the two weighted bases are
\[
 K_{\rm full}(B)C^{6/11}\eps^{9/44},\qquad
 K_{\rm full}(B)\eta_B^{1/4}.
\]
Both are below $1/64$ by the parameter choices. Summing $(k'+1)64^{-k'}$ for $k'\ge N/C_*$ proves \eqref{eq:three-boundary-tail}. The encoded whole path already pays paths selected after seeing innovations.

For raw-only $M$, repeat strict-slab selection without a lower floor. A complete raw contour covers at most $50f+2\le52f$ marked occurrences, giving $N\le104f$. Its disjoint slabs have distinct flags; there are no anchors or intrinsic sources. The same count is an overestimate. This proves the two-sided version directly, without identifying it with an empty boundary initialization.

For domination choose WEST when $x>x_\Gamma$, otherwise SOUTH, routing northeast sites to $\Gamma$ in exactly the time difference. An unflagged live output has $h$ successful providers, all with live selected-index inputs by bottom-spreading. If these inputs are marked, $h>K$ occupies two boxes and marks the child. Flags above $b$ are seeds. Induction also covers ordinary proposals before pinning. The deterministic boundary footprint includes all routed endpoints, so no unseeded boundary point is used.
\end{proof}

\subsection{Halting and the four boundary conditions}

On a halting input, Lemma~\ref{lem:three-chart} and compactness give an effective $r_n$ such that every locally legal chart patch on $[-r_n,r_n]^2$ contains a wrong short or long critical gate with its full ride and recent query segment. Enumerate finite patches, test internal legality and absence of both finite motifs, and stop at the first empty set. Survivors at all radii would give a full chart with neither wrong gate. This tests no infinite extendibility.

For a target of age $t$ and chart position $(x_0,y_0)$ put
\[
 \begin{gathered}
 m_t=\lfloor t/K_0\rfloor,\quad R_n=10(r_n+W_n^*+B+T+r_F+1),\\
 L_t=100(r_F+1)m_t+R_n,\\
 D_t=4L_t+4r_n+4W_n^*+4B+4T+20,\\
 U_t=(r_F+1)(m_t+3)+(V+1)D_t.
 \end{gathered}
\]
Center the motif box at $(x_0-L_t,y_0-L_t)$. Every comparison, query and routed boundary endpoint has age at least $t-D_t$ and physical radius at most $U_t$: the full late cone has chart excursions at most $(r_F+1)(m_t+2)$, forests have depth at most $D_t$ and transverse increment at most $V$, and the fixed $B,T,W_n^*$ terms cover shifts, forecasts, strips and rides, with a compatibility step.

Here is an explicit burn-in. Set
\[
 \begin{gathered}
 d_1=400(r_F+1)/K_0,\quad d_0=4R_n+4r_n+4W_n^*+4B+4T+20,\\
 u_1=(r_F+1)/K_0+(V+1)d_1,\quad u_0=3(r_F+1)+(V+1)d_0.
 \end{gathered}
\]
Then $D_t\le d_1t+d_0$, $U_t\le u_1t+u_0$, and $d_1+8u_1<10^{-3}$. Take
\begin{equation}\label{eq:three-burnin}
 t_0=2+\left\lceil\max\left\{2K_0,\frac{d_0}{1/2-d_1},
 \frac{d_0+8u_0+2}{1-d_1-8u_1}\right\}\right\rceil.
\end{equation}
For $t\ge t_0$, $m_t\ge1$, $D_t\le t/2$ and $t-D_t\ge8U_t+2$. Let $\mathcal G_t$ exclude $\mathcal D_s(z)$ at every deterministic point with $t-D_t\le s\le t$ and physical radius at most $U_t$ about the target. By \eqref{eq:three-density},
\begin{equation}\label{eq:three-good-box}
 \Prob(\mathcal G_t^c)\le6(t+1)(2U_t+1)^3e^{-(t-D_t)/4212}
 \le P_n(t)e^{-t/8424}
\end{equation}
for an explicit effective polynomial. Comparison spans are at most $4U_t$, so full-age pairing applies with its extra compatibility step. This probability is added to the other bounds, never used as a conditioning event.

The four conditions for the single wrong-gate boundary are explicit:
\begin{enumerate}
\item \emph{Southwest of the whole full cone.} The margin $L_t$ puts the motif, ride, queries and strips southwest of every relevant late point, including all $B$ helper shifts and proposal parents. Their query ages are at least $t/2$. This requires $K_0$ proportional to the full read radius.
\item \emph{One coherent same-history footprint.} On $\mathcal G_t$, an unmarked late reading supplies a clean west/south forest reaching that common box. At least $h-K=256$ clean successes remain at every exposure. Coherence labels its whole southwest region legally, including the motif and routed transverse footprint. The extracted gate is fixed for that entire history. Different histories need not agree.
\item \emph{Complete exceptional witnesses.} Marked live boundary data have full raw witnesses. Unmarked live data agree with the reference. A wrong live current $Q$ has a complete deep witness at lag zero; otherwise survival requires an intrinsic failure or a wrong live $Q$ on the short ride or recent $W$ window. Decreasing-size flag-disjoint packing retains complete witnesses, and only retained deep ones incur lag entropy, with $c_D\le f_B/11$.
\item \emph{One final boundary contour.} Positive future edges pay above-plane span. A contour reaching $b$ is the last selected contour, because the whole late path is above $b$. Boundary flags and intrinsic groups are charged once. A below-plane $W$ lag is never counted as $W$ ordinary path steps.
\end{enumerate}

The interval inclusion in item~3 holds for arbitrary initial helpers. At a short wrong flag-zero gate, its true origin starts a set containing the ideal upper bit. UNKNOWN, union with zero, correct copying and common replacements preserve containment. Fewer than $s$ ideal upper ones imply fewer than $s$ actual singleton-one fields, forcing a pin. At a long wrong flag-one gate, last HIGH marks in the recent $W$ segment preserve the possibility of one under UNKNOWN, union and copy; no origin occurs there since $2W\le\ell_a$. At least $s$ such marks force $h_i\ge s$ and pin. Bottom $Q$ adds zero, bottom $A$ gives UNKNOWN, both favorable to these inclusions; bottom current $Q$ already pins. Thus every surviving live gate belongs to the intrinsic/raw/deep union, a noise-only event once gate type and phase are fixed, without reference colors.

\begin{proof}[Halting direction of Theorem~\ref{thm:three}]
On $\mathcal G_t$, a history with no unmarked live data or proposal in the enlarged late cone uses the raw-only field $M$. Otherwise one unmarked reading supplies a coherent reference and a wrong gate $\Gamma$ as above. Pack its witnesses in a deterministic gate box containing all endpoints routed from the full late cone. Lemma~\ref{lem:three-boundary} dominates all relevant live data and ordinary proposals by its $M_\Gamma$. This is one field for the entire history. Possible gate positions, the two types, finite phase residues determining the intrinsic experiment and reference-root positions form a deterministic finite or polynomial-in-$t$ list. Every candidate field is noise-only; no initial words or chart fillings are enumerated.

If the target map is nonconstant, choose a witnessing pair of starts after fixing innovations. A shared fresh replacement cannot differ. At an unflagged differing full output, a full actual parent differs. Follow such parents for $m_t$ steps in the late cone. At each vertex one history has a live ordinary proposal at one of the $B$ shifted bases: if all $A_t(i+j,P)$ were bottom in both, both outputs would be the common $(\botv,\botv,\mathrm{UNKNOWN}^B)$. Each history has one dominating field for all these proposals, so at least half the vertices are shifted-marked in one of its two fixed fields. Apply \eqref{eq:three-boundary-tail}, including its raw-only form, union over the candidate list, and add \eqref{eq:three-good-box}. These bound existence over all starts without a union over configurations or conditioning on selected histories.

Since $m_t\ge t/K_0-1$, an effective polynomial $Q_n$ gives, uniformly on $I_n^{(3)}$ for $t\ge t_0$,
\[
 q_t(\eps)\le Q_n(t)e^{-dt},\quad
 d=\min\{1/8424,1/(2C_*K_0)\}=1/(2^{41}K_0).
\]
Here $q_t$ is the single-site all-input full-symbol coefficient. If $Q_n(t)\le A_n(t+1)^p$, an effective integer
\[
 C_n>3A_np!(2/d)^p+3^{t_0+1}
\]
works for every $t\ge0$ with $c=d/2$. Indeed $(t+1)^pe^{-dt/2}\le e^{d/2}p!(2/d)^p$, $e^{d/2}<3$, and the last term covers earlier times. A fixed decay coefficient does not mean fast forgetting: the prefactor contains $3^{t_0+1}$, and $t_0$ grows linearly in $r_n$ and $W_n^*$. All coefficients and the burn-in are effective once the halting motif search terminates.

Translation and a finite union bound give the estimate on $F$. Shared finite-window pullback images are nested finite nonempty sets; the summable estimate makes them eventually singletons almost surely. Forward blocks have the same distribution. Feller compactness supplies a stationary law, and coupling it to every initial law proves uniqueness and attraction. Together with the nonhalting proof this establishes the theorem.
\end{proof}

\subsection{The fixed interaction on $\Z^4$}

\begin{corollary}[All Gibbs measures in dimension four]\label{cor:three-gibbs}
There is one computably specified translation-invariant finite-range interaction on $\Z^4$, with finite alphabet and individual energies in $\{0,1\}$, and computable pairwise disjoint nonempty rational open inverse-temperature windows $\mathcal J_n^{(4)}$, such that at every $\beta\in\mathcal J_n^{(4)}$ uniqueness among all Gibbs measures is equivalent to halting of $U(n)$. Halting gives exponential weak spatial mixing in arbitrary finite volumes, uniformly over exterior configurations, including holes and disconnected volumes. Nonhalting gives two Gibbs measures invariant under every translation of $\Z^4$. Mixing constants may depend on $\beta$.
\end{corollary}
\begin{proof}
For the fixed rule $f$, actual neighborhood $N$ and $N_\Sigma=|\Sigma_3|$, put on $\Z^3\times\Z$ the anchored energy
\[
 \Phi_{z,t}(\sigma)=\ind\{\sigma_{z,t}\ne f((\sigma_{z+u,t-1})_{u\in N})\}.
\]
Include every factor meeting a finite volume, including exterior-child factors. The normalizer $1+(N_\Sigma-1)e^{-\beta}$ is parent-independent; dividing by it leaves every DLR specification unchanged and gives exactly replacement at
\[
 \eps(\beta)=N_\Sigma/(e^\beta+N_\Sigma-1).
\]
For $I_n^{(3)}=(a_n,b_n)$ enumerate positive rationals $c_n<d_n$ until $a_n<\eps(d_n)<\eps(c_n)<b_n$ is certified by computable exponential approximations. Strict monotonicity and openness guarantee termination. Set $\mathcal J_n^{(4)}=(c_n,d_n)$; the windows are nonempty and pairwise disjoint. Their endpoints are computable but far too large to print: $N_\Sigma\ge100^{4^{43429}}$ and their scale is $\log N_\Sigma+\log(1/\eps)$.

All hypotheses of \cite[Section~2.1 and Corollary~13.7]{DMst} hold: the alphabet has at least two letters; $N$ contains zero because $S_0$ contains zero and SOUTH has logical index $i$; entries are at least $\eps/N_\Sigma>0$; retain-or-replace basic maps are independent; spatial dimension is $d=3\ge1$; and the coefficient of its equation~(4) satisfies $q_t\le C_ne^{-ct}$ for all $t\ge0$ on halting windows. Corollary~13.7, with estimates (7)--(8), therefore gives uniqueness among all DLR measures and arbitrary-volume exponential weak spatial mixing. No invariance is assumed of competing measures. On nonhalting windows \cite[Lemma~2.4]{DMst} sends the two spatially invariant stationary laws to distinct stationary-history Gibbs measures, distinguished by time-zero marginals and invariant under all four translations.
\end{proof}

No dimension-three analogue of Theorem~B.1, Theorems~9.6--9.7 or Corollaries~9.8--9.12 is asserted. These new results, and those of Section~10, are not formalized in Lean.

\setcounter{equation}{\threeoldequation}
\endgroup

\section*{Open problems}
\addcontentsline{toc}{section}{Open problems}
The fixed-rule noise diagram remains open on $\Z^2$ and $\Z$. For equilibrium models, a static finite-energy interaction on $\Z^3$ or $\Z^2$ with halting-labelled all-Gibbs uniqueness windows remains open; the interaction here is a space-time model on $\Z^4$.

For these three constructions, behaviour between the specified windows remains undetermined. Whether nonhalting windows have continuum many extreme translation-invariant stationary laws or Gibbs measures is separate from the extremal conclusion without invariance in \cref{thm:continuum}. The full positive-rational ergodicity set has the recalled $\Pi^0_2$ upper bound of \cite[Remark~10.9]{DMst}; its exact complexity remains open. Robustness under arbitrary unnormalized energy perturbations remains open. No robustness, extremal-continuum, field-response or further effective theorem is asserted for either lower-dimensional rule beyond its proved conclusions.

\section*{How this work was done}
\addcontentsline{toc}{section}{How this work was done}

\paragraph{AI note.} GPT-6.1 Sol derived the construction and proofs in a long working session in which agents of the same model family audited parts of the argument; Claude Opus 5.5 checked the result refutation-first in three reviews from the other developer. GPT-6.1 Sol also searched the literature and wrote this manuscript. Claude Opus 5.5 then checked the written text refutation-first against the derivation and the three reviews. GPT-6.1 Sol then added Corollary~B, which Claude Opus 5.5 checked in a further refutation-first pass. For version~1.1.0, GPT-6 Astra and GPT-6.1 Sol, directed by Nidhal Mghirbi, found the corrections and new results; Claude Opus 5.5 verified that package refutation-first and independently recomputed both scale choices. GPT-6.1 Sol integrated the verified package into version~1.1.0, and Claude Opus 5.5 read the integrated text back refutation-first. The Lean development was written by Claude Opus 5.5 agents. The text of version~1.2.0 was written by GPT-6.1 Sol. For version~1.3.0, the lower-dimensional results were found in a GPT-6.1 Sol research run and proved again in an independent blind GPT-6.1 Sol replication. Claude Opus 5.5 then checked both lower-dimensional results refutation-first across model families; its citation, scope, finite-history exposition and scale repairs are incorporated here. For version~1.4.0, the dimension-three result was found in a GPT-6.1 Sol research run, checked refutation-first by Claude Opus 5.5, and proved again by an independent blind GPT-6.1 Sol replication with a different construction. GPT-6.1 Sol wrote the repaired construction presented here. The authors chose the question, directed the work and take responsibility. Models are not listed as authors.
That check re-derived each load-bearing step from the source papers before reading either run's argument.

\section*{Verification status}
\addcontentsline{toc}{section}{Verification status}

Theorem~E on $\Z^3$ and Corollary~F on $\Z^4$ have written proofs in \cref{sec:three} and are not formalized. A cross-family refutation-first check and an independent blind replication with a different construction both support the new result. The written proof incorporates the direct seed-routing check, full-age coherence, the published grouped-contour constants, complete flag-disjoint boundary packing and the four single-gate conditions. Two independently written exact-arithmetic Python and Node interval certificates check its least $k$, definite $\eps_*$, least $n_0$ and exact $c$, with recorded outputs in \texttt{verification/z3-scale/}. The original Lean appendix is byte-for-byte unchanged, and no Lean run was performed for this release.

Theorem~C on $\Z^4$ and Corollary~D on $\Z^5$ have written proofs in \cref{sec:lower} and are not formalized in Lean. The lower-dimensional proof uses the published physical contour and all-path density lemmas of \cite[Lemmas~12.3--12.4]{DMforgetLower}; it proves the characteristic contour, finite-history coherence, all-input full-symbol estimate and all-DLR transfer needed here. The original Lean appendix is byte-for-byte unchanged, and no new Lean run is claimed. The supplied cross-family check confirms both lower-dimensional results; its repairs have been applied. The new least scale parameters and binding threshold were separately certified by two independent exact-arithmetic implementations.

Theorem~A (in the form of \cref{thm:main}), \cref{cor:undecidable} and the one-rule corollary after Theorem~A are now proved end to end in Lean; \cref{app:lean} gives the formal statements, scope, proof-route differences and release checks.

The construction and component proofs were derived with separate agents re-deriving and challenging them. The first review, from the other developer, re-derived the contour, coherence, birth and pullback arguments of Sections~5 and~8 in a general form, including the grouped potential estimate and the complete-history chart argument. The second review checked the fixed-rule assembly, sensor ownership, direct machine, parameter inequalities, feedback memory estimate, historical compression and shifted-target coupling. It identified the missing written proof for finite-board forcing in exceptional phases. The repaired lemma was then re-derived from the rules of \cite[Section~4]{DMhard} in the third review. The subsequent refutation-first check of the version~1.0.0 text found no mathematical error at that stage and identified credit and presentation changes; the later audit described in the AI note identified the defects corrected in version~1.1.0, and the cross-family verification confirmed them.

The scaffold statements of \cref{prop:scaffold}---desubstitution with a unique tower for every tiling, local phase determination, frames, and the finite-rank red-path and ruler-bit lemmas---are formalized in Lean in \cite[Appendix~D, version~1.2.0]{DMhard}, without \texttt{sorry}. The formal results include \texttt{desub}, \texttt{tower\_unique}, \texttt{phase\_local}, \texttt{red\_col}/\texttt{red\_row}, the \texttt{rank\_bits} lemmas and the frame and laminar lemmas. Their trust base is \texttt{propext}, \texttt{Classical.choice} and \texttt{Quot.sound}. Quarter encoding (\texttt{quarter\_exists}, \texttt{quarter\_unique}) is formalized there as well; the typical-phase routing existence is a proved import from the Lean core of \cite[Propositions~5.1 and 5.7 and Lemma~5.3]{DMhard}, as described in \cref{app:lean}; the black-type comparison is proved in \cref{sec:chart}. The finite-rank statements used in \cref{lem:allphase} are among these formalized inputs and apply even when the other axis has infinite rank.

The first review re-derived the pole/frontier argument and grouped estimates. The second and third reviews checked their unequal-time and zero-potential-link use; they did not repeat all of the original fine combinatorics. The check of the written text re-derived the frontier invariant and the numerical estimates. The Wang set $R$ and the physical alphabet were not enumerated; with the constants of \cref{sec:parameters} they are far too large to list. Nonhalting satisfiability follows from the cited typical-phase construction and the finite computations in \cref{prop:chart}. Corollary~B was added after the check of the written text. A further refutation-first check by the other model family re-derived its proof, checked each hypothesis of \cite[Lemma~2.4 and Corollary~10.4]{DMst} for this automaton, and recomputed the scale of \cref{sec:parameters} with an independent program. The additional conclusion about uniformly recurrent and periodic Gibbs measures was checked against \cite[Theorem~10.10(i)]{DMst} in version~1.0.0; version~1.1.0 drops it because uniqueness among all Gibbs measures subsumes it. For version~1.1.0, a further cross-family refutation-first verification re-derived corrections~(i)--(iv), finite-history argument, extremal cardinality, robustness, field response and effective corollaries, and independently recomputed both scales. It treated the earlier scaffold and stochastic lemmas as inputs, spot-checking their uses, and checked the statement and hypotheses of \cite[Corollary~13.7]{DMst} rather than re-proving it. These are shared dependencies of the verification. Correction~(v) was checked when the integrated text was read back, and its count is proved in Lean (\cref{app:lean}). Theorem~A (Theorem~9.3) and Corollary~9.4 retain their printed statements and original constants. The completed Lean core is reported in \cref{app:lean}, which lists what it does not cover, including Corollary~B and the results added in version~1.1.0. The checks reported there belong to the Lean release package; this text update performed no alphabet enumeration or new Lean run. The authors' role is selection of the question and direction of the work; no additional human line-by-line verification is claimed.

\appendix
\crefalias{section}{appendix}
\section{Finite histories and exponential coalescence}\label{app:finite}

We prove \cref{thm:exponential} quantitatively. To also prepare the robustness proof, allow an independent exceptional flag $A$ of probability $\delta$ at each update, selecting any random local map when present. The baseline maps are used otherwise. All update randomness is independent across space--time. Assume
\begin{equation}\label{eq:overlayconditions}
 0\le\delta\le\min\{\eps_n^-,\widehat\eta_B/(2BW_n^*),1/(64K_{\rm rare})\}.
\end{equation}
The original maps correspond to $\delta=0$. Write $E'=E\vee A$,
$e'=\eps+(1-\eps)\delta\le2\eps$, $A'=4V+1$, and $J=W_n^*+52$.
Define $\MM$ with $E'$ in place of $E$. This two-sided flag field may use auxiliary negative-time flags even though the trajectory starts at zero. Every marked point has a finite flag derivation; an unmarked output has marked Moore parents in at most one provider cluster.

\begin{lemma}[A finite motif certificate]\label{lem:motif}
\proved\ If $U(n)$ halts, a terminating finite search finds $r\ge1$ such that every Moore-compatible live patch on $Q_r=[-r,r]^2$ contains one of the following motifs. A type S motif is a horizontal segment of $\ell_{a-1}+1$ letters, with $\Orig$ at the left endpoint and at no other position of the segment, and $\End$, $p=0$, $H_{\rm row}=0$ at its right endpoint. Type L has length $\ell_a+1$ and right flag one instead.
\end{lemma}
\begin{proof}
By \cref{prop:chart}, every full legal chart has a wrong bottom row at depth $a-1$ or $a$. Its finite ride has only its endpoint walls. Indeed the closed depth-$d$ frame avoids residues $0,1$ modulo $4^{d+1}$, so its relevant ranks are finite, as in \cref{lem:allphase}. For $e<d$, the bottom row has residue $1$ modulo $4^{e+1}$ relative to $\tau_2$, so lies in no shallower row frame. Shallower vertical walls are therefore absent; same-depth frames are disjoint, and strict laminarity excludes a deeper column endpoint inside the frame. The only origin and end on the ride are its endpoints, at distance $\ell_d$. Its phase is zero since $d\ge k$. This gives a motif in every full chart, including exceptional phases.

Enumerate all assignments from the effective finite alphabet $\mathcal L_R$ on $Q_r$ and reject internal incompatibilities and either contained motif. Each test is finite, using the finite-approximant phase decoder. If a survivor existed for every $r$, extend it arbitrarily outside the box and take a product-compactness limit. Every compatibility and motif test eventually lies inside the boxes; the limit is a legal full chart avoiding both motifs, a contradiction. The first empty test therefore terminates on halting inputs. Its termination is a semidecision, not an infinite-extension test.
\end{proof}

\begin{lemma}[Finite-time coherence within one trajectory]\label{lem:finitecoherence}
\proved\ In one fixed trajectory starting at time zero, two unmarked live data or second-coordinate readings at time $s$, at equal or Moore-neighbouring planar sites and at provider $\ell^1$ distance $d$, agree or are compatible, respectively, whenever $s\ge d+1$.
\end{lemma}
\begin{proof}
At $d=0$, $s\ge1$, an unmarked live reading is unflagged, so its letter is the live ordinary proposal; data and second coordinate agree when both are live. For planar neighbours the retained offset sets overlap in at least $2h-q=3K$ offsets. A check at a common offset certifies compatibility of their retained centres. Neither argument imposes a condition on time zero.

For $d>0$, reduce one component of provider separation by one and pair offsets by that unit shift in each provider cube. There are $5K(15/16)$ injective pairs. Each root excludes at most $K$ nonretaining offsets and at most $K$ with marked Moore parents. At least $(11/16)K=45056$ pairs remain. A remaining pair gives unmarked live data parents with the respective labels at time $s-1$, provider distance $d-1$, and the same planar displacement. The induction condition $s-1\ge(d-1)+1$ holds. This proves the claim within the fixed trajectory.
\end{proof}

\begin{lemma}[Exponential unmarked-liveness bound]\label{lem:unmarkedfinite}
\proved\ For a fixed site at time $t$, the probability that some initial configuration produces an unmarked live data or second-coordinate reading there is at most
\[
 C_U\exp[-t/(A'J)].
\]
The prefactor is computable from the motif radius and is uniform over $\eps\in I_n$ and the exceptional maps satisfying \eqref{eq:overlayconditions}.
\end{lemma}
\begin{proof}
Put $R_*=r+T$. From a putative unmarked live root $(u_0,P_0,t)$ propagate $R_*$ ordinary steps backwards, choosing a retained provider outside the possibly marked cluster at each step. At least $h-K$ such providers exist and supply nine unmarked live data parents. This gives witnesses throughout $u_0+Q_{R_*}$ at time $t-R_*$, with providers in $P_0+[0,VR_*]^4$. If $t-R_*\ge4VR_*+1$, \cref{lem:finitecoherence} makes their letters a compatible reference patch. Each witness has a letter-preserving centre-parent path at its planar site, with provider at time $\sigma$ in $P_0+[0,V(t-\sigma)]^4$.

Choose a motif in $u_0+Q_r$, with right endpoint $g$, and follow its witness's unmarked centre path for
\[
 N=\lfloor t/(2A')\rfloor
\]
vertices, beginning at time $t-R_*$. At a sampled time $b$ and lag $\ell$, the deep query reads data at $\sigma=b-\ell-1$. Its actual leaf providers are
\[
 P_b+\ell v-(T-1)v+\sum_{i=1}^T s_i,\qquad s_i\in S.
\]
Relative to $P_0$ these lie coordinatewise between $-VT$ and $V(t-\sigma+T)$. Reference providers at time $\sigma$ lie between $0$ and $V(t-\sigma)$, so their $\ell^1$ distance is at most $4V(t-\sigma+T)$. All query sites lie in the motif, and the full virtual cone is inside $Q_{R_*}$. Type S needs lags at most $\ell_{a-1}<W$; type L needs lags at most $W\le W_n^*<\ell_a$. Here $W\le C/\eps<\ell_a/2$ gives $W\le W_n^*$; in particular no origin is in the long recent window.

It suffices to take the computable burn-in
\begin{equation}\label{eq:burnin}
 t\ge2[A'(R_*+W_n^*)+4VT+2].
\end{equation}
Since $t-\sigma\le R_*+N+W_n^*$, this gives
$\sigma\ge4V(t-\sigma+T)+1$ as well as the patch condition.
Thus every unmarked live leaf agrees with its reference letter by within-trajectory coherence. A wrong live query can be traced using retained centres at all $T$ virtual levels; these keep its planar site and label. Its actual leaves are consequently marked. Append their finite $\MM$ derivations. The formal-clock argument of \cref{lem:deep}, using \cref{lem:pole}, gives a contour with $f_i\ge11$ distinct $E'$ flags, at most $50(f_i-1)$ support edges, and earliest flag time
\[
 t_{\min}\ge b-\ell-1-50(f_i-1).
\]
The derivations may reach auxiliary negative times; their flags are still paid for by their probabilities.

At a sampled vertex, live unmarked data implies a live $A_b$ and $\Pin=0$, hence a live current deep driver. A wrong current driver supplies lag zero. If no wrong live deep instruction occurs on the relevant ride, define $\mathcal X$ to be its baseline intrinsic bad event or any exceptional flag on its lane atoms. Its probability is at most
$\widehat\eta_B+BW_n^*\delta\le2\widehat\eta_B$.
Absent $\mathcal X$, a short motif forces erasure: its origin starts with a set containing zero, bottom ordinary instructions give UNKNOWN, bottom deep instructions insert zero, correct nonorigins copy, and replacements contain their upper scalar mark. Induction makes each actual interval contain the ideal upper bit. Fewer than $s$ upper ones implies fewer than $s$ singleton ones. For a long motif, at least $s$ recent last-HIGH marks write singleton ones. Subsequent correct copies, unions with zero and UNKNOWN preserve the possibility of one. This also forces erasure. Both contradict liveness, so $\mathcal X$ or a wrong live query must occur.

Greedily inspect the $N$ path vertices. For $\mathcal X$, advance one centre step. For a deep contour record its lag and restart one step before $t_{\min}$. Keep the whole contour even if it passes the prefix's lower end, and stop the path there. Different contours' flag times are disjoint. With total flags $f$, contours $c$ and intrinsic tests $m_0$, we have $c\le f/11$. At the same planar endpoint, distinct sampled times have disjoint lane atoms: an atom has planar coordinate $g_x-\ell$ and time $b-\ell$, so equality forces the same $\ell$ and $b$, regardless of the provider path. Hence, for each fixed certificate, the $m_0$ tests are independent.

For completeness the finite encoding includes the truncated last piece. A contour of size $f_i$ has at most $2W$ lags and restart gap at most $52f_iW$. A gap $h_*$ has at most $(Vh_*+1)^4$ provider endpoints. Lag, gap and endpoint together cost at most
\[
 2[54(V+1)]^5 f_i^5 W^6.
\]
A truncated unused endpoint may be recorded as zero. Raw contour traversals and their flag subset cost at most $(2\Delta^2)^{51f_i}$ by \cref{lem:counts}; exact physical times and virtual levels use the same degree bound. Compositions cost $2^f$, interleavings $2^{f+m_0}$, and stop/source annotations fit its slack, using $f_i^5\le2^{5f_i}$ and $q$ choices at an intrinsic step. Thus the number of tests is at most
\[
 K_{\rm hist}^{f+m_0}W^{6f/11},\qquad
 N\le J(f+m_0).
\]
Raw and statistical tests may overlap. H\"older gives joint weight
$(e')^{3f/4}(2\widehat\eta_B)^{m_0/4}$. Since $W\le C/\eps$, the two weighted bases are
\[
 K_{\rm hist}2^{3/4}C^{6/11}\eps^{9/44},
 \qquad K_{\rm hist}(2\widehat\eta_B)^{1/4},
\]
both below $1/32$ by \eqref{eq:Bsearch} and \eqref{eq:epsstar}. Summing masses $k=f+m_0\ge\lceil N/J\rceil$ gives
\[
 \sum_{k\ge\lceil N/J\rceil}(k+1)32^{-k}
 \le(16/15)16^{-\lceil N/J\rceil}.
\]
Take the deterministic union over at most $2|Q_r|(VR_*+1)^4$ endpoint roots and types at time $t-R_*$. Patch fillings are not counted or conditioned on: once these roots and a certificate are fixed, only their specified flags and lane experiments are tested. This is a noise event covering existence of any start. Since $\log16/2>1$, the tail is a computable multiple of $e^{-t/(A'J)}$. Enlarge that multiple for times below \eqref{eq:burnin}.
\end{proof}

\begin{lemma}[Late-cone domination]\label{lem:bulk}
\proved\ Let $r_F\ge1$ bound every read offset, and $m_t=\lfloor t/4\rfloor$. Except on a noise event of probability at most a computable polynomial in $t$ times $e^{-t/(2A'J)}$, simultaneously for every initial row: the relevant live ordinary proposals along the top $m_t$ vertices of a backward dependency path from a fixed time-$t$ target belong to the grouped rare-birth field of \cref{lem:rare}, with $E'$ as flags.
\end{lemma}
\begin{proof}
The finite-prefix version of \cref{lem:rawpaths} bounds an ordinary marked path of $L_*$ vertices from a prescribed root by
\[
 \sum_{f\ge\lceil L_*/50\rceil}(K_{\rm mem}e')^f.
\]
Indeed retain each whole attached contour, including the final truncated piece; it covers at most $50f_i$ vertices, the flags have disjoint time slabs, and the path and support count is at most $101f$. The same degree and traversal count applies. The base is below $1/64$.

Take the part of the target's backward cone with
$\lceil t/2\rceil\le s\le t$, of size at most $(t+1)(2r_Ft+1)^6$.
Union-bound \cref{lem:unmarkedfinite} and the marked-path bound with
$L_*=m_t+1$ at every point. Outside the resulting event every start has all live readings marked there and no marked ordinary path of $m_t$ edges there. The unmarked-liveness exponent is at least $1/(2A'J)$, and the raw-path exponent at least $\log64/200$.

A live data parent of a relevant ordinary proposal has time at least $t-m_t$. Recursively stop it at an $E'$ flag and otherwise use its retained live centre parents. Through depth $m_t$ all sites remain in the cone portion because $t-2m_t\ge\lceil t/2\rceil$. Every node is marked. A branch of $m_t$ edges is excluded, so finite branching gives a finite ancestry terminating at actual $E'$ flags. Induction puts the parent in the minimal scalar $W^{\rm live}$ and the live proposal in $N^{\rm live}$ one time earlier. This establishes domination for the finite trajectory before the rare-birth estimate is used. The grouped seed proof of \cref{lem:rare} then applies with seed weight
$(N_0(e')^M)^{g/16}$.
\end{proof}

\begin{lemma}[Finite discrepancy packing]\label{lem:finitepacking}
\proved\ Under \eqref{eq:overlayconditions}, the all-input full-symbol coefficient is at most
$C_ne^{-t/(4A'J)}$, with a prefactor computable on halting inputs, uniformly over the window and exceptional maps.
\end{lemma}
\begin{proof}
Outside the bad event of \cref{lem:bulk}, two different target outputs give a backward discrepancy path of $m_t$ vertices. At a direct exceptional vertex, record that $A$ flag as a singleton certificate. Otherwise a baseline replacement agrees across starts. A differing deterministic output needs a live $A_s(w-j\mathbf e_y)$ in one run for some $j<B$: if all these proposals are bottom, data and second coordinate are bottom and every helper is UNKNOWN in both runs. The chosen base is in the grouped rare field at actual time $s-1$.

Attach a grouped contour at that actual time and keep its group times unchanged. If its earliest group time is $v_*$, resume the discrepancy path at physical time $v_*-1$, or stop at the prefix end. The connector from time $s$ to $s-1$ is outside the pole contour. Its time change is paid for explicitly. The contour has
\[
 e\le H_0g_i-A_0,\qquad A_0=4000/85>2,\quad H_0=11000/85<130.
\]
It covers at most $e+2\le H_0g_i$ path vertices, including connector and restart. A seed-only contour satisfies the same bound. Its support edges drop physical time by one; group stars preserve time, with oscillation $7000$. Thus this application meets the repaired \cref{lem:pole}.

For $g$ groups and $h_*$ direct exceptions, all selected seed slabs and direct-flag times are disjoint, and
\[
 m_t\le H_0g+h_*.
\]
The connected graph uses fewer than $(2H_0+6)g+h_*<512(g+h_*)$ edges. Project onto physical vertices and one auxiliary centre per group. Edges are actual dependency offsets, ordinary support offsets, $B$ literal shifted-root offsets with time decrement one, and same-time cube stars. A direct exception uses an existing physical source role. The role and degree bounds of \cref{lem:counts} still give $K_{\rm rare}^{g+h_*}$ tests. No new $B$ factor is omitted.

The group parity bound and the disjoint direct-flag times give weight
$(N_0(e')^M)^{g/16}\delta^{h_*}$. Both weighted bases are at most $1/64$:
\[
 K_{\rm rare}(N_0(e')^M)^{1/16}
 \le\frac{(2N_0)^{1/16}}{64(N_0+1)}<1/64,\qquad
 K_{\rm rare}\delta\le1/64.
\]
Here $(e')^M\le e'\le2\eps$ and \eqref{eq:epsstar} were used. Summing $g+h_*\ge\lceil m_t/H_0\rceil$, and using $k+1\le2^k$, gives an exponential tail with rate at least $\log32/(4H_0)$. This exceeds $1/(2A'J)$, as does the raw-path rate in \cref{lem:bulk}. The cone factor is polynomial, so halving the slowest exponent absorbs it, giving $1/(4A'J)$. All finite prefactors, small-time allowances and this polynomial absorption are computable from $r$ and the construction constants. Every test is a noise event; existence of any pair of disagreeing starts implies one of these tests. This proves the all-input quantifier.
\end{proof}
\section{A computable robustness radius for local kernels}\label{app:robust}

This is robustness in stochastic tables and their normalized space--time interactions. Its precursors are \cite[Corollary~8.12 and Remark~8.14]{DMforget} and \cite[Proposition~4.5]{DMst}. Arbitrary perturbations of unnormalized equilibrium energies have different normalization factors and are outside this statement.

\begin{theorem}[A radius preserving both labels]\label{thm:robust}
\proved\ On local stochastic tables on $\Sigma$ with the same neighbourhood $N_F$, use
$d(P,Q)=\max_\eta d_{\rm TV}(P(\cdot\mid\eta),Q(\cdot\mid\eta))$.
Put
\begin{align*}
 \delta_n&=\min\{\eps_n^-,C/W_n^*,
       \widehat\eta_B/(2BW_n^*),1/(64K_{\rm rare}),1/2\},\\
 r_n&=\delta_n\eps_n^-/(2N_\Sigma).
\end{align*}
The radius is positive and computable from $n$ without its halting answer. For every $\eps\in I_n$ and $d(P_\eps,Q)\le r_n$, $Q$ is strictly positive. If $U(n)$ halts, an i.i.d.\ representation of $Q$ has the exponential all-input bound of \cref{thm:exponential}, uniformly over these $Q$ and $\eps$. If $U(n)$ does not halt, $Q$ has two translation-invariant stationary bit laws with error less than $1/31$, and continuum many extremal stationary laws. Its own normalized interaction $-\log Q(c\mid\eta)$ has all-Gibbs uniqueness and arbitrary-volume exponential weak spatial mixing in the halting case, and two translation-invariant Gibbs measures and continuum many extremal Gibbs measures otherwise.
\end{theorem}
\begin{proof}
For each entry, total variation gives
\[
 |Q(c\mid\eta)-P_\eps(c\mid\eta)|
 \le r_n\le\delta_n P_\eps(c\mid\eta)/2,
\]
because $P_\eps(c\mid\eta)\ge\eps_n^-/N_\Sigma$.
Thus $Q\ge(1-\delta_n)P_\eps>0$ entrywise, and
$Q=(1-\delta_n)P_\eps+\delta_n Q'$ for a stochastic table $Q'$.
At each update choose an independent flag $A$ of probability $\delta_n$; when present sample a local map for $Q'$ by independently drawing its value for each parent word. Otherwise use the baseline map. This is an i.i.d.\ representation of $Q$.

The augmented raw flag $E'=E\vee A$ has probability $e'\le2\eps$. The augmented single-gate intrinsic event has probability at most
$\widehat\eta_B+BW_n^*\delta_n\le2\widehat\eta_B$.
All conditions of \eqref{eq:overlayconditions} hold, so \cref{lem:motif,lem:finitecoherence,lem:unmarkedfinite,lem:bulk,lem:finitepacking} prove the halting bound. Their historical bases are explicitly
\[
 K_{\rm hist}2^{3/4}C^{6/11}\eps^{9/44}<1/32,\qquad
 K_{\rm hist}(2\widehat\eta_B)^{1/4}<1/32.
\]
Direct exceptional vertices are charged by
$K_{\rm rare}\delta_n\le1/64$ as in \cref{lem:finitepacking}.
Thus the proof covers parent-dependent exceptional maps without treating them as common replacements.

For nonhalting, fix a correct typical chart and the ideal helper initialization driven by baseline negative-time noise, as in \cref{sec:memory}. A raw characteristic bin now has probability at most
$W\eps+W\delta_n\le C+W_n^*\delta_n\le2C$.
Enlarge each statistical bin event by any exception on the lane supports of its finite gates. Those atoms lie in its own or preceding time bin on $B$ lane characteristics, so there are at most $2BW_n^*$ such atoms. Its probability is at most
$\widehat\eta_B+2BW_n^*\delta_n\le2\widehat\eta_B$.
Raw bins have disjoint supports. Statistical supports on different base characteristics are disjoint by \cref{lem:ownership}, and every third bin on one characteristic is disjoint. The proof of \cref{lem:joint} therefore applies with
\[
 \gamma_{\rm tube}=\sqrt{2C}+(2\widehat\eta_B)^{1/6},\qquad
 K_{\rm mem}\gamma_{\rm tube}
 <\frac{\sqrt2}{128}+\frac{2^{1/6}}{64}<\frac1{32}.
\]

The deterministic error derivation is unchanged. A replacement or exception at the output is a raw source. At an unflagged output the original rule acts. A false erasure either has a recent bad driver or, if a relevant lane is exceptional, the augmented statistical event. With no such exception, the original ideal-helper comparison holds. Thus \cref{lem:decomposition,lem:memorycount} give the same witness count and a geometric error sum
\[
 e_{\rm tube}=
 \frac{K_{\rm mem}\gamma_{\rm tube}}{1-K_{\rm mem}\gamma_{\rm tube}}
 <1/31.
\]
This is weaker than the original-kernel bound $e_{\rm base}<3/125$ and gives separation greater than $29/31$, rather than $119/125$.

Spatial averaging followed by temporal Cesaro limits gives two translation-invariant bit laws. Keeping each chart fixed instead gives an uncountable separated family, so the full proof of \cref{thm:continuum}, now with $e_{\rm tube}$, gives the extremal continuum. Stationary trajectories are DLR states by \cite[Lemma~2.4]{DMst}. On halting inputs $Q$ has finite alphabet, the same finite neighbourhood containing zero, strict positivity, and the exponential all-input coefficient just proved. Applying \cite[Corollary~13.7]{DMst} to its normalized interaction proves all-Gibbs uniqueness and both arbitrary-volume estimates. All quantities defining $r_n$ use only the fixed effective construction and $n$, so their computation needs no halting decision.
\end{proof}

% Scope and checks: SCOPE.md, REPORT.md, STATEMENT-READ.md, STATUS.md,
% UncomputableNoiseDiagram/Statements.lean and Main.lean, README-ARXIV.md.
\section{Lean formalization}\label{app:lean}

Theorem~A (in the form of \cref{thm:main}), \cref{cor:undecidable} and the one-rule corollary after Theorem~A are proved end to end in Lean~4 \cite{Lean4} with its mathematical library Mathlib \cite{Mathlib20}. The proofs contain no unproved step (\texttt{sorry}), take no named hypotheses, and use no axioms beyond \texttt{propext}, \texttt{Classical.choice} and \texttt{Quot.sound}. This appendix says what is stated, what is outside the formalization, where the proof routes differ, and how the proofs were checked.

\paragraph{The statements.} The statement file \texttt{Statements.lean} defines the cellular automata, replacement noise, stationary laws, ergodicity, translation invariance, common-noise pullbacks and rational windows that occur in the results. Its definitions carry comments identifying the printed passages they render. The final proofs of its three propositions are in the development's \texttt{Main.lean}:
\begingroup\small
\begin{verbatim}
theorem UND.theorem_A : UND.TheoremA
theorem UND.corollary_9_4 : UND.CorollaryNineFour
theorem UND.one_rule_undecidable : UND.OneRuleUndecidable
\end{verbatim}
\endgroup
The first is Theorem~A in the form of Theorem~9.3, including the finite-window pullback conclusion; the second is Corollary~9.4; the third is the sentence after Theorem~A, using \texttt{UPD.Compiler.univMachine\_universal} of the hard-model core \cite{DMhard}.

\paragraph{Choices.} Theorem~9.3 is formalized in its existential form, as in Theorem~A: for every machine $U$ there exist an automaton $A$ on $\Z^6$ and rational endpoint maps $\mathit{lo},\mathit{hi}$. They satisfy \texttt{UND.NoiseDiagram U A lo hi}, the statement checked by the comparator. The proof's witnesses are the automaton~\eqref{eq:proposal}, \texttt{UND.theCA U}, and the intervals~\eqref{eq:windows}, \texttt{UND.Param.loF} and \texttt{UND.Param.hiF}; their identification with these printed objects is read from the definitions. The checked statement does not name these witnesses.

The formal Theorem~A holds for every machine $U$, since the construction does not use universality; Corollary~9.4 assumes undecidability of its halting set on positive inputs. The finite alphabet is presented as $\{0,\ldots,|\Sigma|-1\}$. Rational endpoints have computable presentations by pairs of natural numbers $(p,q)$ with $q>0$, and the windows lie in $(0,1)$. Undecidability is the promise problem: no computable answer is correct on every presentation of a rational rate in a window. The adjective ``computably specified'' for one fixed finite automaton is not rendered as a separate condition; the development gives the automaton by explicit definitions and claims no computable family $U\mapsto A_U$ in the statement.

\paragraph{What is proved.} Every step of Sections~2--8 and of the parameter verification of Section~9 used by the three results is proved. The construction parameters, including $k=52386$, are unchanged. In particular, the formal neighbourhood has the corrected $B$ helper sites of \cref{prop:local}. The bound \texttt{card\_nbhdP\_le} proves $|N_F|\le B(n_A+n_Q)+B\le B(n_A+n_Q+2)$, and \texttt{prop\_2\_1} gives the unchanged two-direction bound $\kappa$. The $j=0$ helper parent is counted once.

Two earlier Lean cores are reused, on the same pins: the core \texttt{FixedNoiseErgodicity} of the forgetting paper, version~1.4.0 \cite{DMforget}, and \texttt{UndecidablePhaseDiagram} of the hard-model paper, version~1.2.0 \cite{DMhard}. The former supplies \texttt{FNE.flagLaw}, \texttt{FNE.WeakTendsto}, and process and edge-count layers. The latter supplies the machine and halting definitions, the existence of a universal machine, and the scaffold, tower and background-chart layers. The typical-phase routing existence used in \cref{prop:chart} is a proved import: the routing fields are \texttt{UPD.BG.cfg} unchanged, with (O1)--(E3), activity and carriers supplied by the hard-model core's proved lemmas. The release check compiled from source the four imported modules of the forgetting core and the 36 of the hard-model core; these sources are byte-identical to the released cores.

\paragraph{Proof routes.} The statements are held as printed, with the existential form noted above. Four proof routes differ from the text; the labels below are those of the development's fidelity record.
\begin{itemize}
\item F17: the translation-invariant chart law in \cref{prop:chart} is the pushforward of Haar measure along $\tau\mapsto\texttt{chartOf}\,\tau$, rather than a weak limit of spatial averages of one chart.
\item F19: the copied-arm witnesses in \cref{lem:memorycount,lem:fullcoupling} keep every group link and count them linearly, at most $25g$, rather than extracting the printed $5g$ star sub-witness. The encodings still feed the printed union bounds: at most $256g$ edges against $K_{\rm mem}$ and $512g$ edges against $K_{\rm rare}$.
\item F21: the proof of \cref{lem:fullcoupling} bounds a live $A_t$ by the raw marked set $\MM$ through \cref{lem:eliminate}. It covers discrepancy paths by the raw contours of \cref{lem:rawpaths}, with at most $101f\le512f$ edges for $f$ flags, into $K_{\rm rare}^f$, and base $K_{\rm rare}\eps<1$. It builds no rare-birth contours, so version~1.1.0's birth-contour clock correction is not exercised. Consequently \cref{lem:rare} is not used or formalized (F22). The printed proof of \cref{lem:fullcoupling} is retained here.
\item F23: the historical queries of \cref{lem:eliminate} are packed by their number $r$ of blocks, with H\"older weights using the printed exponents. Their block weights sum to at most $\beta^r$, $\beta<1$, instead of the printed $(f+m_0)$-indexed count.
\end{itemize}

\paragraph{Not formalized.} The following are outside the scope of this formalization:
\begin{itemize}
\item Corollary~B (the fixed interaction on $\Z^7$);
\item Theorems~9.5--9.7 (exponential all-input coalescence, continuum many extremal states, field response) and the finite-history Appendix~A that serves them;
\item Theorem~B.1 (the robustness radius);
\item Corollaries~9.8--9.12;
\item the scale remark, Remark~9.13 (the Chernoff alternative for the scales), and the numerical scale report of Section~9 (\texttt{verification/scale\_check.py});
\item two printed intermediate claims that Theorem~A does not use: the clause ``any other positive count eventually fails'' of Lemma~4.1(b) (F16), and Lemma~8.3 (F22).
\end{itemize}

\paragraph{Toolchain and release.} The pins are Lean~4.35.0 and Mathlib at the commits
\begin{flushleft}\small\leftskip1.5em
Lean: \nolinkurl{eb6cafb7290e8bd5f2937522d8d804f58f02bafc},\\
Mathlib: \nolinkurl{56a005855ea17e0b0958f73ec801ae140b0a1679}.
\end{flushleft}
The sources, pins, scope and check records are in the \texttt{lean/} folder of this paper's release repository. Dependencies and the reused cores under \texttt{vendor/} are fetched, not shipped. The reproduction instructions are in \texttt{README-ARXIV.md} there.
\begin{itemize}
\item \emph{The comparator.} A challenge imports only Mathlib. It contains verbatim excerpts of the two reused statement files and the whole of this development's statement file except its imports, followed by the three open proofs. The excerpts supply \texttt{FNE.flagLaw}, \texttt{FNE.WeakTendsto}, \texttt{UPD.Machine} with its run and halting definitions, and \texttt{UPD.Universal}. Auxiliary commands set the module names, set two elaboration options for the hard-model excerpt, and reset the per-file auxiliary-proof cache before each copy. The last command reproduces the empty cache of a separate module and affects only generated proof names. The comparator compares every constant reachable from the three statements with the development. The actual Linux run of \texttt{lake comparator -{}-config comparator.json}, inside the default bubblewrap sandbox and with the default kernel, accepted all three targets, with only \texttt{propext}, \texttt{Classical.choice} and \texttt{Quot.sound} permitted.
\item \emph{Two kernel replays.} An export of \texttt{Solution} was accepted by \texttt{leanchecker -{}-from-export} and by \texttt{nanoda}, the latter with only the same three axioms permitted. The negative control, with \texttt{Classical.choice} not permitted, was rejected.
\item \emph{Build and audit.} A fresh serial build from a clean copy compiled all 117 reachable modules, including the 40 imported core modules. Lint passed and the \texttt{sorry} census was zero. The fresh axiom audit covered all 3712 development declarations and 4452 declarations of the imported core modules, including private and generated declarations, and checked the types of the three final theorems. Only the three axioms above occur. The dependency sources match the pinned manifest; Mathlib was not rebuilt from source, and both kernel replays checked the exported proof closure.
\end{itemize}

In short form, the reproduction commands below use the recorded package layout and the toolchain, dependency and host-memory setup specified in \texttt{README-ARXIV.md}. They are run from its \texttt{project/} directory.
\begingroup\small
\begin{verbatim}
python3 -I -B scripts/replay.py
bash scripts/lint.sh
python3 -I -B scripts/sorry_census.py
python3 -I -B ../run-linux.py comparator
python3 -I -B ../run-linux.py export
python3 -I -B ../run-linux.py leanchecker
python3 -I -B ../run-linux.py nanoda
python3 -I -B ../run-linux.py nanoda-negative
\end{verbatim}
\endgroup




\begin{thebibliography}{99}\raggedright
\bibitem{BCW21} J. Bausch, T. S. Cubitt and J. D. Watson, \emph{Uncomputability of phase diagrams}, Nature Communications \textbf{12} (2021), 452. \href{https://doi.org/10.1038/s41467-020-20504-6}{doi:10.1038/s41467-020-20504-6}.
\bibitem{BermanSimon88} P. Berman and J. Simon, \emph{Investigations of fault-tolerant networks of computers}, Proceedings of the 20th Annual ACM Symposium on Theory of Computing (1988), 66--77. \href{https://doi.org/10.1145/62212.62219}{doi:10.1145/62212.62219}.
\bibitem{BMM13} A. Bu\v{s}i\'c, J. Mairesse and I. Marcovici, \emph{Probabilistic cellular automata, invariant measures, and perfect sampling}, Advances in Applied Probability \textbf{45} (2013), 960--980. \href{https://arxiv.org/abs/1010.3133}{arXiv:1010.3133}.
\bibitem{CH10} J.-R. Chazottes and M. Hochman, \emph{On the zero-temperature limit of Gibbs states}, Communications in Mathematical Physics \textbf{297} (2010), 265--281. \href{https://doi.org/10.1007/s00220-010-0997-8}{doi:10.1007/s00220-010-0997-8}.
\bibitem{Chernoff52} H. Chernoff, \emph{A measure of asymptotic efficiency for tests of a hypothesis based on the sum of observations}, Annals of Mathematical Statistics \textbf{23} (1952), 493--507. \href{https://doi.org/10.1214/aoms/1177729330}{doi:10.1214/aoms/1177729330}.
\bibitem{Choquet} G. Choquet, \emph{Existence et unicit\'e des repr\'esentations int\'egrales au moyen des points extr\'emaux dans les c\^ones convexes}, S\'eminaire Bourbaki, expos\'e 139 (1956--1958), 33--47.
\bibitem{CPGW15} T. S. Cubitt, D. Perez-Garcia and M. M. Wolf, \emph{Undecidability of the spectral gap}, Nature \textbf{528} (2015), 207--211. \href{https://doi.org/10.1038/nature16059}{doi:10.1038/nature16059}.
\bibitem{CPGW22} T. S. Cubitt, D. Perez-Garcia and M. M. Wolf, \emph{Undecidability of the spectral gap}, Forum of Mathematics, Pi \textbf{10} (2022), e14. \href{https://doi.org/10.1017/fmp.2021.15}{doi:10.1017/fmp.2021.15}.
\bibitem{DMforget} S. Douglas and N. Mghirbi, \emph{Forgetting under noise is undecidable: ergodicity of positive-rate cellular automata at a fixed noise rate}, version 1.4.0, 2026. \href{https://doi.org/10.5281/zenodo.24126025}{doi:10.5281/zenodo.24126025}.
\bibitem{DMwindows} S. Douglas and N. Mghirbi, \emph{A perturbed cellular automaton with infinitely many ergodicity transitions}, version 1.0.0, 2026. \href{https://doi.org/10.5281/zenodo.24122609}{doi:10.5281/zenodo.24122609}.
\bibitem{DMhard} S. Douglas and N. Mghirbi, \emph{An uncomputable phase diagram at positive temperature for one classical lattice model}, version 1.2.0, 2026. \href{https://doi.org/10.5281/zenodo.24144192}{doi:10.5281/zenodo.24144192}.
\bibitem{DMnil} S. Douglas and N. Mghirbi, \emph{A nilpotent cellular automaton that remembers under noise}, version 1.2.0, 2026. \href{https://doi.org/10.5281/zenodo.24120279}{doi:10.5281/zenodo.24120279}.
\bibitem{DMst} S. Douglas and N. Mghirbi, \emph{Exponential ergodicity and the space--time Gibbs states of one-dimensional probabilistic cellular automata}, version 1.2.0, 2026. \href{https://doi.org/10.5281/zenodo.24122307}{doi:10.5281/zenodo.24122307}.
\bibitem{DMthermal} S. Douglas and N. Mghirbi, \emph{Spontaneous magnetization and thermal decay are undecidable at positive temperature}, version 1.2.0, 2026. \href{https://doi.org/10.5281/zenodo.23962706}{doi:10.5281/zenodo.23962706}.
\bibitem{Gacs01} P. G\'acs, \emph{Reliable cellular automata with self-organization}, Journal of Statistical Physics \textbf{103} (2001), 45--267. \href{https://doi.org/10.1023/A:1004823720305}{doi:10.1023/A:1004823720305}.
\bibitem{Gacs21} P. G\'acs, \emph{A new version of Toom's proof}, 2021. \href{https://arxiv.org/abs/2105.05968}{arXiv:2105.05968}.
\bibitem{GST25} L. Gayral, M. Sablik and S. Taati, \emph{Characterisation of the set of ground states of uniformly chaotic finite-range lattice models}, Nonlinearity \textbf{38} (2025), 075011. \href{https://doi.org/10.1088/1361-6544/addbb8}{doi:10.1088/1361-6544/addbb8}.
\bibitem{GKLM89} S. Goldstein, R. Kuik, J. L. Lebowitz and C. Maes, \emph{From PCA's to equilibrium systems and back}, Communications in Mathematical Physics \textbf{125} (1989), 71--79. \href{https://doi.org/10.1007/BF01217769}{doi:10.1007/BF01217769}.
\bibitem{HM10} M. Hochman and T. Meyerovitch, \emph{A characterization of the entropies of multidimensional shifts of finite type}, Annals of Mathematics \textbf{171} (2010), 2011--2038. \href{https://doi.org/10.4007/annals.2010.171.2011}{doi:10.4007/annals.2010.171.2011}.
\bibitem{Hoeffding63} W. Hoeffding, \emph{Probability inequalities for sums of bounded random variables}, Journal of the American Statistical Association \textbf{58} (1963), 13--30. \href{https://doi.org/10.1080/01621459.1963.10500830}{doi:10.1080/01621459.1963.10500830}.
\bibitem{Kari92} J. Kari, \emph{The nilpotency problem of one-dimensional cellular automata}, SIAM Journal on Computing \textbf{21} (1992), 571--586. \href{https://doi.org/10.1137/0221036}{doi:10.1137/0221036}.
\bibitem{Kechris} A. S. Kechris, \emph{Classical Descriptive Set Theory}, Graduate Texts in Mathematics 156, Springer, 1995.
\bibitem{Kur78} G. L. Kurdyumov, \emph{An algorithm-theoretic method of studying homogeneous random media}, Doklady Akademii Nauk SSSR \textbf{238}(5) (1978), 1047--1050.
\bibitem{Lean4} L. de Moura and S. Ullrich, \emph{The Lean 4 theorem prover and programming language}, in Automated Deduction -- CADE 28, LNCS 12699, Springer (2021), 625--635.
\bibitem{LanfordRuelle69} O. E. Lanford III and D. Ruelle, \emph{Observables at infinity and states with short range correlations in statistical mechanics}, Communications in Mathematical Physics \textbf{13} (1969), 194--215.
\bibitem{Mathlib20} The mathlib Community, \emph{The Lean mathematical library}, in Proceedings of the 9th ACM SIGPLAN International Conference on Certified Programs and Proofs (CPP 2020), ACM (2020), 367--381.
\bibitem{MST19} I. Marcovici, M. Sablik and S. Taati, \emph{Ergodicity of some classes of cellular automata subject to noise}, Electronic Journal of Probability \textbf{24} (2019), article 41, 1--44. \href{https://doi.org/10.1214/19-EJP297}{doi:10.1214/19-EJP297}.
\bibitem{MT26} I. Marcovici and S. Taati, \emph{Positive-rate PCA and IPS with stationary Bernoulli measures are rapidly forgetful}, 2026. \href{https://arxiv.org/abs/2605.16904}{arXiv:2605.16904}.
\bibitem{MS24} H. Marsan and M. Sablik, \emph{Characterization of the set of zero-noise limits measures of perturbed cellular automata}, 2024. \href{https://arxiv.org/abs/2412.04672v2}{arXiv:2412.04672v2}.
\bibitem{MS25} H. Marsan and M. Sablik, \emph{Flexibility versus genericity of phase diagrams of perturbed continuous maps on the Cantor set}, 2025. \href{https://arxiv.org/abs/2508.00461}{arXiv:2508.00461}.
\bibitem{MST26} H. Marsan, M. Sablik and I. T\"orm\"a, \emph{A perturbed cellular automaton with two phase transitions for the ergodicity}, Journal of Statistical Physics \textbf{193} (2026), article 6. \href{https://doi.org/10.1007/s10955-025-03564-0}{doi:10.1007/s10955-025-03564-0}. \href{https://arxiv.org/abs/2507.03485}{arXiv:2507.03485}.
\bibitem{Ollinger08} N. Ollinger, \emph{Two-by-two substitution systems and the undecidability of the domino problem}, in \emph{Logic and Theory of Algorithms (CiE 2008)}, Lecture Notes in Computer Science \textbf{5028}, Springer, 2008, 476--485. \href{https://doi.org/10.1007/978-3-540-69407-6_51}{doi:10.1007/978-3-540-69407-6\_51}.
\bibitem{PLC24} J. Purcell, Z. Li and T. Cubitt, \emph{Chaitin Phase Transition}, 2024. \href{https://arxiv.org/abs/2410.02600}{arXiv:2410.02600}.
\bibitem{PW96} J. G. Propp and D. B. Wilson, \emph{Exact sampling with coupled Markov chains and applications to statistical mechanics}, Random Structures Algorithms \textbf{9} (1996), 223--252.
\bibitem{Toom80} A. L. Toom, \emph{Stable and attractive trajectories in multicomponent systems}, in R. L. Dobrushin and Ya. G. Sinai (eds.), \emph{Multicomponent Random Systems}, Advances in Probability and Related Topics \textbf{6}, Marcel Dekker, New York, 1980, 549--575.
\bibitem{Toom00a} A. Toom, \emph{Algorithmical unsolvability of the ergodicity problem for binary cellular automata}, Markov Processes and Related Fields \textbf{6} (2000), 569--577. \href{https://math-mprf.org/journal/articles/id882/}{Journal article page}.
\bibitem{Toom00b} A. Toom, \emph{Algorithmical unsolvability of the ergodicity problem for locally interacting processes with continuous time}, Journal of Statistical Physics \textbf{98} (2000), 495--501. \href{https://doi.org/10.1023/A:1018699527637}{doi:10.1023/A:1018699527637}.
\bibitem{Toometal90} A. L. Toom, N. B. Vasilyev, O. N. Stavskaya, L. G. Mityushin, G. L. Kurdyumov and S. A. Pirogov, \emph{Discrete local Markov systems}, in R. L. Dobrushin, V. I. Kryukov and A. L. Toom (eds.), \emph{Stochastic Cellular Systems: Ergodicity, Memory, Morphogenesis}, Manchester University Press, 1990.
\bibitem{vdBS99} J. van den Berg and J. E. Steif, \emph{On the existence and nonexistence of finitary codings for a class of random fields}, Annals of Probability \textbf{27} (1999), 1501--1522.
\bibitem{DMforgetLower} S. Douglas and N. Mghirbi, \emph{Forgetting under noise is undecidable}, version 2.1.0, 2026, Section~12. \href{https://doi.org/10.5281/zenodo.24178276}{doi:10.5281/zenodo.24178276}.
\bibitem{LGO12} B. Le Gloannec and N. Ollinger, \emph{Substitutions and Strongly Deterministic Tilesets}, in \emph{How the World Computes}, LNCS \textbf{7318} (2012), 462--471. \href{https://doi.org/10.1007/978-3-642-30870-3_47}{doi:10.1007/978-3-642-30870-3\_47}. Theorem~4; author manuscript HAL hal-00661887, version~1.
\bibitem{DMnoise13} S. Douglas and N. Mghirbi, \emph{One noisy cellular automaton with an uncomputable noise diagram}, version 1.3.0, 2026. \href{https://doi.org/10.5281/zenodo.24178271}{doi:10.5281/zenodo.24178271}.
\end{thebibliography}
\end{document}











